\def\ChapterCode{ANA}
\chapter
[\CO{[ANA]} Anatomie und Physiologie]
{Anatomie und Physiologie}
\label{chp:ANA}

\ChapterIntro
{%
Die Anatomie ist die Lehre vom Bau des Körpers, die
Physiologie die Lehre von dessen Funktionweise. Dieses
Kapitel stellt daher die Grundlage für die in den
Folgekapiteln beschriebenen Erkrankungen und Verletzungen
dar.
} {%
\begin{description}
  \item[Maintainer:] Lena Hirtler
  \item[Autoren:] Diverse
  \item[Reviewer:] Standard-Reviewprozess
  \item[Version:] {\DocVersion} ({\DocVersionInfo})
  \item[SHA1:] {\DocGitCommit}
\end{description}

{\TxTeilkapitelAASS}
}

\nocite{AumuellerAnatomie2}%
\nocite{BenninghoffAnatomie17Bd1}%
\nocite{LeohnhardtAnatomie3}
\nocite{LippertAnatomie7}%
\nocite{PlatzerTaschenatlasAnat8Bd1}%
\nocite{WaldeyerAnatomie17}%

% \AuthNotes{
% Changelog: 
% 
% \begin{description}
%     \item[2009-02-25] Start Changelog.
%     \item[2009-05-06] GABS: Korrekturen von KOCH
%     eingearbeitet.
% \end{description}
% }

\Text{Richtungs- und Lageangaben}{
Für die Beschreibungen von Richtungen und Lagebeziehungen
wird eine systematische Terminologie verwendet. Die
einschlägigen Bezeichnungen sind auch regelmäßig in Arzt-
und Entlassungsbriefen anzutreffen.

\begin{multicols}{2}
\Z{%
\begin{description}
  \item[lateral] \Lang{Abk.} lat.; seitlich
  \index{lateral|TLike}
  \index{lat.|see{lateral}}
  \item[medial] \Lang{Abk.} med.; mittig
  \index{medial|TLike}
  \index{med.|see{medial}}
  \item[dorsal] rückenwärts
  \index{dorsal|TLike}
  \item[ventral] bauchwärts
  \index{ventral|TLike}
  \item[cranial] kopfwärts
  \index{cranial|TLike}
  \item[caudal] fußwärts \Z{(eigentlich: schwanzwärts)}
  \index{caudal|TLike}
  \item[distal] \Lang{Abk.} dist.; körperfern
  \index{distal|TLike}
  \index{dist.|see{distal}}
  \item[proximal] \Lang{Abk.} prox.; körpernah
  \index{proximal|TLike}
  \index{prox.|see{proximal}}
  \item[dexter, -tra, trum] \Lang{Abk.} dext.; rechts
  \index{dexter|TLike}
  \index{dext.|see{dexter}}
  \item[sinister, -tra, trum] \Lang{Abk.} sin.; links
  \index{sinister|TLike}
  \index{sin.|see{sinister}}
\end{description}
}

\Cave{Seiten- und Richtungsangaben (rechts, links, \ldots)
beziehen sich immer auf die betreffende Person (Patienten)!}
\end{multicols}
}{}{}{}{}


\TextSlide{Gliederung des Körpers}{%
Der menschliche Körper wird grob in den \EE{Stamm} (Rumpf),
die \EE{Gliedmaßen} (\EE{Extremitäten}) sowie in \EE{Hals
und Kopf} unterteilt. Der Stamm gliedert sich in Brust
(Thorax) und Bauch (Abdomen). Der Schultergürtel und das
Becken stellen jeweils  den Übergang zwischen Rumpf und
oberen und unteren Extremitäten dar. 
\TODO{GABS}{2011-03-20}{Blumig ausformulieren}{Erfüllt Scutz Anknüpfungspunkte, }

}{
\begin{itemize}
    \item Stamm (Rumpf)
    \begin{itemize}
      \item Brust
      \item Bauch
    \end{itemize}
    \item Gliedmaßen (Extremitäten)
    \item Hals und Kopf
\end{itemize}
}{}{}{}

\section{Die Zelle}\index{Zelle}
\TextSlide{\TxBeschreibung}
{%
Die Zelle ist die kleinste lebende Einheit in unserem
Körper. Sie besteht u.\,a. aus einem
\E{Zellkern}\index{Zellkern}, der die gesamten
Erbinformationen in Form des Riesenmoleküls
\E{DNS}\ZFootnote{Desoxyribonukleinsäure; \Lang{engl.}
\E{DNA}.} enthält, sowie verschiedenen
Zellorganellen\index{Zellorganellen}. Die Zellorganellen
sind wichtig für den Zellstoffwechsel, man kann sie mit der
Funktion der Organe im menschlichen Körper vergleichen. Der
überwiegende Bestandteil der Zelle ist \E{Flüssigkeit}, das
Zellplasma. Die Zelle wird von der
\E{Zellmembran}\index{Zellmembran} umhüllt.
\cite{AlbertsMolZellbiologie2,LoefflerBiochemie7}
}{
\begin{itemize}
    \item Kleinste lebende Einheit
    \item Zellkern
    \begin{itemize}
        \item Erbinformationen (DNS)
    \end{itemize}
    \item Zellorganellen
\end{itemize}
}
{}{}{}



\begin{figure}[H]
    \includegraphics[width=\linewidth]{./images/hirtler-aass/Zelle-edited}
    \Captionof{figure}{Schema einer Zelle \SourcePicture{Hirtler}}
\end{figure}

\TextSlide{Gewebe}{%
Artverwandte Zellen können einen Zellverband bilden, welcher
\T{Gewebe} genannt wird. Diese Gewebe können dann größere
Strukturen bilden, wenn sich
verschiedene Gewebearten ergänzen entstehen hieraus 
Organe, wie z.\,B. Herz, Hirn, Lunge, Knochen oder Muskel. 
Organe, die dem gleichen Ziel dienen, können zu
Organsystemen zusammengefasst werden (z.B. Verdauungstrakt,
Harnbildendes und -leitendes System, Kreislauf, \ldots).

\VARIANT{VAR-ERNAEHRUNGSPAED}{
Man kann allgemein 4 Gewebearten unterscheiden:
\begin{itemize}
  \item Epithelgewebe: bildet die Oberfläche (=Haut) und
  kleidet Hohlorgane aus, Sinneszellen gehören auch in diese
  Gruppe.
  \item Binde- und Stützgewebe: Es verbindet
  andere Gewebe miteinander. Hierzu gehört u.\,a. auch das
  Fett. Knochen und Knorpel gehören zum Stützgewebe.
  \item Muskelgewebe: die verschiedenen Muskelarten werden
  an anderer Stelle beschrieben.
  \item Nervengewebe: Hierzu gehören die Zellen zur elektrischen
  Informationsübertragung des Körpers. 
\end{itemize}
} 
  \A{ Diskussion: Was ist ein
Organ?

\ACh{GABS}{2010-03-09}{Ist Muskel nicht ein Organ?}
\ACh{HIRT}{2010-03-09}{Wie du sicher weißt gibt es verschiedene Definitionen
von Organen.Bei den gängigsten fallen die Muskeln nicht hinein.}
\ACh{GABS}{2010-03-12}{Bei welcher? 

Pschyrembel: >>{aus Zellen u. Geweben
zusammengesetzte Teile des Körpers, die eine Einheit mit best. Funktionen
bilden.<< 

Wollen wir's nicht allgemein halten? Zumal die Definition}Organ} noch
ausständig ist \ldots{}
\ACh{HIRT}{2010-03-13}{Damit ist die Diskussion beendet.}
\ACh{GABS}{2010-03-14}{Neue Formulierung passt, weil unmissverständlich.}
}
}{
\begin{itemize}
    \item Zellverband
    \item → größere Strukturen
\VARIANT{VAR-ERNAEHRUNGSPAED}{    \item Gewebearten
    \begin{itemize}
      \item Epithelgewebe
      \item Binde- und Stützgewebe
      \item Muskelgewebe
      \item Nervengewebe 
    \end{itemize} 
}
\end{itemize}
}{}{}{}

\nocite{UlfigHistologie1,LippertAnatomie7,GantenMolGrund97}


\section{Das Skelett und der Bewegungsapparat}
\index{Skelett}

% \clearpage % TEMP temp clearpage
\TextSlide
{Funktionen des Skeletts}
{
Das Skelettsystem besteht aus \E{Knochen} und \E{Knorpeln} und dient 
hauptsächlich als
   Stützgerüst der \EE{Haltung} des Körpers (aufrechte Haltung)
    \ACh{GABS}{2010-03-09}{Auch für nicht-aufrechte Haltung wichtig?}
    \ACh{HIRT}{2010-03-09}{Hier geht es um die grundsätzliche Vorstellung,
    dass wir ohne Knochen nur ein Haufen Gatsch sind.}
    und als \E{Ursprungs- und Ansatzfläche für die Muskulatur}. Daher wird
    es oft auch als passiver Bewegungsapparat bezeichnet. Die Befestigung
    der Muskulatur an den Knochen ist Voraussetzung für die
    Wirkungsentfaltung der Muskeln. Weiters dient es als
\EE{Kalziumspeicher} (\ce{Ca++}) und der \EE{Blutbildung}
(Bildung von roten und weißen Blutkörperchen sowie den
Blutplättchen im roten Knochenmark).
     \ACh{GABS}{2010-03-09}{Wenn man die Körperchen aufzählt müssen die
    Thrombos aber auch hinein \ldots}
    \ACh{HIRT}{2010-03-09}{Vielen vielen Dank für den Hinweis\ldots}
}
{
\begin{itemize}
  \item Haltung
  \item Ursprungs- und Ansatzfläche für Muskulatur
  \item Kalziumspeicher
  \item Blutbildung
\end{itemize}
}
{}{}{}

\begin{figure}[H]
    \includegraphics[width=10cm]{./images/anatomie/Human_skeleton_front-edited.pdf}
    \Captionof{figure}{\AuthNotesPicLegalOk{}{ PD }Das Menschliche Skelett.
    \SourcePicture{Mariana Ruiz Villarreal, Public domain}}

\end{figure}




\subsection{Der Knochen}
\index{Knochen}

\A{Vorschlag: \protect\enquote{Der menschliche Körper besteht aus ???
Knochen.} Lena: nein, weil unklar wegen Sesambeinchen und z.B. Steißbein. }

\TextSlide
{Knocheneinteilung}
{%
Man kann die Knochen anhand ihrer Form einteilen in: 
\begin{itemize}
    \item \T[Knochen!platte]{Platte Knochen} (z.\,B. Beckenknochen,
    Schulterblatt, Brustbein, Schädel),
    \item \T{Röhrenknochen}\index{Knochen!Röhren-} (Oberarmknochen, Schienbein)
    und
    \item \T{Würfelknochen}\index{Knochen!Würfel-} (z.\,B. Hand-, Fußwurzelwurzelknochen,
    Wirbelknochen)
\end{itemize}
}
{
\begin{itemize}
  \item Platte Knochen
  \item Röhrenknochen
  \item Würfelknochen
\end{itemize}
}
{}{}{}

\AuthNotes{GABS \& KOCH: evtl. umstellen auf deutsch -- latein, aber was heißt
Kortikalis auf deutsch?}\AuthNotes{HIRT: Man könnte sie als Knochenrinde
bezeichnen.}

\TextSlide
{Aufbau eines Knochens}
{%
\begin{itemize}
  \item \T{Periost}: \E{Beinhaut}, überzieht den Knochen und
  ist sehr gut mit Nerven versorgt. Es ist daher die einzige Struktur des Knochens die
  schmerzempfindlich ist. Sie ist verantwortlich für
  Dickenwachstum und Frakturheilung\footnote{\T{Fraktur}:
  Bruch.}.
  \item \T{Kompakta} (Kortikalis, \protect\enquote{Knochenrinde}):
  \index{Kortikalis}feste äußere Schicht, äußere Form, verantwortlich für Stabilität des Knochens
  \item \T{Spongiosa}: \index{Knochenbälkchen}Knochenbälkchen
  ${\rightarrow}$ Zug- und Druckfestigkeit, rotes Knochenmark
  \item \T{Markhöhle}: Gefüllt mit Knochenmark
  \item \T{Epiphysenfuge}:  Wachstumsfuge
  \item \T{Gelenkfläche}: \index{Gelenkfläche!Aufbau des
  Knochens}mit knorpeligen Überzug als Teil eines Gelenkes.
  Manche Knochen haben keine Gelenksfläche.
\end{itemize}
Lange Röhrenknochen kann man außerdem in 
\begin{itemize}
    \item \T{Kopf}\index{Kopf!Knochenaufbau},
    \item \T{Hals}\index{Hals!Knochenaufbau} und
    \item \T{Schaft}\index{Schaft!Knochenaufbau} 
\end{itemize}
unterteilen.
}{
\begin{itemize}
  \item Allgemein
  \begin{itemize}
    \item Periost
    \item Kompakta
    \item Spongiosa
    \item Markhöhle
    \item Epiphysenfuge
    \item Gelenkfläche
  \end{itemize}
  \item Langer Röhrenknochen
  \begin{itemize}
    \item {Kopf}
    \item {Hals}
    \item {Schaft}
  \end{itemize}
\end{itemize}
}
{./images/anatomie/roehrenknochen-querschnitt.png}
{Ein langer Röhrenknochen im Querschnitt.}
{Hirtler}

% \Layout{27}
% {}
% {}
% {}
% {./images/hirtler-aass/needs-work/Knochenbau}
% {}
% {Hirtler}

\begin{table}[H]
\includegraphics[angle=90,width=\linewidth]{./images/hirtler-aass/needs-work/Knochenbau}
\caption{Ein langer Röhrenknochen im Querschnitt \SourcePicture{Hirtler}}
\end{table}



% \subsubsection{Das Knochenmark}
\index{Knochenmark}


% \fxnote[author=HIRT]{Review: ANA: Knochenmark, rotes: Entweder lange oder kurze
% Überschrift.} 
% GABS 2010-03-12: ok.

\TextSlide
{Rotes Knochenmark}
{%
% 
% Bitte Index-Einträge in den Fließtext hineingeben, nicht in Überschrift
% (es gibt noch alte Stellen wo die Einträge in der Überschrift stehen,  aber
% das wird nach und nach geändert!) GABS 2010-03-12
% 
\index{Knochenmark!rotes}
Es ist das wichtigste blutbildende Organ
des Körpers. Hier werden die \EE{zellulären Bestandteile des Blutes
gebildet}. Dazu gehören die \E{roten Blutkörperchen}
(\E{Erythrozyten}), die
\E{weißen Blutkörperchen} (\E{Leukozyten}) sowie die \E{Blutplättchen}
(\E{Thrombozyten}) -- \prettyref{topic:blut}. 

Bei Neugeborenen findet man das rote Knochenmark in
sämtlichen Knochen und Markhöhlen, im Laufe des Lebens bildet es sich in den
Markhöhlen zurück und verwandelt sich in {\quotedblbase}gelbes
Fettmark{\textquotedblleft}. Rotes Knochenmark findet man dann lediglich in platten Knochen wie Becken, Brustbein und
Schädel sowie in bereichen mit Spongiosa, wo es für die Blutbildung
verantwortlich ist und normalerweise zeitlebens erhalten bleibt.
%
}
{%
\begin{itemize}
  \item Bildung der zellulären Blutbestandteile
  	\subitem rote Blutkörperchen
  	\subitem weiße Blutkörperchen
  	\subitem Blutplättchen
  \item im Erwachsenenalter nur noch in platten Knochen  
\end{itemize}
}
{}{}{}
\Fazit{Das \E{rote Knochenmark} ist für die Blutbildung zuständig}

\AuthNotes{KOCH: wichtig? ev. weglassen?

GABS: Ja, wichtig wegen KT, onko-Patienten, etc. Soll
bleiben. 

KOCH: ok.} 

Einige medizinische Behandlungsmethoden verursachen als
Nebenwirkung eine massive \EE{Schädigung des
Knochenmarks}\index{Knochenmark!Schädigung} (z.\,B.
Bestrahlung, Chemotherapien bei Krebspatienten
(\protect\enquote{onkologische Patienten}), etc.). Diese
Patienten leiden neben einer Blutarmut auch an einem
schwachen Immunsystem, da nicht genügend weiße
Blutkörperchen gebildet werden. Sie können daher sehr leicht
schwerwiegende Infektionen bekommen.
Diese Patienten sind im Krankentransport häufig anzutreffen.

\Fazit{Es gibt Patienten, bei denen aufgrund ihrer
Grunderkrankung oder deren Behandlung eine
Knochenmarkschädigung und in Folge eine Immunschwäche
besteht.}

% \fxnote[author=HIRT]{Review: ANA: Knochenmark, gelbes: Entweder lange oder kurze
% Überschrift.} 
% GABS 2010-03-12: ok.

\TextSlide
{Gelbes Knochenmark\index{Knochenmark!gelbes}}
{%
Im Laufe des Lebens bildet sich in den
Röhrenknochen das rote Knochenmark  zu \E{fettreichem},
gelben Knochenmark um.
Das gelbe Fettmark kann bei schweren Knochenverletzungen von Bedeutung
sein, da verletzungs\-bedingt kleine Fetttröpfchen in die
Blutbahn gelangen können, die dann als Fettgerinnsel zu ernsten
Komplikationen führen können
(Fettembolie\AuthNotes{+ KOCH}\footnote{\I{Fettembolie}:
Vertopfungs eines Blutgefäßes durch ein Fettgerinsel.}).

Dieses Knochenmark ist nicht unbedingt für den Körper
notwendig. Deshalb kann man nach Knochenbrüchen auch Nägel in die
Markhöhle einsetzen ohne größere Auswirkungen auf den
Knochen zu verursachen.
%
}
{%
\begin{itemize}
  \item fettreich
  \item Grund für Komplikationen bei Frakturen
\end{itemize}
}
{}{}{}
\Fazit{Das \E{gelbe Knochenmark} besteht aus Fett}

% \subsubsection{Knochenwachstum}

\TextSlide
{Knochenentwicklung}
{%
Das Skelett eines Neugeborenen besteht bei der Geburt zum
größten Teil aus Knorpel. Erst mit zunehmendem Lebensalter
beginnt dieser von Knochenkernen ausgehend, zu verknöchern.
Eine
\EE{Wachstumszone}
(\T{Epiphysenfuge}) bleibt dabei allerdings erhalten.

Die Epiphysenfuge enthält beim Kind/Jugendlichen bis etwa
zum Ende der Pubertät \E{teilungsfähiges Knorpelgewebe},
daher ist ein \EE{Längenwachstum} möglich. Beim Erwachsenen
sind diese knorpeligen Fugen hingegen vollständig
\E{verknöchert}, sodass keinerlei Wachstum mehr stattfinden
kann.
}
{%
\begin{itemize}
  \item Geburt: Knorpel
  \item Knochenkerne
  \item Wachstumszone
  \item Verknöcherung der Wachstumszone
\end{itemize}
}
{}{}{}

\subsection{Gelenke}
\index{Gelenke}



\TextSlide
{Allgemeine Gelenklehre}
{\DefinitionInPara{Gelenke sind (meistens bewegliche) Verbindung zwischen
mindestens zwei Knochen.} Man kann zwischen Gelenken \E{mit
Gelenksspalt} (\protect\enquote{echtes Gelenk},
\Z{Diarthrose}) und Gelenken \E{ohne Gelenksspalt}
(\protect\enquote{unechtes Gelenk}, \Z{Synarthrose})
unterscheiden.
Zu den Gelenken ohne Gelenksspalt gehören alle
Knochenverbindungen die \E{nicht} frei beweglich sind,
z.\,B. die Schädelnähte, die Schambeinfuge und die
Rippen-Brustbein-Verbindungen des Brustkorbes (Thorax,
\prettyref{topic:thorax}) \Z{(Syndesmose, Synchondrose,
Synostose)}.

\VARIANT{VAR-ERNAEHRUNGSPAED}{
\begin{itemize}
  \item Syndesmose: \protect\enquote{Bandhaft}, bindegewebige Verbindung zwischen Knochen, z.\,B.
  zwischen Elle und Speiche und Schienbein und Wadenbein.
  \item Synchondrose: \protect\enquote{Knorpelhaft}, knorpelige Verbindung zwischen Knochen,
  z.\,B. die Bandscheiben und die Symphyse.
  \item Synostose: \protect\enquote{Knochenhaft}, knöcherne Verbindung zwischen Knochen, z.\,B.
  die Schädelnähte.
\end{itemize}
}

\ACh{HIRT}{2010-03-11}{Um allen Diskussionen zuvorzukommen: ich finde es wichtig, dass zumindest diese Tatsachen erwähnt werden.}
\ACh{GABS}{2010-06-07}{Fachbegriffe wie besprochen grau.}

% orientierende Zusammenfassung (GABS, 2010-06-07)
\Fazit{Das, was man umgangssprachlich unter dem Begriff
\protect\enquote{Gelenk} versteht, sind
\protect\enquote{echte} Gelenke mit Gelenksspalt.}
}
{%
\begin{itemize}
  \item Verbindung zwischen
  mind. zwei Knochen
  \item \protect\enquote{echtes Gelenk}
  \item \protect\enquote{unechtes Gelenk}
\end{itemize}
}{}{}{}

\Layout{20}
{Prinzipieller Aufbau eines \protect\enquote{echten Gelenks}
und Gelenksarten}
{%
Ein Gelenk ist von einer bindegewebigen \EE{Kapsel} umgeben
und durch Sehnen und Bänder stabil verbunden.
Die Gelenksflächen sind zur Reibungs\-ver\-minderung mit 
\EE{Knorpel} überzogen und der \EE{Gelenksspalt} ist mit
einer dünnen Schicht \EE{Schmierflüssigkeit}
\Z{(Synovialflüssigkeit)} gefüllt. 

\index{Gelenke!Arten}%
\index{Gelenk!Kugel-}%
\index{Gelenk!Ei-}%
\index{Gelenk!Sattel-}%
\index{Gelenk!Scharnier-}%
\index{Gelenk!Zapfen-|Z}%
Es gibt verschiedene Arten von Gelenken: Klassische
Beispiele sind das \T{Kugelgelenk} (Hüftgelenk),
 das \T{Eigelenk} (z.\,B. Handwurzelgelenk),
das \T{Sattelgelenk} (Daumengrundgelenk), das
\T{Scharniergelenk} (Fingergliedgelenke, Gelenk zwischen
Oberarmknochen und Elle) \Z{sowie das Zapfengelenk (z.\,B.
Radioulnargelenk)}.
\Commentary[Gelenksarten, Beispiel Scharniergelenk]{Scharniergelenk ist nur das Gelenk, welches von Elle und Oberarmknochen gebildet wird.}


\bigskip

\bigskip

}{%
\begin{itemize}
  \item Kapsel
  \item Knorpelüberzug der Gelenkflächen
  \item Gelenksflüssigkeit
  \item Gelenkarten
  \begin{itemize}
  \item Kugelgelenk (Hüftgelenk, \ldots)
  \item Eigelenk (Handwurzelgelenk, \ldots)
  \item Sattelgelenk (Daumengrundgelenk, \ldots)
  \item Scharniergelenk (Fingergliedgelenke, Gelenk zwischen Elle und
  Oberarmknochen, \ldots)
  \item Zapfengelenk (Radioulnargelenk, \ldots)
\end{itemize}
\end{itemize}
}
{./images/hirtler-aass/Gelenk-edited.pdf}
{\AuthNotesPicLegalOk{}{Lena }Ein Gelenk im Querschnitt.}
{Hirtler}




\Layout{27}
{}
{}
{}
{./images/hirtler-aass/Gelenkarten.pdf}
{\EE{Gelenksarten}: \\
  \EE{A:} Kugel\-gelenk (z.\,B. Hüft\-gelenk),\\
  \EE{B:} Ei\-gelenk (z.\,B. Hand\-wurzel\-gelenk),\\
  \EE{C:} Sattelgelenk (Daumen\-grund\-ge\-lenk),\\
  \EE{D:} Scharniergelenk (z.\,B.
  Finger\-glied\-ge\-lenke, Ge\-lenk zwi\-schen Elle und
  Ober\-arm\-kno\-chen),\\
  \EE{E:} Zapfengelenk (z.\,B.
  Radioulnargelenk). }
{Hirtler}



\subsection{Der Schädel} 
%%% Index und Label
\index{Schädel}
\label{topic:schaedel}

\HeadingSub{Kopf (\T{Caput}) / Schädel (\T{Cranium})}

 


\begin{WidePar}
\begin{tabularx}{\linewidth}{XX}
\includegraphics[width=\linewidth]{./images/anatomie/Human_skull_side_simplified_hirngesicht_de-edited.pdf}
\Captionof{figure}{\AuthNotesPicLegalOk{}{PD}
\index{Hirnschädel}\index{Gesichtsschädel}Der Schädel
gliedert sich in den \T{Hirnschädel} und den \T{Gesichtsschädel}. 
\SourcePicture{Mariana Ruiz Villarreal, Public Domain}}
&
\includegraphics[width=\linewidth]{./images/anatomie/Human_skull_front_simplified_bones_de-edited.pdf}
\Captionof{figure}{\AuthNotesPicLegalOk{}{PD}Schädel Vorderansicht.
\SourcePicture{Mariana Ruiz Villarreal, Public Domain}}
\\
\end{tabularx}
\end{WidePar}



\TextSlide{Aufbau des Hirnschädels}
{%
Der Hirnschädel besteht aus der \EE{Schädelbasis} und dem
\EE{Schädeldach} (Kalotte). Beide Strukturen werden durch
platte Knochen gebildet, welche beim Erwachsenen durch
\E{Schädelnähte} \index{Schädelnähte} miteinander verzahnt
und verknöchert sind.

\bigskip

\bigskip
}{%
\begin{itemize}
    \item Schädelbasis
    \item Schädeldach (Kalotte)
    \item Platte Knochen, durch Schädelnähte verzahnt und
    verknöchert
\end{itemize}
}
{}{}{}



\Layout{20}{Fontanelle}{%
Die \T[Fontanelle]{Fontanellen} sind \EE{knorpelige 
Verbindungen} zwischen Schädelknochen beim Neu\-geborenen
(der Schädel muss sich bei der Geburt verformen
können!). \Z{Sie schließen  sich bis zum 18.
Lebensmonat,} verknöchern  aber erst vollständig mit
zunehmendem  Lebensalter \Z{(ca. 20 Lebensjahr
abgeschlossen)}.

Es gibt zwei Fontanellen: Die \E{große Fontanelle} befindet
sich zwischen dem Stirnbein und dem Scheitelbeinen, die
\E{kleine} Fontanelle zwischen den Scheitelbeinen und dem
Hinterhauptsbein.
Die Unterscheidung ist für die Geburtshilfe interessant, da
man dadurch unterscheiden kann in welche Richtung der Kopf
des Kindes schaut.

\Z{Die \T{Sutur} ist die endgültig fest verknöcherte Naht
zwischen den Schädelknochen beim Erwachsenen ${\rightarrow}$
starrer  stabiler Schädel, Schutz  für Gehirn.}

}
{
\begin{itemize}
  \item Verschluss bis 18. Lebensmonat
  \item \Z{große Fontanelle}
  \item \Z{kleine Fontanelle}
\end{itemize}
}
{./images/anatomie/fontanelle-dt.jpg}
{Fontanellen.}
{Gray's Anatomy, Copyright expired.}



\Layout{20}{Schädelbasis}{%
Das Gehirn liegt auf der Schädelbasis \index{Schädelbasis}
auf \Z{(außer wenn man auf dem Kopf steht)}. Hirnnerven
treten durch diverse Löcher aus dem Schädel aus. \EE{Das
\T[Großes Hinterhauptsloch]{große
Hinterhauptsloch}\index{Hinterhauptsloch!Großes} 
\Z{(Foramen magnum)} ist die Durchtrittsstelle des
Rückenmarks.}

\subparagraph{Exkurs: \T{Hirnstammeinklemmung}}
\index{Hirnstammeinklemmung} 

Das \EE{große Hinterhauptsloch} ist die Durchtrittsstelle
des Rückenmarks. Bei einer \EE{Hirndrucksteigerung}
(\E{Hirnödem, Hirnblutung}) ist es die einzige Möglichkeit
zur Ausdehnung des Gehirns! \EE{Dabei drückt sich das Gehirn
allerdings selbst ab}, es kommt zu einer tödlichen
\Ca{Hirnstammeinklemmung!} ${\rightarrow}$
\Ca{Lebensgefahr}, da das Kreislauf- und Atemzentrum
abgedrückt wird.
}
{
\begin{itemize}
  \item Auflagefläche des Gehirns
  \item Hirnnerven
  \item Großes Hinterhauptsloch \Z{(Foramen magnum)}
  \begin{itemize}
    \item Durchtritt Rückenmark
    \item Hirnstammeinklemmung
  \end{itemize}

\end{itemize}
}
{./images/anatomie/cavernous_sinus.png}
{Aufsicht auf die Schädelbasis.}
{Gray's Anatomy, Copyright expired.}


\clearpage % TEMP temp clearpage
\subsection{Die Wirbelsäule}
\index{Wirbelsäule}
\TextSlide{\TxBeschreibung}
{%
Die Wirbelsäule bildet das knöcherne Rückgrat des Körpers
und ist damit eine der wichtigsten Stützen des Körpers.
Andererseits ist sie auch Ansatzfläche vieler wichtiger
Muskeln. Sie hat eine \E{doppelt S-förmig gekrümmte Form},
welche den \E{aufrechten Gang} ermöglicht.

\subparagraph{Bestandteile}

Die Wirbelsäule besteht aus mehreren Komponenten. Die
knöchernen, übereinander \protect\enquote{gestapelten} \T{Wirbel}
\Z{(Vertebrae)} geben der Wirbelsäule ihre Form und
Stabilität. Die Wirbel sind i.\,d.\,R. zueinander im
gewissen Umfang \E{beweglich}.
Zwischen den einzelnen Wirbeln finden sich \T{Bandscheiben}
(\T{Discus}), welche als \EE{Stoßdämpfer} dienen. Mit
zunehmendem Alter können diese \E{verschleißen}.
Von den Wirbeln wird ein Kanal, der sog. \T{Wirbelkanal},
gebildet. In diesem ist das \EE{Rückenmark} eingebettet.

\subparagraph{Gliederung} Die Wirbelsäule gliedert sich in
verschiedene Abschnitte: 
\index{Halswirbel}
\index{Brustwirbel}
\index{Lendenwirbeln}
\index{Kreuzbeinwirbel}
\index{Steißbeinwirbel}
Es gibt  7 \EE{Hals}wirbel \Z{(Cervical-Wirbel,
C\,\VonBis{1}{7})},  12 \EE{Brust}wirbel 
\Z{(Thorakal-Wirbel, Th\,\VonBis{1}{12})}, 5
\EE{Lenden}wirbel \Z{(Lumbal-Wirbel, L\,\VonBis{1}{5})}, 5
verschmolzene \EE{Kreuzbein}wirbel \Z{(Os sacrum,
S\,\VonBis{1}{5})}, sowie \VonBis{3}{4}
\EE{Steißbein}\protect\enquote*{wirbel} \Z{(Os coccygis)}.
Analog zu den o.\,g. Wirbeln benennt man die ersten drei Teile 
der Wirbelsäule als \I[Halswirbelsäule]{Hals-},
\I[Brustwirbelsäule]{Brust-} und \I{Lendenwirbelsäule}
(\I{HWS}, \I{BWS}, \I{LWS}).


\subparagraph{Klinischer Bezug} 
Wenn die Wirbelsäule beschädigt wird, kann es gleichzeitig
zu einer Beschädigung des Rückenmarks kommen
(\E{Wirbelsäulentrauma}\VARIANT{VAR-STANDARD}{,
\FullPrettyRef{topic:wirbelsaeulentrauma}}).
Bandscheiben können aufplatzen und der weichere Kern \E{in
den Wirbelkanal rutschen}. Durch den Druck auf das
Rückenmark kann das starke Schmerzen und Nervenausfälle
verursachen (\E{Bandscheibenvorfall}\VARIANT{VAR-STANDARD}{,
\FullPrettyRef{topic:lumbago}}).


}{%
\begin{itemize}
  \item Doppelt S-förmig gekrümmtes
Achsenskelett, ermöglicht  aufrechten Gang 
\item \TabSub{Bestandteile}
\begin{itemize}
    \item Wirbel \Z{(Vertebrae)}
    \item \index{Bandscheiben}\T{Bandscheiben}
    (\T{Discus}): dienen der \protect\enquote{Stoßdämpfung}
    \item \EE{Im Wirbelkanal liegt das Rückenmark
    eingebettet.}
\end{itemize}
\item \TabSub{Gliederung}
\begin{itemize}
    \item 7 \EE{Hals}wirbel\\
    \Z{(Cervical-Wirbel, C\,\VonBis{1}{7})}
    \index{Halswirbel}
    \item 12 \EE{Brust}wirbel\\
    \Z{(Thorakal-Wirbel, Th\,\VonBis{1}{12})}
    \index{Brustwirbel}
    \item 5 \EE{Lenden}wirbel\\
    \Z{(Lumbal-Wirbel, L\,\VonBis{1}{5})}
    \index{Lendenwirbeln}
    \item 5 verschmolzene \EE{Kreuzbein}wirbel\\
    \Z{(Os sacrum, S\,\VonBis{1}{5})} \index{Kreuzbeinwirbel}
    \item \VonBis{3}{4}
    \EE{Steißbein}\protect\enquote*{wirbel}\\
    \Z{(Os coccygis)} \index{Steißbeinwirbel}
\end{itemize}
\item Klinisch: Bandscheibenvorfall, Wirbelsäulentrauma
\end{itemize}

}
{./images/hirtler-aass/Wirbelsaeule-Ueberblick.pdf}
{Die Wirbelsäule.}
{Hirtler. }

\noindent\Parallel{\includegraphics[width=\linewidth]{./images/hirtler-aass/Wirbelsaeule-Ueberblick.pdf}\\
\captionof{figure}{Die Wirbelsäule.
\SourcePicture{Hirtler}}
}
{\flushleft\begin{itemize}
    \item 7 \EE{Hals}wirbel\\
    \Z{(Cervical-Wirbel, C\,\VonBis{1}{7})}
    \index{Halswirbel}
    \item 12 \EE{Brust}wirbel\\
    \Z{(Thorakal-Wirbel, Th\,\VonBis{1}{12})}
    \index{Brustwirbel}
    \item 5 \EE{Lenden}wirbel\\
    \Z{(Lumbal-Wirbel, L\,\VonBis{1}{5})}
    \index{Lendenwirbeln}
    \item 5 verschmolzene \EE{Kreuzbein}wirbel\\
    \Z{(Os sacrum, S\,\VonBis{1}{5})} \index{Kreuzbeinwirbel}
    \item \VonBis{3}{4}
    \EE{Steißbein}\protect\enquote*{wirbel}\\
    \Z{(Os coccygis)} \index{Steißbeinwirbel}
\end{itemize}}




\clearpage

\Layout{20}{Die Wirbel: Aufbau und Arten}
{\index{Wirbel}%
Ein Wirbel besteht aus einem 
\T{Wirbelkörper}\index{Wirbelkörper} und einem
 \T{Wirbelbogen}\index{Wirbelbogen}, welcher das
  Rückenmark umschließt.
Vom Wirbelbogen entspringen
 2 \T{Querfortsätze}\index{Querfortsätze} und
 1 \T{Dornfortsatz}\index{Dornfortsatz}.


\VARIANT{VAR-ERNAEHRUNGSPAED}{
Ausnahmen hierzu bilden der 1. Wirbel
(\protect\enquote{Atlas}, hat keinen Wirbelkörper) und das
Kreuzbein (besteht aus zusammengewachsenen Wirbeln,
Querfortsätze sind zu Gelenkfächen mit den Beckenschaufeln
ausgewachsen, Wirbelloch ist durchgängig verknöchert) bzw.
das Steißbein (verkümmerte Wirbeln, kein Wirbelloch, meist
auch keine Fortsätze mehr erkennbar).
}

Der Wirbekörper trägt die Last des Rumpfes, zwischen ihnen
liegen die \E{Bandscheiben}. Der Wirbelbogen \EE{schützt das
Rückenmark} vor Verletzungen, die Querfortsätze dienen
z.\,B. zur Befestigung der Rippen. Die Dornfortsätze sind
jene Teile der Wirbelsäule, die am Rücken tastbar sind.
In den verschiedenen Abschnitten der Wirbelsäule sehen die
Wirbel unterschiedlich aus.
}
{%
Aufbau:
\begin{itemize}
  \item Wirbelkörper
  \item Wirbelbogen
  \begin{itemize}
    \item 2 Querfortsätze
    \item 1 Dornfortsatz
  \end{itemize}
\end{itemize}
}{./images/hirtler-aass/Wirbel-allgemein.pdf}
{Wirbel.}
{Hirtler}






\subsection{Der Brustkorb}
\index{Brustkorb}
\index{Thorax}
\index{Brustbein}
\index{Sternum}
\label{topic:brustkorb-skelett}
\label{topic:thorax}

\HeadingSub{\Lang{Term.} \T{Thorax}}

\Layout{20}{\TxBeschreibung}
{%
Der  Brustkorb besteht \E{vorne} aus dem \T{Brustbein}
 (\T{Sternum}) und \EE{12
Rippenpaaren}, welche \E{hinten} mit der
\EE{Brustwirbelsäule} (12 Brust\-wirbel) in gelenkiger
Verbindung stehen.

\index{Rippen}
\index{Rippen!echte}
\index{Rippen!unechte}
\index{Rippen!fliegende}
Die obersten 7 Rippen-Paare werden als
{\quotedblbase}\T{echte Rippen}{\textquotedblleft}
bezeichnet. Sie erreichen mit ihren Spitzen das Brustbein
und bilden praktisch {\quotedblbase}Ringe{\textquotedblleft}
um die Brustorgane. Die 5 unteren Rippenpaare bilden den
bogenförmigen unteren Rand des Brustkorbes. Die Rippenpaare
\VonBis{8}{10} sind dabei nur indirekt über einen Knorpel
mit dem Brustbein verbunden ({\quotedblbase}\T{unechte
Rippen}{\textquotedblleft},
\emph{{\quotedblbase}echt{\textquotedblleft}} bedeutet, die
Rippe hat eine direkte Verbindung zum Sternum), das 11. und
das 12. Paar endet mit knorpeligen Spitzen frei zwischen den
Bauchmuskeln({\quotedblbase}\T{fliegende
Rippen}{\textquotedblleft}).

Zwischen den
Rippen spannt sich die
\T{Zwischenrippenmuskulatur}\index{Zwischenrippenmuskulatur}
(wichtig für \EE{Atemmechanik}), in diese sind segmentale
Nerven- und Gefäßbündel eingelagert. Das \EE{Zwerchfell}
\index{Zwerchfell} schließt den Brustraum nach unten hin ab
und trennt somit den Brustraum vom Bauchraum (Vgl.
Atemmechanik, \FullPrettyRef{topic:atemmechanik}).

Der Brustkorb enthält und schützt die Brust\-organe (Lunge,
Herz) sowie die Oberbauchorgane (Leber, Magen, Milz). 
Nach Rippenfrakturen kann es durch spitze
Fragmente zu einer Lungenfell- oder Lungenverletzung kommen,
die lebensgefährlich werden kann.

\Fazit{Die Intaktheit des Brustkorbes ist für die Atmung von
großer
Bedeutung.}
}
{
\begin{itemize}
    \item \EE{Brustbein} (\T{Sternum})
    \item \EE{12  Rippenpaare}
    \begin{itemize}
        \item \Z{7 \protect\enquote{echte}}
        \item \Z{3 \protect\enquote{unechte}}
        \item \Z{2 \protect\enquote{fliegende}}
    \end{itemize}
    \item \EE{Brustwirbelsäule} (12 Brust\-wirbel), mit
    Rippen in gelenkiger Verbindung
\end{itemize}
}
{./images/anatomie/Thorax.pdf}
{Der knöcherne Thorax.}
{Mariana Ruiz Villarreal, Public Domain}



\subsection{Der Schultergürtel und die obere Extremität}
\index{Schultergürtel}
\label{topic:schulterguertel}

\Layout{27}
{}
{}
{}
{./images/anatomie/Human_arm_bones_diagram_de-edited.pdf}
{Der
Schultergürtel und die obere Extremität.}
{Mariana Ruiz Villarreal, Public Domain}


\TextSlide{Der Schultergürtel}
{
Der Schultergürtel besteht (von vorne nach
hinten) aus:

\begin{itemize}
    \item \index{Brustbein}\T{Brustbein}
    (\index{Sternum}\T{Sternum})
    \item \index{Schlüsselbein}\T{Schlüsselbein}
    \Z{(Clavicula)}
    \item \index{Schulterblatt}\T{Schulterblatt} \Z{(Scapula)}
\end{itemize}

Der Schultergürtel ermöglicht die Bewegung der Arme. Über
eine Bandverbindung des Schlüsselbeines am Brustbein und
wiederum eine Bandverbindung zwischen Schlüsselbein und
Schulterblatt sind alle Knochen miteinander verbunden. Die
Gelenkspfanne für den Oberarmknochen befindet sich am
Schulterblatt.

Der Oberarm ist hauptsächlich durch Muskelzug \Z{(M.
deltoideus, Rotatorenmanschette)} am Schultergürtel
befestigt.

\Fazit{Bänder und Muskel sind für die Stabilität des Schultergelenkes wichtig.}
}
{
\begin{itemize}
    \item Brustbein (Sternum)
    \item Schlüsselbein \Z{(Clavicula)}
    \item Schulterblatt \Z{(Scapula)}
    \item Befestigung des Oberarmes durch Muskelzug
\end{itemize}
}{}{}{}

\clearpage
\subsubsection{}


\TextSlide{Die obere Extremität}
{\index{Extremität!obere}%
\label{topic:obere-extremitaet}%
Der Arm beginnt mit dem \T{Schultergelenk}. Dieses ist ein
klassisches \E{Kugelgelenk}, es ermöglicht die Bewegung in
allen drei Raumebenen. Die beteiligten Knochen im
Schulterglenk sind das \E{Schulterblatt} (Scapula) und
\T{Oberarmknochen} (Humerus). Letzterer bildet als einzelner
Knochen den Oberarm.

Das \T{Ellenbogengelenk} liegt zwischen Oberarm und
Unterarm. Es wird aus \EE{Oberarmknochen} und \EE{Elle}
(Ulna) sowie \EE{Speiche} (Radius) gebildet. Das
Ellenbogengelenk ist ein zusammengesetztes Gelenk. Die
Verbindung zwischen Elle und Oberarmknochen ist ein
Scharniergelenk und ermöglicht das Abwinkeln des Armes. Die
Verbindung zwischen Speiche und Oberarmknochen hingegen ist
ein Kugelgelenk. In Kombination mit dem Zapfengelenk
zwischen Elle und Speiche ermöglicht dieses Gelenk das
Drehen des Unterarmes und der Hand.

Die beiden Knochen des \EE{Unterarmes} sind die \T{Elle}
(Ulna) und die \T{Speiche} (\T{Radius}), \Z{sie liegen
bei Drehung der Handfläche nach oben parallel, bei Drehung
der Handfläche zum Boden hingegen überkreuz.} Die
\EE{Speiche} (Radius) liegt \EE{daumenseitig}\footnote{Daher
ist die Radialarterie daumenseitig zu tasten!}, die Elle an
der Seite des kleinen Fingers.

Es folgt das \EE{Handwurzelgelenk}, welches hauptsächlich
aus der Speiche und den \E{Handwurzelknochen} gebildet wird.
Das Handwurzelgelenk ist ein \EE{Eigelenk}, \Z{es hat fast
so viele Freiheitsgrade wie ein Kugelgelenk, in einer Ebene
ist jedoch die Bewegung durch die Ei-Form eingeschränkt
(Kippbewegungen in Richtung von Daumen und Kleinfinger).}
Die \EE{Hand} beginnt auf Höhe des Gelenkes.

Es gibt insgesamt 8 \T{Handwurzelknochen},
diese sind \E{Würfelknochen}, \Z{welche in zwei Reihen zu
jeweils 4 Knochen angeordnet sind.}

Hierauf folgen \EE{5 \T{Mittelhandknochen}} und dann die 5
Finger mit jeweils 2 (Daumen) oder 3 (Finger \VonBis{II}{V})
\T{Fingergliedknochen} (\E{körpernahe}, \E{mittlere} und
\E{körperferne} Fingergliedknochen). Die Grundgelenke
(Mittelhandknochen $\leftrightarrow$ körpernahe
Fingergliedknochn) sind Sattelgelenke, die
\EE{Fingergliedgelenke} sind \EE{Scharniergelenke}.
}{
\begin{itemize}
    \item \EE{Schultergelenk} \index{Schultergelenk}
    \item \EE{Oberarmknochen} \index{Oberarmknochen} 
    (\T{Humerus}\index{Humerus})
    \item \EE{Ellenbogengelenk} \index{Ellenbogengelenk}
    \item (2) \EE{Unterarmknochen} \index{Unterarmknochen} 
    \begin{itemize}
        \item \EE{Elle} \index{Elle}
        (\T{Ulna}\index{Ulna})
        \item \EE{Speiche} \index{Speiche}
        (\T{Radius}\index{Radius}), daumenseitig!
    \end{itemize}
    \item Handwurzelgelenk \index{Handwurzelgelenk}
    \item (8) \EE{Handwurzelknochen}
    \index{Handwurzelknochen}
    \item (5) \EE{Mittelhandknochen}
    \index{Mittelhandknochen}
    \item \EE{Fingergliedknochen}
    \index{Fingergliedknochen} mit \VonBis{2}{3} Gliedern
\end{itemize}
}
{}{}{}

\AuthNotes{GABS: Bild evtl. nur mit einer Seite, dafür aber
größer und neben Text?

KOCH: passt so.}

\Fazit{Achtung: Unterscheide: Hand --
Arm!\index{Hand}\index{Arm}}

\subsection{Der Beckengürtel und die untere Extremität}
\index{Beckengürtel}
\index{Extremität, untere}
\label{topic:beckenguertel}

\TextSlide{Der Beckengürtel}
{
Der Beckengürtel besteht (von hinten nach vorne) aus dem
\begin{itemize}
    \item \T{Kreuzbein} (Os sacrum) und dem 
    \item \T{Beckenknochen}\Z{, welcher wiederum aus dem 
\begin{compactitem}
        \item \Z{\index{Darmbein|Z}Darmbein, dem}
        \item \Z{\index{Sitzbein|Z}Sitzbein und dem }
        \item \Z{\index{Schambein|Z}Schambein besteht.}
    \end{compactitem}
    }
\end{itemize}
Die \T{Symphyse} stellt eine straffe
Knorpelverbindung vorne zwischen den beiden Beckenknochen
dar (Dehnungsfuge, für Geburtsvorgang notwendig).
Die Stabilität des Beckengürtels ist für Stehen, Gehen und
alle sonstigen Bewegungen der Beine unabdingbar. Die
Gelenkspfanne des Hüftgelenkes befindet sich im
Beckenknochen. Der Oberschenkelknochen wird hauptsächlich
durch Bänder im Gelenk gehalten.

Nach einem Unfall kann es zu einer Fraktur des
Beckengürtels kommen (\E{Beckentrauma}\VARIANT{VAR-STANDARD}{,
\FullPrettyRef{topic:beckentrauma}}). Stehen und Gehen ist dann
nicht mehr möglich, der Patient verspürt starke Schmerzen
bei jeder passiven Bewegung \E{beider} Beine.

}
{
\begin{itemize}
    \item Kreuzbein (Os sacrum)
    \item Beckenknochen
    \begin{itemize}
        \item \Z{Darmbein}
        \item \Z{Sitzbein}
        \item \Z{Schambein}
    \end{itemize}
    \item Symphyse
\end{itemize}
}
{}{}{}

\clearpage
\Layout{20}{Die untere Extremität}
{%
\label{topic:untere-extremitaet}%
\index{Schenkelhals}%
Das Bein beginnt mit dem \T{Hüftgelenk}, einem
\EE{Kugelgelenk}, gebildet aus dem \E{Beckenknochen} und dem
\E{Oberschenkelknochen}. Im Gegensatz zum Schultergelenk ist
die Pfanne im Hüftgelenk besser ausgebildet und die
Befestigung des Gelenkskopfes in der Pfanne wird durch
starke Bänder bewerkstelligt.

Der \T{Oberschenkelknochen} (Femur) bildet als einziger
Knochen den Oberschenkel. Er ist gekenzeichnet durch einen
ausgeprägten \T{Schenkelhals} mit \protect\enquote{Knick}
zwischen Kopf-Hals-Bereich und Körper. Körperfern endet der
Oberschenkelknochen im Kniegelenk.

Das \T{Kniegelenk} verbindet den Ober- und Unterschenkel und
wird aus dem Oberschnkelknochen und dem Schienbein
(Unterschenkel) gebildet. Es ist \E{kein} reines
Scharniergelenk, Drehungen sind bei gebeugtem Knie möglich.
% Im Kniegelenk sind nur Oberschenkelknochen und Schienbein
% beteiligt, nicht das Wadenbein.
Die Besonderheit des Kniegelenkes besteht einerseits in den
Bandbefestigungen: Kreuzbänder (2) und Seitenbänder (2)
werden gemeinsam mit den Menisken (2) häufig bei diversen
Sportarten verletzt. Andererseits gibt es die
\EE{Kniescheibe}, sie dient der besseren Kraftübertragung
des Kniestreckers.
\cite{PetersenVordereKreuzband1,MuellerKnie1,SchabusKnie1}

Der \T{Unterschenkel} wird vom \T{Schienbein} (Tibia; dick,
vorne) und \T{Wadenbein} (Fibula; dünn, seitlich hinten)
gebildet. Sie gehen im Sprunggelenk in den Fuß über und
bilden dort den Innen- (Schienbein) und Außenknöchel
(Wadenbein).

Im \T{Sprunggelenk} gibt es, ähnlich wie bei der Hand,
\T{Fußwurzelknochen} (7, \Z{der bekannteste ist das
Fersenbein}), \T{Mittelfußknochen} (5) und 5 Zehen mit
jeweils 2 (Großzehe) bis 3 Gliedern.
}{
\begin{itemize}
    \item \EE{Hüftgelenk} \index{Hüftgelenk} 
    \item \EE{Oberschenkel} \index{Oberschenkel}  
    (Femur)
    \index{Femur} 
    \item \EE{Kniegelenk} \index{Kniegelenk} 
    \item \EE{Kniescheibe} \index{Kniescheibe} 
    \item \EE{Unterschenkelknochen} (2)
    \index{Unterschenkelknochen}

    \begin{itemize}
        \item \EE{Schienbein}  (Tibia) , dick, vorne
        \index{Schienbein}  \index{Tibia} 
        \item \EE{Wadenbein}  (Fibula), dünn, seitlich
        hinten \index{Wadenbein}  \index{Fibula} 
    \end{itemize}
    \item \EE{Sprunggelenk} \index{Sprunggelenk} 
    \item 7 \EE{Fußwurzelknochen}, \Z{u.a. Sprungbein,
    Fersenbein} \index{Fußwurzelknochen} 
    \item 5 \EE{Mittelfußknochen}  
    \index{Mittelfußknochen}
    \item \EE{Zehengliederknochen} mit je 2-3 Gliedern
    \index{Zehengliederknochen} 
\end{itemize}
}
{./images/anatomie/Human_skeleton_lower-extr-edited.pdf}
{\AuthNotesPicLegalOk{}{}Der Beckengürtel und die untere Extremität.}
{Mariana Ruiz Villarreal, Public Domain}


% GABS, 2010/05/04: Die Beschreibung bezgl. Trauma und Massnahmen ist glaub ich
% im Trauma-Kapitel besser aufgehoben, der Text in TRA ist eh noch nciht berühmt
% → dorthin verschoben.

% % GABS: Ad Trauma: 
% Eine Kompression des Beckens ist wichtig um den
% möglichen großen Blutverlust (einige Liter!) einzudämmen.

Als \EE{Knöchel} bezeichnet man die Knochengabel, welche aus
Schien- und Wadenbein auf Höhe des Sprunggelenks gebildet
wird. Es gibt einen Innen- und einen Außenknöchel (jeweils
körperfernes Ende von Schien- bzw. Wadenbein). Bei einer
Verletzung durch Umknicken des Fußes nach außen (häufiger)
bzw. innen kann einerseits dieser Knochen-Fortsatz brechen,
oder andererseits die Bandfixierung dieser Strukturen zu den
Mittelfußknochen reißen.

% \A{STEU: Knöchel erwähnen.}
% \A{HIRT: Review: ANA: Knöchel erwähnt, gut so?} 
% \A{GABS: Vl. könnte man Innen- und Aussenknöchel erwähnen.} 
% \A{HIRT: Review: ANA: gemacht!} 

 
\Fazit{Achtung: Unterscheide Fuß -- Bein!}

\subsection{Die Muskulatur}
 \index{Muskel}  
 \index{Muskulatur!glatte} 
 \index{Muskulatur!quergestreifte}
 \index{Muskulatur!Skelett-}
 \index{Muskulatur!Herz-}

\TextSlide{Muskelarten}
{
Prinzipiell gibt es 3  Arten von Muskelzellen:

\begin{itemize}
  \item \T{Quergestreifte (Skelett-)Muskulatur}:
  willkürlich steuerbar, trainierbar, er\-müd\-bar,  schnell und
  kräftig. Sie heißt deswegen
  \enquote{quergestreift}, da man unter dem
  Mikroskop tatsächlich eine Streifung sehen kann, allerdings nicht mit
  dem freien Auge.
  \item \T{Glatte  Muskulatur}: Nicht willentlich
  kontrollierbar, nicht ermüdbar, ausdauernd.
  Sie kommt u.\,a. in vielen
  Verdauungsorgangen (z.\,B. Darm) und den Blutgefäßen vor.
  \item \T{Herzmuskulatur}: Quergestreifter Muskel, der
  jedoch besondere Eigenschaften hat: nicht willentlich
  steuerbar, nicht ermüdbar und verfügt über eigenen
  Impulsgeber (Reizleitungssystem). Weiters ist er natürlich
  trainierbar.
\end{itemize}

}
{
\begin{itemize}
  \item Quergestreifte Muskulatur
  \item Glatte Muskulatur
  \item Herzmuskulatur
\end{itemize}
}
{}{}{}

\Layout{20}{Gegenspieler}
{%
Muskel können nur Zug und keinen
Druck ausüben, daher sind \EE{Gegenspieler} (Agonist und
Antagonist) erforderlich (z.\,B. ein Muskel streckt den Arm,
ein \E{anderer} beugt ihn).

Jeder Muskel hat einen \EE{Ursprung}
(\Speech{\protect\enquote{Wo kommt er her?}}) und einen
\EE{Ansatz} (\Speech{\protect\enquote{Wo geht er hin?}}).
Diese beiden Punkte bestimmen die Bewegung, die der Muskel
durchführen kann. Der Ursprung ist definitionsgemäß immer
der Punkt, der stillgehalten wird. Der Ansatz dagegen wird
durch den Muskelzug auf ihn zu bewegt. Normalerweise liegt
der Ursprung eines Muskels näher beim Körper, der Ansatz ist
der Punkt, der weiter entfernter ist. Für alle dieser
Erklärungen gibt es natürlich Ausnahmen, konkret die
Unterstützung der Atemmuskulatur als
\T{Atemhilfsmuskulatur}.
}
{
\begin{itemize}
  \item Agonist -- Antagonist
  \item Ursprung -- Ansatz
\end{itemize}
}
{./images/anatomie/Gray361.png}{\AuthNotesPicLegalOk{}{}Muskelzug. }{Gray's Anatomy, Copyright expired.}






%%%%%%%%%%%%%%%%%%%%%%%%%%%%%%%%%%%%%%%%%%%%%%%%%%%%%%%%%%%
%%%%%%%%%%%%%%%%%%%%%%%%%%%%%%%%%%%%%%%%%%%%%%%%%%%%%%%%%%%
%%%%%%%%%%%%%%%%%%%%%%%%%%%%%%%%%%%%%%%%%%%%%%%%%%%%%%%%%%%
%%%%%%%%%%%%%%%%%%%%%%%%%%%%%%%%%%%%%%%%%%%%%%%%%%%%%%%%%%%
\bigskip
% \clearpage
\section{Das Herz}
 \index{Herz} 
 \label{topic:herz}

\Layout{20}{\TxBeschreibung}
{%
Das Herz ist ein etwa faustgroßer\ZFootnote{Faustgröße des
\E{Patienten!}} \EE{Hohlmuskel} und hat die Funktion einer
\EE{Pumpe}. Es ist somit Teil des \E{Blutkreislaufs}. 
Es liegt \EE{hinter dem Sternum\footnote{Brustbein}}
\A{\footnote{\A{EMHO: statt Fußzeile: Brustbein (Sternum).
damit einheitlich.\par GABS: nein. Hier sollen sie bereits
primär den Begriff Sternum verwenden und daran gewöhnt
werden.}}} und teilweise im linken Brustkorb.
Unten liegt es auf \EE{dem Zwerchfell auf}.

\index{Scheidewand}  
\index{Septum}  
\index{Vorhof}
\index{Atrium}  
\index{Herz!-kammer}
\index{Ventrikel}  
Das Herz ist durch eine \EE{Scheidewand} (Septum) in eine
\E{rechte} und eine \E{linke} Seite geteilt.
Umgangssprachlich spricht man dabei jeweils vom
\emph{{\quotedblbase}rechten{\textquotedblleft}} oder
\emph{{\quotedblbase}linken Herzen{\textquotedblleft}}. Jede
Hälfte wird durch \EE{Klappen} in einen \T{Herzvorhof}
(\T{Atrium}) und eine \T{Herzkammer} (\T{Ventrikel})
geteilt.

\subparagraph{Blutversorgung}
\index{Herzkranzgefäße}
\index{Koronargefäße}

Die Blutversorgung des Herzens erfolgt über die
\T{Herzkranzgefäße} (\T{Koronargefäße}), welche direkt aus
der Aorta (Hauptschlagader) -- gleich oberhalb der
Aortenklappe beim Ausgang des linken Ventrikels -- abgehen.

\subparagraph{Schichten}

Das Herz ist in mehreren Schichten aufgebaut:
\AuthNotes{\footnote{GABS: Schichten relevant? KOCH: ja!}}
\begin{itemize}
\item Endokard (innen) \index{Endokard}     
\item \EE{Myokard} (\EE{Muskel}) \index{Myokard|EE} 
\item Epikard (außen) \index{Epikard} 
\item Perikard (ganz außen), bildet den Herzbeutel.
\end{itemize}

\noindent Im Herzbeutel befindet sich eine geringe Menge
Flüssigkeit, die eine reibungsarme Bewegung des Herzens
ermöglicht. Wenn durch Eintritt von z.\,B. Blut nach einer
Verletzung der Herzbeutel zu viel gefüllt ist, kann es
dadurch zu einer Einschränkung der Bewegung und
schlimmstenfalls zum Tod kommen.
}
{
\begin{itemize}
  \item Hohlmuskel
  \item Aufbau
  \begin{itemize}
    \item Scheidewand (Septum) - rechts / links
    \item Vorhof (Atrium)
    \item Kammer (Ventrikel)
    \item Herzklappen
  \end{itemize}
  \item: Blutgefäße: Herzkrankgefäße
  \item Schichten
  \begin{itemize}
    \item Endokard
    \item Myokard
    \item Epikard
    \item Perikard
  \end{itemize}
\end{itemize}
}
{./images/hirtler-aass/Herzvorderflaeche.pdf}{Das
Herz.}
{Hirtler}


\Layout{20}{Herzklappen}
{%
\index{Herzklappen}%
\index{Segelklappen}%
\index{Taschenklappen}%
%
\marginline{\Commentary[Herzklappen]{Die Unterscheidung von
Segel- und Taschenklappen wird nur der Vollständigkeit
halber angeführt. Interessanter, aber nicht essentiell, ist
die Nennung der Aorten- und
Mitralklappe.}}% 
Insgesamt \EE{4} \T{Herzklappen} sorgen als \EE{Ventile}
dafür, dass das Blut nur \E{in eine Richtung} strömt. Es
gibt 2 Segelklappen und 2 Taschenklappen, jeweils eine auf
jeder Seite des Herzens. Die Segelklappen befinden sich
zwischen den Vorhöfen und den Herzkammern, die
Taschenklappen befinden sich an den Ausflussöffnungen der
Kammern am Übergang zur Lungenarterie (rechte Seite) bzw.
zur  Aorta (Hauptschlagader, linke Seite).

Rechts trennt die \Z{Trikuspidalklappe (Segelklappe)} den
Vorhof von der Kammer, am Ausfluss der Kammer trennt die
\T{Pulmonalklappe} (Taschenklappe) die Kammer von der
Lungenarterie. Auf der linken Seite trennt die
\T{Mitralklappe} (Segelklappe) den Vorhof von der Kammer, am
Ausfluss trennt die \T{Aortenklappe} (Taschenklappe) den
Ventrikel von der Aorta.
}{
\begin{itemize}
  \item 4 Stück, je 2 pro Seite:\begin{itemize}
    \item Segelklappen: Vorhof → Kammer
    \item Taschenklappen: Kammer → Kreislauf
  \end{itemize}
  \item Rechts: \Z{Trikuspidalklappe,} Pulmonalklappe
  \item Links: Mitralklappe, Aortenklappe
\end{itemize}
}
{./images/hirtler-aass/Herz-Schema_edited.pdf}
{Das Herz - schematisch.}
{Hirtler}

\bigskip

\Layout{47}{Übersicht Herzklappen}{}{}{%
\begin{tabularx}{\linewidth}{XXX}\toprule\addlinespace
\noindent & \noindent\TabSub{Rechts} & \noindent\TabSub{Links}
\\\addlinespace\midrule\addlinespace
\noindent\TabHeading{Vorhof → Kammer}

\noindent\Z{(Segelklappen)}
&
\noindent\Z{Trikuspidalklappe}
&
\noindent\Z{Mitralklappe}
\\\addlinespace
\noindent\TabHeading{Kammer → Ausfluss}

\noindent\Z{(Taschenklappen)}
&
\noindent\Z{Pulmonalklappe}
&
\noindent\Z{Aortenklappe}
\\\addlinespace\bottomrule
\end{tabularx}
}
{Übersicht Herzklappen}{} 



 
% \Captionof{figure}{ Schema der
% Herz-Anatomie.  }

\TextSlide{Mechanische und elektrische Herzfunktion}{%
\label{topic:herzfunktion}%
Die Arbeit des Herzens wird durch zwei wesentliche
Funktionen gewährleistet: Die \EE{mechanische} und die
\EE{elektrische} Herzfunktion. Normalerweise arbeiten beide
Komponenten Hand-in-Hand, diese Zusammenarbeit kann aber
krankheitsbedingt gestört sein\footnote{Eine typische
Störung des Zusammenspiels von mechanischer und
elektrischer Herzfunktion ist zum Beispiel die \I{Pulslose
Elektrische Aktivität} (\I{PEA}). Hierbei sendet das
Reizleitungssystem Impulse an den Herzmuskel, der diese
jedoch ignoriert und in Folge nicht kontrahiert. Die PEA
ist ein reanimationspflichtiger, nicht-schockbarer
Rhythmus.}.

\begin{itemize}
  \item \EE{Mechanische Herzfunktion}:
  Der Herzmuskel pumpt und entspannt sich beim Erwachsenen
  in Ruhe ca. 60 bis 100~mal in der Minute. Bei jeder
  Anspannung (Kontraktion) wirft das Herz ca. 70{\Ml*} Blut
  in die Arterien aus. Dabei pumpt es ca.
  \VonBis{4}{6}\Liter*  Blut (bzw. 1\,x das gesamte
  Blutvolumen) in der Minute.
  
  Durch die Pumpfunktion des Herzens entsteht ein
  lebenswichtiger Druck im (arteriellen) Gefäßsystem, der
  (arterielle systolische) \E{Blutdruck} \index{Blutdruck}.
  \VARIANT{VAR-STANDARD}{Vgl. \FullPrettyRef{topic:blutdruck}. }
  \item \EE{Elektrische Herzfunktion}:
  Die Steuerung der mechanischen Herzfunktion erfolgt durch
  elektrische Aktionen. Dafür verfügt das Herz über ein
  eigenes \EE{Reizleitungssystem}, siehe
  \prettyref{topic:reizleitungssystem}.
\end{itemize}
}{
\begin{itemize}
  \item \EE{Mechanische}
  \begin{itemize}
    \item Pumpt \VonBis{60}{100}\PerMin* ca. 70{\Ml*}
    \item \VonBis{4}{6}\Liter* pro Minute
    \item Blutdruck
  \end{itemize}
  \item \EE{Elektrische}
  \begin{itemize}
    \item Reizleitungssystem
  \end{itemize}
\end{itemize}

}{}{}{}

%
%
\clearpage
\subsection{Reizleitungssystem}
\index{Reizleitungssystem} 
\label{topic:reizleitungssystem}



\TextSlide{\TxBeschreibung}{%
Das Herz besitzt ein eigenständiges Erregungszentrum,
welches elektrische Impulse über ein Reizleitungssystem zum
Herzmuskel leitet (\E{elektrische Herzfunktion}). Jeder
dieser Impulse wird normalerweise vom Herzmuskel mit einer
Kontraktion beantwortet (\E{mechanische Herzfunktion}).
\Commentary[Herz / Reizleitungssystem]{Das
Herzreizleitungssystem soll speziell in Hinblick auf den
AED-Teil unterrichtet werden. Wesentlich ist die
Unterscheidung von \E{mechanischer} und \E{elektrischer}
Herzfunktion, auf die nochmal nachdrücklich hingewisen
werden muss.}
}{%
\begin{itemize}
  \item Elektrische
Herzfunktion → mechanische
Herzfunktion
\end{itemize}
}{}{}{}

\Layout{20}{Aufbau und Verlauf}
{%
Der \T{Sinusknoten} \EE{(1)} ist normalerweise der
Impulsgeber (\E{natürlicher Schrittmacher}). Er feuert mit
einer Eigenfrequenz von ca.
\VonBis{60}{100}\PerMin* elektrische Impulse (Erregungen)
ab.

Eine solche Erregung wird über den rechten Vorhof zum 
\T{AV-Knoten} \EE{(2)} geleitet, dort gebündelt, und über
das \T{His-Bündel}, welches sich im Bereich der
Vorhof-Kammer-Grenze befindet, in Richtung der Kammern
weitergeleitet. Der AV-Knoten und das His-Bündel \E{filtern}
 und \E{verlangsamen} die Erregungsleitung, um zu
verhindern, dass zu viele Impulse die Herzkammern erreichen.

Im nächsten Schritt gelangen die Impulse über die
Tawara-\EE{Schenkel} in die rechte und linke Herzkammer, wo
sie über die Purkinje-\EE{Fasern} die Muskulatur erreichen.

Bei einem \E{Ausfall} eines Knotens, bzw. bei Unterbrechung
der Erregungsleitung (z.\,B. durch Gewebeschädigung nach
einem Herzinfarkt), kann die jeweils \E{nachfolgende Station
als Impulsgeber übernehmen} -- allerdings mit deutlich
langsamerer Frequenz\ZFootnote{D.\,h. jede der Stationen ist
zur eigenständigen Erregungsbildung fähig, sofern ihr nicht
die übergeordnete Station zuvor kommt.}. Im Extremfall, wenn
das gesamte Leitungssystem ausfällt, können auch die
Herzmuskelzellen selbst einen minimalen Rhythmus
($\sim$\,30\PerMin*) aufrecht erhalten.

\Fazit{Der normale Schrittmacher ist der Sinusknoten,
welcher die elektrischen Impulse für die Vorhöfe und die
Kammern erzeugt. Man spricht dann vom \T{Sinusrhythmus}.}

Beispiele von Rhythmen werden unter \prettyref{topic:herzrhythmusstoerungen}
vorgestellt. 
}
{
\begin{itemize}
  \item Abfolge:
  \begin{itemize}
    \item Sinusknoten \EE{(1)}
    \item AV-Knoten \EE{(2)}
    \item His-Bündel
    \item Tawara-Schenkel
    \item Purkinje-Fasern
  \end{itemize}
  \item Normaler Impulsgeber: Sinusknoten
  \item Nachfolgende Stationen können bei Bedarf übernehmen (HF\,\Sign{↓})
\end{itemize}

}{./images/anatomie/Reizleitungssystem_1.png}{\AuthNotesPicLegalOk{}{CC-BY}
Reizleitungssystem, Schema.}{J. Heuser,
basierend auf der Arbeit von Patrick J. Lynch; illustrator;
C. Carl Jaffe; MD; cardiologist Yale University Center for
Advanced Instructional Media. Lizenz: CC-BY}


\VARIANT{VAR-ERNAEHRUNGSPAED}{

\subsection{Messung der Herzfrequenz -- HF}
\label{topic:pulsmessung}
 
\Definition{Als \T{Herzfrequenz} (\T{HF})
bezeichnet man die Anzahl der Schläge des Herzens pro Minute.} 

Sie wird mittels Fühlen des Pulses ermittelt.
\marginpar{\E{\scriptsize Die Blutgefäße werden im Detail
unter
\prettyref{topic:blutgefaesse-wichtige} besprochen.}} 
Dieser wird idealerweise an der \EE{Radialarterie} (A.
\E{radialis}) bzw. bei nicht ansprechbaren oder schockierten
Patienten an der \EE{Halsschlagader} (A. \E{carotis})
ertastet. Bei Neugeborenen und Säuglingen kann man auch an
der \EE{Oberarmarterie} (A. \E{brachialis}) fühlen. Gemessen
wird immer mit zwei bis drei Fingern. Nie jedoch mit dem
Daumen, da hier meistens nur der eigene Puls gefühlt wird
und nicht der des Patienten. Wenn an der Halsschlagader (A.
\E{carotis}) die Herzfrequenz ermittelt wird, darf nie auf
beiden Seiten gleichzeitig gemessen werden, da sonst die
Gefahr besteht, dass es zu einer Unterversorgung des Gehirn
durch Blut kommen kann.

\Text{Durchführung}{
Es wird an der Radialarterie (A. \E{radialis}) der Puls mit
zwei bis drei Fingern ertastet. Nun werden 15 Sekunden lang
die Schläge gezählt. Multipliziert man den Messwert mit vier,
erhält man  die Herzfrequenz (HF) des Patienten. Der
Normalwert liegt beim Erwachsenen zwischen 60 und 100
Schlägen pro Minute. Dokumentiert wird sie mit einem HF
(z.\,B.: HF\,=\,70)
}{
}{}{}{}

}

\VARIANT{VAR-ERNAEHRUNGSPAED}{
\subsection{EKG}
\label{topic:ekg}
\index{EKG}

Das \E{Elektrokardiogramm} (\Abk{EKG}) zeigt mittels eines
entsprechenden Gerätes die elektrische Herzaktivität des
Reizbildungs- und Reizleitungssystems an. Dazu werden auf
dem Körper nach einem bestimmten Muster Elektroden
aufgeklebt und mittels eines Kabels mit dem \Abk{EKG}-Gerät
verbunden. Im Rettungsdienst werden fast ausschließlich
Klebeelektroden für den einmaligen Gebrauch verwendet. 

\AuthNotes{GABS: Sollte im Defi-Teil durchgenommen
werden?:} 

\Layout{22}{\Z{EKG -- Ablauf einer Erregung}}
{\Z{\begin{itemize}
    \item P-Welle $\approx$ Vorhoferregung
    \item PQ-Zeit $\approx$ AV-Überleitung
    \item QRS-Komplex $\approx$ Kammererregung
    \item ST-Strecke + T-Welle $\approx$ Erregungsrückbildung
\end{itemize}}
}
{}
{./images/wikipedia-cc-by-sa/EKG_Komplex}
{Eine elektrischer Herzaktion}
{WMC:Hank, CC-BY-SA-2.0-DE}



\Cave{Das EKG zeigt nur die \E{elektrische} Aktivität,
nicht aber die tatsächliche Muskelarbeit an! 

\EE{Es erlaubt
keine Aussage über die Auswurfleistung! }}

\Fazit{Grundsatz: \EE{Selbst ein pulsloser Patient kann ein
unauffälliges EKG haben} (allerdings nicht
lange)!}

\subsubsection{EKG -- Ableitungen}
\label{topic:ekg-ableitungen}

\TextSlide{\TxBeschreibung}{
Grundsätzlich unterscheidet man zwischen den Extremitätenableitungen
und den Brustwandableitungen. 

\begin{itemize}
    \item Extremitätenableitungen\Z{: \EKG{I}, \EKG{II},
    \EKG{III}, \EKG{aVL}, \EKG{aVR}, \EKG{aVF}}
    \item Brustwandableitungen\Z{: \VonBis{\EKG{V1}}{\EKG{V6}}}
\end{itemize}
}{
\begin{itemize}
    \item Extremitätenableitungen
    \item Brustwandableitungen
\end{itemize}
}{}{}{}

Es haben sich folgende Begriffe eingebürgert:

\begin{description}
    \item[Extremitätenableitung]
    \enquote*{\I{Rhythmusstreifen}} zur Beurteilung des
    Herzrhythmus und zum Monitoring
    \item[Extremitätenableitung\,+\,Brustwandableitung]
    (\T{12-Kanal-EKG}) -- Damit können Schäden der
    jeweiligen Herzregion zugeordnet werden. Wenn vorhanden,
    ist das 12-Kanal-\Abk{EKG}
    (\Speech{\enquote{12er}}) für die Diagnostik das Mittel der Wahl. 
\end{description}

\TextSlide{Extremitätenableitungen}
{
Die Extremitätenableitungen werden mittels vier Elektroden
abgeleitet. Je eine Elektrode wird an den Enden der
Extremitäten angebracht. Zur Vereinfachung können die
Elektroden auch am Übergang vom Rumpf zu den Extremitäten
angebracht werden (z.\,B. statt dem rechten Arm an der
rechten Schulter). Wichtig ist, dass die Elektrodenpärchen
immer auf \EE{gleicher Höhe} angebracht werden!

Aus den Informationen der vier Elektroden erzeugt das
EKG-Gerät \EE{sechs Ableitungen} \Z{(\EKG{I}, \EKG{II},
\EKG{III}, \EKG{aVL}, \EKG{aVR} und \EKG{aVF}).}
}{
\begin{itemize}
    \item 4 Elektroden an den Extremitäten
    
     (bzw. am Übergang Rumpf -- Extremitäten).
    \item 6 Ableitungen.
\end{itemize}
}{}{}{}


\begin{table}[H]
\FormatTable
\Captionof{table}{Farbe und Positionen der EKG-Elektroden
für die Extremitätenableitung.}

\begin{tabular}{lll}
\TabHeading{Elektrode} 
& 
\TabHeading{Position} 
&
\TabHeading{Alternative Position}
\\\midrule
\colorbox{red}{\sffamily\bfseries\textcolor{white}{Rot}} & Rechte Hand & Rechte
Schulter
\\
\colorbox{yellow}{\sffamily\bfseries Gelb} & Linke Hand & Linke
Schulter
\\
\colorbox{green}{\sffamily\bfseries\textcolor{white}{Grün}} & Linker Fuß  &
Linke Leiste
\\
\colorbox{black}{\sffamily\bfseries\textcolor{white}{Schwarz}} & Rechter Fuß &
Rechte Leiste \\\bottomrule
\end{tabular}
\end{table}


\begin{figure}[H]
\begin{center}
\includegraphics[width=7cm]{./images/geraete/ECGcolor.png}
\Captionof{figure}{\AuthNotesPicLegalOk{}{PD}Extremitätenableitungen
(schwarze Kabel) und Brustwandableitungen (blaue Kabel).
\SourcePicture{WmCo/Madhero88, PD}\AuthNotes{GABS: Bessere Beschriftung wäre
wünschenswert, Dringlichkeit: niedrig}}
\end{center}
\end{figure}




\TextSlide{Brustwandableitungen}
{
Die Brustwandableitungen werden mittels sechs Elektroden
abgeleitet und erlauben eine genauere Zuordnung von
Störungen zu der jeweiligen Region des Herzens. Aus den
Informationen der sechs Elektroden erzeugt das
\Abk{EKG}-Gerät \EE{sechs Ableitungen}:
\Z{\VonBis{\EKG{V1}}{\EKG{V6}}}.
}{
\begin{itemize}
    \item 6 Elektroden am Brustkorb.
    \item 6 Ableitungen.
\end{itemize}
}{}{}{}





\begin{table}[H]
\FormatTable
\Z{
\Captionof{table}{Position der EKG-Elektroden bei der Brustwandableitung.}

\begin{tabulary}{\linewidth}[H]{CL}
\TabSub{Elektrode}
 &
\TabSub{Position}
\\\midrule
\EKG{V1} & 4. Zwischenrippenraum rechts vom Sternum
\\
\EKG{V2} & 4. Zwischenrippenraum links vom Sternum
\\
\EKG{V3} & Mitte zwischen \EKG{V2} und \EKG{V4}
\\
\EKG{V4} & 5. Zwischenrippenraum links in der Verlängerung der
Schlüsselbein-Mittellinie
\\
\EKG{V5} & Mitte zwischen \EKG{V4} und \EKG{V6}
\\
\EKG{V6} & 6. Zwischenrippenraum in der mittleren Achsellinie
\\\bottomrule
\end{tabulary}
}
\end{table}


Die richtige Positionierung der \Abk{EKG}-Elektroden ist die
Voraussetzung um ein sinnvoll auswertbares \Abk{EKG} zu
erstellen. Falsch angebrachte Elektroden können ein
\Abk{EKG} derart verfälschen, dass es zu kompletten
Fehldiagnosen kommen kann\ZFootnote{In \cite{Schnelle200607}
wurde untersucht, wie sich die Anbringungsstelle der
Elektroden auf das \Abk{EKG} auswirken kann. }.


}


\clearpage % TEMP temp clearpage
\section{Der Kreislauf }
\index{Kreislauf}
\label{topic:blutkreislauf} 

\TextSlide{\TxBeschreibung}
{Der Kreislauf besteht aus dem Herz, dem Gefäßsystem und dem Blut.
Es wird der \T{Körperkreislauf} (\I[Kreislauf!großer]{Großer
Kreislauf}, \I{Hochdrucksystem}) vom
\T{Lungenkreislauf} (\I[Kreislauf!kleiner]{Kleiner Kreislauf},
\I{Niederdrucksystem}) unterschieden. Das Herz ist für 
beide Systeme die gemeinsame Pumpe.
Gesteuert wird der Kreislauf vom \EE{Stammhirn}
(Herzfrequenz, Blutdruckregulation, etc.).

\VARIANT{VAR-STANDARD}{
Der Blutdruck und die Pumpfunktion des Herzens werden unter 
\FullPrettyRef{topic:blutdruck},
und \FullPrettyRef{topic:herzfunktion}
besprochen.
}
}
{
\begin{itemize}
  \item großer Kreislauf
  \item kleiner Kreislauf
  \item Pumpe: Herz
  \item Rohrsystem: Gefäße
  \item Flüssigkeit: Blut
  \item Steuerung: Hirnstamm
\end{itemize}
}
{}{}{}

\Fazit{\EE{Kreislauf = Herz + Gefäßsystem + Blut}}

\VARIANT{VAR-ERNAEHRUNGSPAED}{
\subsection{Messung des Blutdrucks -- RR}
\label{topic:rr-messung}
\label{topic:blutdruckmessung}

\TextSlide{\TxBeschreibung}{
Der Blutdruck ist jener Druck, welcher in unseren
Blutgefäßen herrscht. Wenn wir vom
\emph{\protect\enquote{Blutdruck}} sprechen beziehen wir uns
auf den Blutdruck in den Arterien\footnote{\emph{Arterie:}
Vom Herzen wegführendes Blutgefäß.}. Er gehört in
Kombination mit der Herzfrequenz zu den grundlegenden
Vitalwerten.

\E{Der Druck in der Arterie schwankt in Abhängigkeit zum
Herzauswurf}: Wenn das Herz Blut auswirft
(\E{Auswurfphase}, \T{Systole}) erreicht er sein Maximum,
und sinkt danach in der \E{Füllungsphase} des Herzens
(\T{Diastole}) wieder ab, bis das Herz neuerlich Blut auswirft.
Dementsprechend können zwei Werte erhoben werden: Der

\begin{itemize}
  \item \T{Systolische Blutdruck} (Spitzendruck), und der
  \item \T{Diastolische Blutdruck} (Minimaldruck).
\end{itemize}


Im Normalfall
werden beide Werte ermittelt.
% Der \EE{systolische Wert} (\protect\enquote{obere Wert}) ist jener Druck
% im Blutgefäßsystem, der  während der \E{Auswurfphase} des
% Herzens, in der sich die Herzkammern anspannen,
% zusammenziehen und Blut auswerfen, herrscht. Der
% \EE{diastolische Wert} (\protect\enquote{untere Wert}) ist jener Druck im
% Blutgefäßsystem, der während der \E{Entspannungsphase} des
% Herzens, in der sich die Herzkammern mit Blut füllen,
% herrscht.
Die \EE{Einheit}, in der der Blutdruck angegeben wird, ist
\T{Millimeter Quecksilbersäule}
(\T{mm\,Hg}).\index{Millimeter Quecksilbersäule}\index{mm\,Hg}\footnote{\ce{Hg}
ist das chemische Symbol für Quecksilber.}
}{
\begin{itemize}
  \item Druck in den Arterien
  \begin{itemize}
    \item systolischer Wert (Auswurfphase des Herzens)
    \item diastolischer Wert (Entspannungsphase des Herzens)
  \end{itemize}
  \item Einheit \MmHg
\end{itemize}
}{}{}{}

\TextSlide{Arten der Blutdruckmessung}{%
Der Blutdruck kann auf verschiedene Arten gemessen werden.
Üblich ist die unblutige, \E{indirekte Methode}, mittels
einer an der Extremität angelegten
Druckmanschette\ZFootnote{Mittels einer Druckmanschette wird
gemessen, bei welchem Druck ein mit einem Stethoskop
hörbarer, durch Pulstasten fühlbarer oder durch sonstige
technische Hilfsmittel messbarer Blutfluss stattfindet.
Dies erlaubt Rückschlüsse auf den Druck im Blutgefäß.}. Dazu
gibt es unterschiedliche Messgeräte, grundsätzlich
unterscheidet man zwischen elektronischen und rein
mechanischen Messgeräten.

\EE{Elektronische Messgeräte}  werden oft von Patienten zur
Selbstkontrolle des Blutdrucks verwendet. Diese Geräte sind
jedoch vergleichsweise fehleranfällig und sind für einen
professionellen Einsatz nicht geeignet.

Für den professionellen Einsatz gibt es spezielle
elektronische Blutdruckmessgeräte, welche zum Beispiel in
Überwachungsmonitoren zum Einsatz kommen. Allerdings sind
auch diese Geräte mitunter fehleranfällig, wenn die
Messbedingungen nicht optimal (z.\,B. bei Erschütterungen)
sind.

In der Praxis wird die Blutdruckmessung mit \EE{mechanischen
Messgeräten} durchgeführt. Dabei wird eine aufblasbare
Manschette mit angeschlossenem Dosenmanometer und Pumpballen
mit einem Ablassventil benötigt. Im Idealfall wird auch noch
ein Stethoskop verwendet.

\begin{center}
     
%   \includegraphics[width=4cm]{./images/noauth/GER-rr-elektrisch-oberarm.png}
%   
%    \includegraphics[width=4cm]{./images/noauth/GER-rr-elektrisch-unterarm.jpg}
%    
%    \includegraphics[width=4cm]{./images/noauth/GER-rr-messgeraet.jpg}
\TODO{GABS}{2010-04-27}{Blutdruckmesser: Bilder aus Serie
Pallinger}{}
\end{center}
}{
\begin{itemize}
  \item unblutige, indirekte Messung
  \begin{itemize}
    \item mechanische Messung
    \item elektronische Messung
  \end{itemize}
  \item invasive Messung
\end{itemize}
}
{}
{}
{}


\subsubsection{Manuelle Blutdruckmessgeräte}
\Text{Manuelles Blutdruckmessgerät}{%
Ein manuelles Blutdruckmessgerät besteht aus einer
Manschette, einem Schlauch, einem Auslassventil, einem
Pumpballen und einer Druckanzeige.

Es gibt Manschetten in unterschiedlichen Größen.
}
{}
{}
{}
{}

\Layout{37}{Abbildung}
{}
{}
{./images/pallinger-christoph-aass/Blutdruck_32768_export_01024.jpg}
{Blutdruckmanschetten mit Hak- und Klettverschluss, sowie
einem Doppel- und einem Einzelkopfstethoskop.}
{Ch. Pallinger}

\subsubsection{Exkurs: Stethoskop}
\label{topic:stethoskop}
\index{Stethoskop}

\Text{\TxBeschreibung}
{Ein Stethoskop ist eine Hörvorrichtung. Es besteht
patientenseitig aus einem Kopf mit einer Membran oder einem
Trichter, einem Schlauch und untersucherseitig aus einem
Bügel mit zwei Oliven (\protect\enquote{Ohrenstöpsel}).
Die Oliven sollen \E{nach vorne} gerichtet sein, da auch der
äußere Gehörgang nach vorne verläuft.
}{}{}{}{}

\Text{Funktionen}
{%
Manche Typen verfügen über einen zweiseitigen Kopf. Meistens
haben sie sowohl eine Membran, als auch einen
Schalltrichter.~\ZFootnote{Der Unterschied besteht in einer
unterschiedlichen Verstärkung von hohen oder tiefen
Frequenzen, was für den fortgeschrittenen Anwender beim
Abhören von Herz oder Lunge eine Rolle spielt.
Teurere Modelle verfügen mitunter über eine
Dual-Frequency-Membran, mit welcher man durch den
Auflegedruck wahlweise hohe oder tiefe Frequenzen verstärken
kann.} Durch Drehen am Kopf kann man die gewünschte Seite
einstellen.
}{}{}{}{}

\Text{Funktionsprüfung}
{%
Vor der Verwendung muss mittels Beklopfen der Membran oder
des Schalltrichters geprüft werden, ob der Kopf des
Stethoskops auf die gewünschte Seite eingestellt ist und ob
die Schallleitung funktioniert.

}{}{}{}{}
}

\subsection{Abschnitte des Kreislaufs}
\index{Gasaustausch} 
\index{Kreislauf!Abschnitte}

\begin{WidePar}
\Parallel
{\subsubsection{Kleiner (Lungen-)Kreislauf}
Die obere und untere Hohlvene münden in den \EE{rechten
Vorhof} des Herzens und liefern sauerstoffarmes Blut aus
dem Körper an. Von hier wird das Blut in die \EE{rechte
Herzkammer} gepumpt, um dann in die \EE{Lungenarterie} zu
gelangen. In der Lunge verästeln sich die Lungenarterien
zu haardünnen Gefäßen, den \EE{Kapillaren}. Diese
umspannen die Lungenbläschen. Es findet ein
\EE{Gasaustausch} statt: Sauerstoff wird aufgenommen und
Kohlendioxid abgegeben. Das jetzt mit Sauerstoff
angereicherte Blut wird über die \EE{Lungenvenen} in den
\EE{linken Vorhof} geleitet, der wiederum das Blut in die
\EE{linke Herzkammer} pumpt.

\vspace{1em}

\subsubsection{Großer (Körper-)Kreislauf}
Das Blut wird aus der linken Herzkammer ausgestoßen und
gelangt über die Hauptschlagader
(\II{Aorta}) in den
Körperkreislauf. Aus der Aorta entspringen \EE{Arterien}
({\quotedblbase}Schlagadern{\textquotedblleft}), die sich
immer weiter verzweigen und an Durchmesser abnehmen. Sie
versorgen den gesamten Körper mit Blut. Der durch die
Pumpleistung des Herzens erzeugte Puls ist an oberflächlich
gelegenen Arterien tastbar (Radialis-Puls, Carotis-Puls,
...). Daher rührt auch der Name
\protect\enquote{Schlagader}.

Die Arterien verzweigen sich, weiter an Durchmesser
abnehmend, ähnlich dem Geäst eines Baumes, bis hin zu
\EE{Arteriolen} und haardünnen Gefäßen, den \EE{Kapillaren}.
Hier findet wieder ein \EE{Gasaustausch} statt, diesmal
allerdings spiegelbildlich zum Lungenkreislauf:
\EE{Sauerstoff wird an das Gewebe abgegeben, im Gegenzug
wird Kohlendioxid (\ce{CO2}) aufgenommen und
abtransportiert}.

Das nun sauerstoffarme (und \ce{CO2}-reiche) Blut gelangt
über kleine \EE{Venolen} in das Venensystem und fließt in
den \EE{Venen} herzwärts. Am Weg zum Herzen nehmen die Venen
an Durchmesser immer mehr zu, bis sie letztendlich in die
großen Venen ({\quotedblbase}\EE{obere} und \EE{untere
Hohlvene}{\textquotedblleft}, Vena cava \Z{superior und
inferior}) münden.
}
{
\includegraphics[width=\linewidth]{./images/hirtler-aass/Kreislauf-Schema.pdf}

\Captionof{figure}{\AuthNotesPicLegalOk{}{Blutkreislauf, Schema.}Der
Blutkreislauf. \SourcePicture{Hirtler}}



\Fazit{Rechter Vorhof ${\rightarrow}$ rechte Herzkammer
${\rightarrow}$ Lungenarterie(n) ${\rightarrow}$
Lungenkapillaren (\ce{CO2}-Abgabe, \ce{O2}-Aufnahme) ${\rightarrow}$ Lungenvenen ${\rightarrow}$ linker Vorhof 
${\rightarrow}$ linke Herzkammer ${\rightarrow}$ Aorta ${\rightarrow}$
abzweigende Arterien ${\rightarrow}$ Kapillaren
(\ce{O2}-Abgabe, \ce{CO2}-Aufnahme)  ${\rightarrow}$ Venen
${\rightarrow}$ Hohlvene ${\rightarrow}$ Rechter Vorhof
 ${\rightarrow}$ alles beginnt wieder von vorne
 $\circlearrowright$}}
\end{WidePar}


\clearpage % TEMP temp clearpage
\subsection{Die Blutgefäße}
\index{Arterien}
\index{Venen}
\index{Kapillargefäße}

\TextSlide{Arten von Blutgefäßen}
{
Man unterscheidet zwischen:

\begin{itemize}
    \item \T{Arterien} führen \E{vom Herzen \EE{weg}} und
    haben eine dicke
    % GABS, 2010/05/04 Arterien und Muskelschicht: Erwähnung der dickeren
    % Muskelschicht ist wichtig, weil PRüfungsfrage.
    Wand und Muskelschicht. 
    \item \T{Venen} führen \E{zum Herzen \EE{hin}} und
    besitzen Venen\E{klappen}, um Rückfluss\footnote{Rückfluss:
    Fluss entgegen der Flussrichtung} zu vermeiden.     
    Venen haben eine deutlich dünn\-ere Wand als die
    muskulösen Arterien.     
    \item \T{Kapillargefäße} sind dünne \protect\enquote{Haargefäße},
    in denen der \E{\EE{Gasaustausch}} stattfindet; sie \E{verbinden Arterien
    und Venen}. 
\end{itemize}
}
{
\begin{itemize}
  \item Arterien
  \item Venen
  \item Kapillaren
\end{itemize}
}{}{}{}

\TextSlide{Sonderfall Shunt}
{%
\label{topic:shunt}
\index{Dialyse-Patient!Shunt}
Ein \T{Shunt} \index{Shunt}  ist eine Verbindung zwischen
Arterien und Venen. Er kann künstlich (chirurgisch)
angelegt, angeboren oder
akzidentiell\TODO{GABS}{2011-03-07}{akzidentiell erklären}{}
entstanden sein. Bei \EE{Dialyse-Patienten}  wird oft
künstlich -- mittels einer Operation -- ein Shunt
gelegt\ZFootnote{Dadurch dass durch die Vene umgeleitetes,
arterielles, Blut mit höherem Druck fließt, verdickt sich
die Vene und man kann sie mit dickeren Kanülen (wie sie für
die Dialyse benötigt werden) punktieren.}.
Ein derartger Dialyse-Shunt ist \Ca{verletzlich}, daher darf
an einem Arm mit Shunt \EE{nicht gestaut} werden; das
betrifft insbesonders die \E{Blutdruckmessung} (vgl.
\FullPrettyRef{topic:blutdruckmessung}).

\bigskip
\bigskip

}{%
\begin{itemize}
  \item Künstliche Verbindung Arterie -- Vene
  \item Dialyse
  \item \Ca{Verletzlich}
  \item Nicht stauen!
\end{itemize}
}{}{}{}


\TextSlide{Blutgefäße sind in Schichten aufgebaut}
{%
Mit Außnahme der dünnsten Gefäße (Kapillaren) sind alle
Gefäße annähernd gleich und \EE{in mehreren Schichten}
aufgebaut. Von innen nach außen gibt es drei Schichten:

\Z{
\begin{itemize}
  \item Tunica intima: Grenze zwischen Blut und Blutgefäß,
  enthält elastische Fasern (bei Arterien dicker) und
  längsgerichtete Muskelfasern.
  \item Tunica media: besteht häuptsächlich aus ringförmigen
  Muskelfasern, sie ist bei Arterien dicker als bei Venen
  \item Tunica adventitia: lockeres Bindegewebe
\end{itemize}
}

\Z{Die größten Gefäße, wie zum Beispiel die Aorta, haben zur
geeigneten Blutversorgung ihrer Wand dort eigene Gefäße
\Z{(Vasa vasorum)}. Arterien die herznah sind, besitzen in
ihrer Wand eher elastische Fasern. Arterien die der
Organversorgung dienen sind muskulärer aufgebaut.}
\Commentary[Blutgefäße / Aufbau / Schichten]{Der Erwähnung
des schichtweisen Aufbaus der Blutgefäße ist wichtig für das
Verständnis eines Aneurysmas.}
}{
\begin{itemize}
  \item 3 Schichten
  \item Unterschied Arterien -- Venen
  \item elastische und muskuläre Arterien
\end{itemize}
}
{}{}{}





\Layout{57}{}{}{}
{\noindent\begin{tabulary}{\linewidth}{lllL}
\addlinespace\toprule\addlinespace
\noindent\TabHeading{Gefäß}
&
\noindent\TabHeading{Fachbegriff}
&
\noindent\TabHeading{Typ}
&
\noindent\TabHeading{Bemerkungen}
\\\midrule\addlinespace
\EE{Hauptschlagader} \index{Hauptschlagader}
&
\T{Aorta} \index{Aorta}
&
Arterie
&
\protect\enquote{Hauptarterie}, geht von der linken
Herzkammer ab, von dort zweigen alle
anderen Körperarterien ab.
% GABS: Aorta steht schon in der Fachbegriff-Spalte.
\\\addlinespace
\EE{Halsschlagader} \index{Halsschlagader}
&
\T{Karotis} \index{Karotis}
&
Arterie
&
1 Paar (2 Stk.), versorgen den Kopf und das Gehirn, Puls tastbar seitlich vom
Kehlkopf
\\\addlinespace
\EE{Herzkranzgefäße} \index{Herzkranzgefäße}
&
\T{Koronararterien} \index{Koronararterien}
&
Arterie
&
Versorgen das Herz mit Blut, gehen aus der Aorta direkt über ihrer Herzklappe ab
\\\addlinespace
\EE{Lungenarterie} \index{Lungenarterie}
&
\T{Pulmonalarterie} \index{Pulmonalarterie}
&
Arterie
&
\noindent\EE{\ce{O2}-arm.} Geht von der rechten Herzkammer weg in
die Lunge zu den Lungenkapillaren
\\\addlinespace
\EE{Oberarmarterie} \index{Oberarmarterie}
&
Brachialarterie
&
Arterie
&
Blutversorgung des Armes, Taststelle zum Puls-fühlen beim Baby, Blutdruckmessung
\\\addlinespace
\EE{Radialarterie} \index{Radialarterie}
&
\T[Arteria!radialis]{A. radialis} 
&
Arterie
&
Verläuft am Unterarm speichenseitig (dort wo der Daumen ist). Puls
fühlbar.
\\\addlinespace
\EE{Leistenarterie} \index{Leistenarterie}
&
\Z{A. femoralis} \index{Arteria!femoralis}
&
Arterie
&
verläuft am Oberschenkel, Puls tastbar in der Leistenbeuge, wird für
Herzkatheter als Zugang verwendet. 
\\\addlinespace
\EE{Obere Hohlvene} \index{Obere Hohlvene}
\index{Hohlvene!obere}
&
\T{Vena cava} sup. \index{Vena cava}
&
Vene
&
Obere Vene die in den rechten Vorhof mündet
\\\addlinespace
\EE{Untere Hohlvene}
\index{Untere Hohlvene}
\index{Hohlvene!untere}
&
\T{Vena cava} inf. \index{Vena cava}
&
Vene
&
Untere Vene die in den rechten Vorhof mündet.
\\\addlinespace
\EE{Lungenvene} \index{Lungenvene}
&
\T{Pulmonalvene} \index{Pulmonalvene}
&
Vene
&
\EE{\noindent\ce{O2}-reich.} Kommt von den Lungenkapillaren und
mündet in den linken Vorhof
\\\addlinespace
\EE{Pfortader} \index{Pfortader}\label{topic:pfortader}
&
\Z{Vena portae}
&
Vene
&
Führt \ce{O2}-armes, aber nährstoffreiches Blut von Teilen des Darms zur
Leber (zur Entgiftung)
\\\addlinespace
\EE{Haargefäße} \index{Haargefäße}
&
\T{Kapillargefäße} \index{Kapillargefäße}
&
Kapillare
&
Dünne, dicht verzweigte Gefäße, Verbindung zwischen Arterien
und Venen. Hier findet der Gas- und Nährstoffaustausch statt. Es gibt
Körperkapillaren (Körperkreislauf) und Lungenkapillaren
(Lungenkreislauf)
\\\addlinespace\bottomrule
\end{tabulary}
% \end{minipage}

% STEU: Leistenarterie! Puls
% GABS: Wozu?
}
{Übersicht wichtiger Blutgefäße\label{topic:blutgefaesse-wichtige}}
{}

\clearpage % TEMP temp clearpage
% \bigskip

\subsection{Das Blut}\index{Blut}\label{topic:blut}

\dictum[Goethe]{Blut ist ein ganz besond'rer Saft.}

\TextSlide{\TxBeschreibung}
{%
Das Blut macht ca. \EE{\VonBis{7}{8}\,\PC* des Körpergewichtes}
aus, das entspricht in etwa \EE{\VonBis{5}{6}\Liter*} beim
Erwachsenen. Es setzt sich aus verschiedenen Bestandteilen
zusammen und hat weitreichende und wichtige Aufgaben. 
Das Blut besteht aus
\EE{festen} Bestandteilen und \EE{flüssigen Bestandteilen}
(Blutplasma).
Zu
den festen Bestandteilen gehören die \E{roten
Blutkörperchen} (Erythrozyten), die \E{weißen
Blutkörperchen} (Leukozyten) und die \E{Blutplättchen}
(Thrombozyten).
Das Blutplasma besteht zum größten Teil aus \E{Wasser}.
Außerdem enthält es \E{Proteine} und andere \E{gelöste
Stoffe}.
}{
\begin{itemize}
  \item \VonBis{5}{6} Liter
  \item Feste Bestandteile (Zellen, Blutkörperchen)
  \item Flüssige Bestandteile (Plasma)
  \item \prettyref{tab:blut-bestandteile}
\end{itemize}
}{}{}{}



\TextSlide{Aufgaben}{%
Die wichtigste Aufgabe des Blutes ist der \EE{Stoff- und
Gastransport}. Es ist z.\,B. der \E{Sauerstofflieferant} des
Körpers. Es holt sich den Sauerstoff in der Lunge ab, bringt
ihn zu den Zellen und nimmt auf dem Rückweg das abzuatmende
\E{Kohlendioxid} wieder mit. Weiters transportiert es die
über den Darm aufgenommenen \E{Nährstoffe} zu den
entsprechenden Weiterverarbeitungsstellen und sorgt für den
Abtransport von \E{Abfallprodukten} hin zu Leber oder Niere.

Eine weitere wichtige Aufgabe ist die \EE{Regulation der
Körpertemperatur} gemeinsam mit den Blutgefäßen. Bei Kälte
werden die Blutgefäße in den Extremitäten enger gestellt --
weniger Wärme geht nach außen verloren, das Blut und
dementsprechend die Wärme wird
\protect\enquote{zentralisiert}. Andersrum funktioniert es
bei großer Hitze. Über eine Weitstellung der Gefäße versucht
der Körper soviel Hitze wie möglich loszuwerden.

Zu guter Letzt sorgt das Blut auch noch für einen
ausgeglichenen \EE{Säure-Basen-Haushalt} im Körper -- der pH
des Körpers sollte immer zwischen 7,35 und 7,45 liegen. Bei
Abweichungen sorgen \E{Puffersysteme}
\VARIANT{VAR-STANDARD}{(\prettyref{topic:bikarbonatpuffer})}
im Blut für eine Wiederherstellung dieses pH-Wertes.
}{
\begin{itemize}
  \item Sauerstofftransport zu den Körperzellen
  \index{Sauerstofftransport}
  \item \ce{CO2}-Abtransport 
  \index{CO2-Abtransport@\ce{CO2}-Abtransport}
  \item Transport von Nährstoffen und Abfallprodukten
  \item Wärmeverteilung
  \item pH-Pufferungssystem (Säure/Basen-Gleichgewicht)
\end{itemize}
}{}{}{}

\Layout{57}
{}
{}
{}
{\begin{tabularx}{\linewidth}{XXXXXX}
\hline\addlinespace
\multicolumn{6}{c}{\TabHeading{Blut}}
\\\addlinespace\cmidrule(lr){1-3}\cmidrule(lr){4-6}\addlinespace
\multicolumn{3}{c}{\TabSub{Feste Bestandteile}}
&
\multicolumn{3}{c}{\TabSub{Flüssige Bestandteile}}
\\
\multicolumn{3}{c}{$\sim$\,42\,\PC*}
&
\multicolumn{3}{c}{$\sim$\,58\,\PC*}
\\\addlinespace\cmidrule(lr){1-1}\cmidrule(lr){2-2}\cmidrule(lr){3-3}\cmidrule(lr){4-4}\cmidrule(lr){5-5}\cmidrule(lr){6-6}\addlinespace
\TabSub{Erythrozyten}
&
\TabSub{Leukozyten}
&
\TabSub{Thrombozyten}
&
\TabSub{Wasser}
&
\TabSub{Proteine}
&
\TabSub{Gelöste Stoffe}
\\
\noindent{\scriptsize\Z{($\sim$\,4\,000\,000\,/\,{\Mul})}}
&
\noindent{\scriptsize\Z{($\sim$\,4\,000\,--\,8\,000\,/\,{\Mul})}}
&
\noindent{\scriptsize\Z{($\sim$\,150\,000\,--\,300\,000\,/\,{\Mul})}}
&
\noindent$\sim$\,90\,\PC*
&
\noindent$\sim$\,8\,\PC*
&
\noindent$\sim$\,2\,\PC*

{\scriptsize%
Ionen,

Glukose, 

Hormone,

Kreatinin, 

Harnstoff,

{\ldots}%
}
\\\addlinespace\hline
\end{tabularx}}
{Übersicht: Die Bestandteile des
Blutes.\label{tab:blut-bestandteile}}
{}











\TextSlide{Feste (zelluläre) Bestandteile}
{
Die festen Bestandteile des Blutes beinhalten rote Blutkörperchen
(Erythrozyten), weiße Blutkörperchen (Leukozyten) und Blutplättchen
(Thrombozyten).
\begin{itemize}
  \item
  \label{topic:erythrozyten}%
  \label{topic:haemoglobin}
  \index{Erythrozyten}%
  \index{Rote Blutkörperchen}%
  \index{Blutkörperchen!rote}%
  \index{Blutkonserve!Erythrozytenkonzentrat}
  \index{Erythrozytenkonzentrat}
  Die \StepHeading{roten Blutkörperchen} (\I{Erythrozyten}) sind quasi die wichtigsten des
  ganzen Körpers -- ohne ihren Sauerstofftransport
  funktioniert im Körper gar nichts. Das Protein, das für die
  Bindung des Sauerstoffes an die Erythrozyten verantwortlich
  ist, heißt \T{Hämoglobin}. Es ist auch der
  \E{Blutfarbstoff}. Ein Mangel an Hämoglobin wird als
  \T{Anämie} bezeichnet.
  
  Die Erythrozyten sind außerdem die häufigsten festen
  Blutbestandteile (4\,000\,000\,/\,{\Mul}) und tragen die
  \EE{Blutgruppenmerkmale} (A, B bzw. 0 bei Fehlen der
  Merkmale) und die Rhesusfaktoren. 
  \T{Erythrozytenkonzentrate}
  (\protect\enquote{Ery-Konzentrate}) sind
  \EE{Blutkonserven}, welche fast ausschließlich Erythrozyten
  enthalten.
  
  \Z{Das Kennzeichen der Erythrozyten ist die Abwesenheit des
  Kernes\ZFootnote{Der Kern der Erythrozyten geht bei der
  Reifung verloren.} bzw. jeglicher Zellorganellen. Nach einer
  Lebensdauer von ca. 100 Tagen werden die
  \protect\enquote{Senioren} der roten Blutkörperchen in der
  Milz abgebaut.}
  \item \label{topic:leukozyten}%
  \index{Leukozyten}%
  \index{Weiße Blutkörperchen}%
  \index{Blutkörperchen!weiße}%
  \index{Immunsystem!Leukozyten}
  Die \StepHeading{weißen Blutkörperchen} (\I{Leukozyten}) sind ein wichtiger \EE{Teil des Immunsystems} und
  dienen der Infekt- und Fremdkörperabwehr des Körpers.
  
  \Z{Es befinden sich
  zwischen 4\,000 und 8\,000 Leukozyten pro {\Mul} im Blut
  des Körpers; es gibt verschiedene Arten von Leukozyten\ZFootnote{\T{Granulozyten}: Sie bilden den Beginn der
  Immunabwehr.
  % Es   gibt je nach Erkrankung verschiedene
  % \protect\enquote{Spezialisten} unter den Granulozyten:
  %   neutrophile G. (bei Keimen und Fremdkörpern), eosinophile G. (bei Allergien
  %   und Parasiten) und basophile G.
  \T{Lymphozyten}: In diese Gruppe fallen
  Antikörper-bildende Zellen.
  \T{Monozyten}: Sie Können sich durch zelleigene
  Bewegungen im Gewebe   bewegen und Keime bzw. Fremdkörper
  verdauen.}}
  \item %
  \label{topic:thrombozyten}%
  \index{Thrombozyten}%
  \index{Thrombozytenkonzentrat}%
  Die \StepHeading{Blutplättchen} (\I{Thrombozyten}) sind kleine kernlose Scheibchen. Sie
  dienen der \EE{Blutgerinnung} und dem Wundverschluss.
  \Z{Ihre Anzahl wird mit 150\,000 bis 300\,000 pro {\Mul}
  beziffert.} \T{Thrombozytenkonzentrate} sind
  \EE{Blutkonserven}, welche fast ausschließlich Thrombozyten
  beinhalten.
\end{itemize}
}
{
\begin{itemize}
  \item Rote Blutkörperchen (Erythrozyten)
  \begin{itemize}
    \item \ce{O2}-Transport
    \item Farbstoff Hämoglobin (\ce{O2}-Transport)
    \item Blutgruppenmerkmale
    \item Erythrozyten-Konzentrate
  \end{itemize}
  \item Weiße Blutörperchen (Leukozyten)
  \begin{itemize}
    \item Abwehr-/Immunsystem, Entzündungszellen
  \end{itemize}
  \item Blutplättchen (Thrombozyten)
  \begin{itemize}
    \item Zelluläre Blutstillung
    \item Wundverschluss
    \item Thrombozytenkonzentrat
  \end{itemize}
\end{itemize}
}{}{}{}





\TextSlide{Die flüssigen Bestandteile --
Das Plasma}
{%
\index{Plasma}%
Das Blutplasma besteht zum größten Teil aus \EE{Wasser}. Der
restliche Anteil enthält \EE{Proteine}, \EE{Enzyme},
\EE{Hormone}, \EE{Elektrolyte}, \EE{Glukose} und \EE{Fette}.
Die wesentliche Aufgabe des Plasmas ist der \EE{Transport}
von Nährstoffen an ihren Zielort im Körper. Es enthält
weiters die Bestandteile für die plasmatische Blutgerinnung
und das Immunsystem (Antikörper, \Z{Immunglobuline}). Durch
die Elektrolyte ist es am Wasserhaushalt beteiligt und sorgt
mit Puffersubstanzen für einen ausgeglichenen pH-Wert im
menschlichen Körper.
Schlußendlich sorgt es auch für die Wärmeverteilung im
Körper.
\SlideCompensationNoParskip{2}
}
{
\begin{itemize}
    \item 10\PC* Proteine, Enzyme, Hormone, Elektrolyte,
    Glukose, Fette
    \item 90\PC* Wasser
    \item Aufgaben:
    \begin{compactitem}
        \item Nährstofftransport
        \item Wärmeverteilung
    	\item Proteine und Enzyme für \EE{Blutgerinnung} (Serum)
    	\item Pufferung des pH-Wertes
    \end{compactitem}
\end{itemize}
}{}{}{}

\subsection{Die Blutstillung}
\index{Blutstillung} 

\Layout{60}{\TxBeschreibung}{%
Unser Organismus verfügt über eine körpereigene
Blutstillung, welche aktiviert wird, sobald es im Körper zu
einer Verletzung kommt. In einem ersten Schritt
\glqq{}verkleben\grqq{} Blutblättchen an der verletzten
Stelle des Blutgefäßes  und bilden so einen
\EE{Thrombozytenpfropfen}, der die Wunde grob verschließt
(vorläufiger Wundverschluss, \Z{primäre Blutstillung}). Erst
in einem zweiten Schritt kommt es durch die
\EE{Blutgerinnung} zum endgültigen Wundverschluss
(\Z{sekundäre Blutstillung}). Für die Bildung der
Thrombozytenpfropfen sind die festen Bestandteile des
Blutes und für die Blutstillung das Blutplasma (flüssiger
Bestandteil) zuständig. Am Ende der Wundheilung werden
defekte Zellen abgebaut und Bindegewebe aufgabaut.

\Z{Bei der
Blutgerinnung werden nacheinander mehrere im Blutplasma
gelöste Gerinnungsfaktoren \label{topic:gerinnungsfaktoren}
aktiviert. Dadurch verwandelt sich das zunächst ebenfalls im
Plasma gelöste Fibrinogen in einen klebrigen und später in
einen festen, die Wunde verschließenden Blutpfropfen
(Thrombus).} Bei Patienten mit Bluterkrankheit (Hämophilie)
fehlen Gerinnungsfaktoren und es kommt daher
nicht zur Blutstillung.

\Z{Das Gegenstück zur Blutgerinnung, bei der es zur Bildung
von Thromben kommt, ist die Fibrinolyse, bei der es zur
Auflösung von Thromben kommt. Normalerweise sind
diese beiden Mechanismen im Gleichgewicht, sodass das Blut
weder zu schnell noch zu langsam gerinnt.}
% Die Blutstillung umfasst alle Vorgänge, die zwischen dem Entstehen und dem
% Verschluss einer Wunde ablaufen. Sie läuft in zwei Phasen ab: vorläufiger
% Wundverschluss durch Bildung eines \E{Thrombozytenpfropfens} 
% \index{Thrombozytenpfropf} und endgültiger Wundverschluss
% durch Blutgerinnung\index{Blutgerinnung}. 
}
{
\begin{itemize}
  \item vorläufiger Wundverschluss (Thrombozytenpfropf)
  \item endgültiger Wundverschluss (Blutgerinnung)
\end{itemize}
}
{./images/anatomie/Bleeding_finger.jpg}
{}
{\SourcePicture{Crystal
(Crystl) from Bloomington, USA (\url{http://www.flickr.com/people/crystalflickr/}),
Lizenz: CC-BY}}



%\paragraph{Bildung eines Thrombozytenpfropfens}

% Zunächst lagern sich an die verletzte und defekte Stelle des Blutgefäßes Thrombozyten an
% und verkleben miteinander. Die verklebenden
% Blutplättchen bilden einen Thrombozytenpfropf.

% \Fazit{Die Bildung des Thrombozytenpfropfes wird von den zelluären Bestandteilen
% des Blutes erledigt.}


% \TextSlide{Blutgerinnung}
% {%
% \Z{Die Blutgerinnung ist der wichtigste
% Prozess bei der Blutstillung. Sie wird entweder durch die
% Gewebeverletzung selber oder durch ein eigenes System von Gerinnungsfaktoren
% ausgelöst. Der Vorgang der Blutgerinnung hat zum Ziel,
% das im Blutplasma gelöste \Z{Fibrinogen} zu aktivieren und
% in einen gallertartigen, klebrigen und später festen Zustand
% zu überführen. Um dieses Ziel zu erreichen, steht dem Organismus ein
% Gerinnungsystem bestehend aus im Plasma gelösten Gerinnungsfaktoren zur
% Verfügung, das kaskadenartig aufgebaut ist. Dieses System funktioniert wie
% eine Kette aufgestellter Dominosteine, die nacheinander umfallen, wenn der
% erste Stein angestoßen wird.}
% 
% Neben dem System der Blutgerinnung
% gibt es auch ein System zur Auflösung von Thromben, die Fibrinolyse. Beide
% Systeme müssen sich die Waage halten.

% \Fazit{Die Blutgerinnung kommt durch im Plasma gelöste Gerinnungsfaktoren
% \label{topic:gerinnungsfaktoren} zustande. An der Blutstillung sind sowohl die
% festen Bestandteile des Blutes als auch das Plasma beteiligt.}
% }
% {
% \begin{itemize}
%   \item Thrombozyten
%   \item Fibrinogen
%   \item Gerinnungsfaktoren
%   \item Fibrinolyse
% \end{itemize}
% }
% {}{}{}

\clearpage

\section{Das Nervensystem}
 \index{Nervensystem} 


\TextSlide{Unterteilung}
{%
Man kann das Nervensystem nach seinem Aufbau und nach seiner
Funktion unterteilen.
Vom Blickwinkel des \underline{Aufbaus} wird das Nervensystem grob
unterteilt in ein

\begin{itemize}
\item \T{Zentrales Nervensystem} (\T{ZNS}) und ein 
\item \T{Peripheres Nervensystem}(\T{PNS}).
\end{itemize}

\noindent Das ZNS ist im Wesentlichen für \E{Steuer-, Regel-
und Denkfunktionen} zuständig, sowie teilweise für die
Weiterleitung von Impulsen und Informationen. Zum ZNS
gehören \EE{Gehirn} und \EE{Rückenmark}.

Das periphere Nervensystem ist fast ausschließlich für die
\E{Weiterleitung von Informationen} und Impulsen zu den
Endorganen (oder umgekehrt, von den Endorganen zum ZNS)
zuständig. Zum peripheren Nervensystem gehören alle peripheren Nerven.
Betrachtet man die \underline{Funktion} kann man das
Nervensystem in ein
\begin{itemize}
  \item \T[Nervensystem!motorisches]{Motorisches}
  (willkürliche\footnote{\emph{Willkür}: Dem \E{freien
  Willen} unterworfen.} Bewegung),
  \item \T[Nervensystem!sensorisches]{Sensorisches}
  (\protect\enquote*{fühlend}, Reizweiterleitung von den
  Sinnesorganen) und ein
  \item \T[Nervensystem!vegetatives]{Vegetatives}
  (autonomes, bzw.
  unwillkürliches) Nervensystem
\end{itemize}
unterteilen. 
}
{
Unterteilung nach Aufbau
\begin{itemize}
  \item zentrales Nervensystem
  \item peripheres Nervensystem
\end{itemize}
Unterteilung nach Funktion
\begin{itemize}
  \item motorisches Nervensystem
  \item sensorisches Nervensystem
  \item vegetatives Nervensystem
\end{itemize}
Die Begriffe sind recht unscharf definiert und geben nur
eine grobe, aber wichtige Übersicht
wieder.
}
{}{}{}

\subsection{Das Zentralnervensystem -- ZNS}
% \TagNo{Unterteilung nach Lage und Aufbau}
 \index{Zentralnervensystem}  
 \index{ZNS} 
\index{Nervensystem!Zentral-}
\label{topic:zns}

\Text
{\TxBeschreibung}
{Das \EE{Gehirn} und das \EE{Rückenmark} bilden
morphologisch\TODO{GABS}{2011-03-07}{morphologie erklären}{}
und funktionell eine Einheit. Gemeinsam verwerten und
verarbeiten sie die Signale aus der Peripherie und wählen
diejenigen aus, die entweder unterbewusst verarbeitet oder
bewusst wahrgenommen werden sollen. Das Zentralnervensystem
erzeugt schließlich \EE{Signale}, die -- oft als Antwort auf
einen Reiz -- in die Peripherie abgegeben werden.

Zu den Aufgaben des ZNS gehören die Aufnahme  und Verarbeitung  von äußeren und
    internen  Reizen
    (Schmerz, Fühlen, Sehen, Hören, {\dots}), das
\EE{Bewusstsein} und
    \enquote*{Bewusstmachung}, sowie die  Abgabe  von zweckentsprechenden Impulsen
    (Steuerung, Reaktion, {\dots}). Ebenso fällt das Gedächtnis 
    in den Verantwortungsbereich des ZNS.
}
{\begin{itemize}
  \item 
    \item Aufnahme  und Verarbeitung  von äußeren und
    inneren  Reizen
    \item Bewusstsein 
    \item Abgabe  von Impulsen
    \item Gedächtnis
\end{itemize}}
{}
{}
{}



\Layout{20}{Unterteilung nach Lage und Aufbau}
{%
Das Zentralnervensystem (ZNS) besteht aus dem \EE{Gehirn}
und dem \EE{Rückenmark}. 
Das Gehirn besteht aus dem
\E{Großhirn}, \Z{dem Zwischenhirn,} dem Kleinhirn und dem
\E{Hirnstamm}. Das Großhirn ist in zwei Großhirnhälften
geteilt\Z{, welche über den Balken miteinander verbunden
sind}.

\bigskip

\bigskip

\bigskip

\bigskip
}
{
\begin{itemize}
    \item Gehirn
    \begin{itemize}
        \item Großhirn
        \item \Z{Zwischenhirn}
        \item {Kleinhirn}
        \item Hirnstamm
    \end{itemize}
    \item Rückenmark
\end{itemize}
}
{./images/anatomie/Brain_human_sagittal_section.png}
{Das Hirn im seitlichen Querschnitt.}
{Patrick J. Lynch, medical illustrator,
\url{http://patricklynch.net}; C. Carl Jaffe, MD,
cardiologist. Lizenz: CC-BY 2.5}



%
\TextSlide{Großhirn}
{%
\index{Großhirn}%
Das Großhirn ist verantwortlich für Hören, Sehen, Sprache
sowie Riechen.
Außerdem sind dort Zentren für Gefühlsregungen und
Gedächnis. Auch bewusste Bewegungen und die Sensorik werden
im Großhirn verarbeitet.

\Z{Es besteht aus 2 Hälften, sie stehen über
        den Balken miteinander in Verbindung. An der
        Oberfläche sind Furchen und Windungen zu erkennen.
        Das Großhirn ist in Lappen gegliedert.}
}
{
\begin{itemize}
    \item Bewusstes Denken
    \item Sitz des Bewusstseins
    \item Gefühle
    \item Gedächtnis
\end{itemize}
}
{}{}{}

\TextSlide{\Z{Zwischenhirn}}
{%
\index{Zwischenhirn}%
\Z{Im  Zwischenhirn laufen alle Signale aus dem Körper
zusammen, hier befinden sich wichtige
\E{Regulationseinheiten}\ZFootnote{Der wichtigste Kern, der
Thalamus, ist das wichtigste Integrations-, Koordinations-
und Modulationszentrum sowohl für Sensorik als auch für
Motorik. Über den Hypothalamus, ein weiterer Kern des
Zwischenhirns, werden alle vegetativen Funkionen des Körpers
geregelt.}. Es ist eine \E{Schaltinstanz} und ein \E{Filter}
zwischen Rückenmark und Großhirn, auch die \E{Regulation der
vegetativen Funktionen} findet hier z.\,T. statt.}

\bigskip
}
{
\Z{
\begin{itemize}
\item Schaltinstanz und Filter Rückenmark - Großhirn
\item Regulation der vegetativen Funktionen
\end{itemize}}
}{}{}{}

\SlideCompensation{1}

\TextSlide{Kleinhirn}
{%
\index{Kleinhirn}%
Das Kleinhirn dient der Koordination- und Regulation der
Motorik. Zusammen mit dem Labyrinth des Innenohres sorgt es
für das Körpergleichgewicht. Die Bewerkstelligung von
komplexen Bewegungen hat auch hier ihren Ursprung.
\Z{Das Kleinhirn hat zwei Hälften, sowie feine Windungen und
Furchen auf der
Oberfläche.}
}
{
\begin{itemize}
    \item Bewegungskoordination ({\quotedblbase}Makros{\textquotedblleft})
    \item Gleichgewicht
\end{itemize}
}{}{}{}

\TextSlide{Hirnstamm und verlängertes Rückenmark}
{%
\index{Hirnstamm}%
\index{verlängertes Rückenmark}%
Der Hirnstamm ist das Regulationszentrum aller
lebenswichtigen Funktionen (Atmung, Blutdruck, Schlaf usw.).
Außerdem ist es die Ursprungsstelle (bzw. in manchen Fällen
die Ankunftstelle) der Hirnnerven.
}
{
\begin{itemize}
    \item Atem-, Kreislaufzentrum
    \item Hirnnerven
\end{itemize}
}{}{}{}

\SlideCompensation{1}

\TextSlide{Rückenmark}
{%
\index{Rückenmark}%
Das Rückenmark bildet die \EE{Verbindung zwischen Gehirn und
Peripherie}.
Funktionell steht es unter der Kontrolle übergeordneter
Zentren. Es können über einen Eigenapparat eingehende
Signale selbstständig verarbeitet werden (hierunter versteht
man Eigen- und Fremdreflexe). Diese Eigenschaft hat Einfluß
auf das Verhalten wenn übergeordnete Hirnzentren ausfallen,
z.\,B. nach einem Schlaganfall oder einer
Querschnittsverletzung: Nach einiger Zeit (Tage bis Wochen)
übernimmt dann das Rückenmark die Steuerung der
Muskelspannung (Entwicklung einer \E{Spastik}); eine durch
den Menschen bewusst gesteuerte Bewegung ist allerdings
dadurch nicht möglich.

\Z{Das Rückenmark lässt sich durch abgehende Spinalnerven in
\E{Segmente} (insg.
31) einteilen. Es gibt 8 Halssegmente (es gibt 8
Halssegmente, aber nur 7 Halswirbel!), 12 Thorakalsegmente,
5 Lumbalsegmente und 5 Sakralsegmente. Das Rückenmark selbst
endet auf Höhe des 1. bzw. 2. Lumbalwirbels, da es kürzer
als die Wirbelsäule ist; es ziehen jedoch die Nervenstränge
entsprechend der zugehörigen Wirbelkörper weiter nach unten
und verlassen erst dort den Wirbelkanal. Diese Konstruktion
erinnert an eine Quaste
 bzw. einen Pferdeschweif, daher leitet sich auch der
 Fachausdruck \I[Cauda equina|Z]{Cauda equina} (\Lang{lat.}
 Pferdeschwanz) ab.}
 }{
\begin{itemize}
    \item Weiterleitung von Impulsen (Leitungsbahnen)
    \item Vom Rückenmark gehen die peripheren Nerven aus
    \item (Muskeleigen-)Reflexe
    \item \Z{31 Segmente, Ende auf Höhe L1-L2}
\end{itemize}
}{}{}{}

\clearpage
\Layout{20}{Exkurs: Der Kopf besteht aus Schichten}
{%\index{Epiduralraum}
%\index{Dura mater}
\index{Harte Hirnhaut}%
%\index{Subduralraum}
%\index{Arachnoidea}
%\index{Subarachnoidalraum}
%
Unter dem Haupthaar und der damit verbundenen Haut folgt die
Kopfschwarte. Die Besonderheit der Kopfschwarte besteht
darin, dass Blutgefäße nicht wie sonst im Körper bei
Verletzung kolabieren. Dadurch kommt es \E{auch bei relativ
geringen Kopfverletzungen} zu vergleichsweise \E{starken
Blutungen}.

Hierauf folgt die Schädelkalotte. \Z{Diese ist innen fest
mit der harten Hirnhaut (Dura mater) verbunden - sie ersetzt
hier das Periost des Knochens. In die Dura sind die großen
venösen Abflussgefäße des Gehirns eingebaut.

Auf die harte Hirnhaut folgt die Spinngewebshaut
(Arachnoidea mater). In sie sind ebenfalls Gefäße
eingebettet.

Darunter folgt die weiche Hirnhaut (Pia mater). Sie liegt
den Furchen und Windungen des Gehirns untrennbar an.}
}
{
\begin{itemize}
    \item Kopfschwarte
    \item Knochen
%    \item Epiduralraum
    \item Harte Hirnhaut \Z{(Dura mater)}
%    \item Subduralraum
    \item Arachnoidea mater (Spinngewebshaut)
%    \item Subarachnoidalraum
%    \item eingebettete Gefäße
    \item Weiche Hirnhaut \Z{(Pia mater)}
\end{itemize}
}
{./images/anatomie/Gray1196-edited.png}
{\AuthNotesPicLegalOk{}{PD-CE}Querschnitt durch
die Schädeldecke. }
{Gray's Anatomy, Copyright expired.}



\noindent Das Hirn und das Rückenmark schwimmen in einer
Flüssigkeit, dem \T{Liquor}. Die
Schädelknochen und der Liquor bieten eine Schutzfunktion. 
Wenn Liquor austritt (z.\,B. durch Nase, Ohr, ...)  ist das
ein Zeichen für eine schwerwiegende Verletzung!

\Fazit{Tritt Liquor in die Umwelt aus, ist das ein Zeichen
für eine schwerwiegende Verletzung.} 

\AuthNotes{KOCH: doppelt? GABS: ja. Schlüsselsätze sollen
doppelt (normales Format + Fazit-Box) vorkommen}

% \paragraph{\Z{Blut-Hirn-Schranke}} 
% 
% \Z{Die Blutgefäßwand ist im Hirn besonders dicht. Sie
% filtert daher besonders gut und nur für bestimmte Moleküle (\ce{O2},
% \ce{CO2}, Wasser, \ldots) leicht passierbar. Sie bietet dem Hirn
% Schutz vor schädlichen und giftigen Substanzen. 
% 
% Die Blut-Hirn-Schranke bietet dem Hirn besonderen
% Schutz vor schädlichen Substanzen im Blut.} 

\subsection{Das Periphere Nervensystem -- PNS}
\index{Nervensystem!Peripheres}
\index{PNS}
% \TagNo{Unterteilung nach Lage und Aufbau}

\Definition{Unter dem Begriff \T{Peripheres Nervensystem} fasst man
die Verbindungsnerven zwischen dem ZNS und den Endorganen zusammen.}

\Layout{20}{Unterteilung nach Lage und Aufbau}
{%
Das periphere Nervensystem dient der Kommunikation zwischen
zentralem Nervensystem und den peripheren Organen. Es werden
dabei Informtaionen in \EE{beide Richtungen} übertragen,
d.\,h. vom ZNS zu den Organen (\Speech{\protect\enquote{Hirn
an Hand: Leg' dich auf die Herdplatte}}) und von den Organen
zum ZNS (\Speech{\protect\enquote{Haut von Hand an Hirn:
Hier ist es heiß!}}).

Es werden Hirnnerven und Spinalnerven unterschieden.
Hirnnerven (insgesamt 12 Paare) sind Nerven, die ihren
Ursprung im Hirnstamm haben. Sie sind grundsätzlich für die
Versorgung vom Kopf zuständig.

Die \EE{meisten Nerven entspringen segmentweise dem
Rückenmark} und werden Spinalnerven genannt. Sie verbinden
das Rückenmark mit den jeweiligen Organen. \Z{Für die
Extremitäten bilden sie jeweils Nervengeflechte, sogenannte
Plexus (Plexus cervicalis, Plexus brachialis, Plexus
lumbalis und Plexus sacralis).}

Die Fortleitung innerhalb eines Nerven funktioniert über
\E{elektrische Impulse}. \Z{Zwischen Nerven und anderen
Nerven bzw. den Zielorganen dienen chemische Signale der
Weiterleitung (in sog. Synapsen).}
}
{
\begin{itemize}
    \item Kommunikation ZNS ${\longleftrightarrow}$
    periphere  Organe
    \item Hirnnerven (12)
    \item Spinalnerven (gehen vom Rückenmark segmental ab)
    \item elektrische Impulse, chemische Signale (Synapsen)
\end{itemize}
}
{./images//hirtler-aass/WS_Aufbau.pdf}
{\AuthNotesPicLegalOk{}{}Rückenmark
mit abgehenden peripheren Nerven. Beachte dass die Nerven
geordnet nach Rückenmarkssegmenten abgehen.}
{Hirtler.}



\TextSlide{Arten von Nerven}
{%
\index{Nerven!motorische}
\index{Nerven!sensorische}
\index{Nerven!vegetative} 
\T{Motorische Nerven} sorgen dafür, dass ein
\protect\enquote{Bewegungswunsch} in den entsprechenden
Muskeln ankommt. Auf das Nervensignal hin kommt es zur
Kontraktion des Muskels.

\T{Sensorische Nerven} liefern dem Gehirn
\EE{Sinneseindrücke}. Diese können von den diversen
Sinnesorganen kommen. So gehören Tastempfindungen,
Temperaturempfindungen, Gerüche, Geschmacksempfindungen usw.
dazu. Aber auch weniger offensichtliche Dinge, wie z.\,B.
die Gelenksstellung gehören dazu.

\T{Vegetative Nerven} sind sozusagen die Drahtzieher im
Hintergrund. Sie sorgen für den reibungslosen \EE{Ablauf
unbewusster Funktionen}, wie z.\,B. die Herzfunktion, die
Verdauung und die Nierenfunktion. Die vegetativen Nerven
können je nach ihrer Funktion -- global stimulierend oder
hemmend --  in zwei Systeme unterteilt werden: In den
\EE{Sympathikus} und \EE{Parasympathikus}.
}
{
\begin{itemize}
    \item motorisch  
    \item sensorisch 
    \item vegetativ 
    \begin{itemize}
        \item sympathisch
        \item parasympathisch
    \end{itemize}
\end{itemize}
}{}{}{}

\subsection{Das vegetative (autonome)  Nervensystem}
\index{Nervensystem!Vegetatives} 
\index{Nervensystem!autonomes}
\index{Vegetatives Nervensystem}
\label{topic:vegetatives-nervensystem}
\label{topic:nervensystem-vegetatives}


 
\paragraph{Sympathikus  und Parasympathikus  sind Gegenspieler}~
\index{Sympathikus}  
\index{Parasympathikus}
\label{topic:sympathikus}
\label{topic:parasymapthikus} 
%  \TagNo{Unterteilung nach Funktion}

\vspace{\baselineskip}

\noindent\Parallel{
\T{Sympathikus}

Aktiviert den Aufmerksamkeits- und Kampfmodus.

\emph{\protect\enquote{Fight and flight}} -- \emph{\protect\enquote{Kampf und Flucht}}
}{%
\T{Parasympathikus} 

Aktiviert den Ruhemodus.

\emph{\protect\enquote{Rest and digest}} - \emph{\protect\enquote{Rast und Verdauung}}
}

z.\,B.: RR + HF

\noindent%
\Parallel{
Sympathikus:  HF ${\uparrow}$, RR ${\uparrow}$
}{
Parasympathikus: HF ${\downarrow}$, RR ${\downarrow}$
}




\begin{figure}[H]
    \begin{center}
    \includegraphics[width=\linewidth]{./images/hirtler-aass/Parasymp-Symp-Cartoon.pdf}
    \Captionof{figure}{\protect\enquote{Fight or flight, rest and digest}
    \SourcePicture{Hirtler.}}
    \end{center}
\end{figure}

\opt{STATUS-EXPERIMENTAL}{
\begin{figure}[H]
    \begin{center}
   
%     \includegraphics[width=\linewidth]{./images/noauth/AASdev20090228-img37.png} 
\includegraphics[width=\linewidth]{./icons/under-construction.png}

\TODO{GABS}{2010-04-11}{Übersicht vegetatives
Nervensystem: Bild fehlt.}{}

    \Captionof{figure}{\AuthNotesPicLegalNotOk{img37}{}
    Nur zur Illustration: Das vegetative Nervensystem hat
    Einfluss auf viele Organe. Dieses Bild
    \underline{nicht lernen}. Es dient nur der
    Verdeutlichung wie wichtig und einflussreich das vegetative
    Nervensystem im Körper ist. \SourcePicture{Quelle nicht eruierbar.}} 
    \end{center}
\end{figure}
}
 
 
 
\clearpage % TEMP temp clearpage
\section{Das Auge}
\index{Auge}
\label{topic:auge} 

\TextSlide{\TxBeschreibung}
{%
Das Auge als Sehorgan ist ein spezialisiertes Sinnesorgan,
dass sensibel auf Lichtreize reagiert. Es besteht aus:

\begin{itemize}
  \item Augapfel
  \item Hilfseinrichtungen: Augenmuskeln, Augenlider und
  Tränendrüsen
\end{itemize} 

Das gesamte Sehorgan befindet sich in der \I{Augenhöhle}. Die
lichtempfindlichen Zellen liegen auf der Netzhaut des Auges
und werden nach einer Passage durch die lichtbrechenden
Strukturen des Augapfels erreicht.
Der \I{Augapfel} hat eine annähernde Kugelform. Von außen nach
innen hat er verschiedene Schichten:

\begin{itemize}
  \item \T{Lederhaut} \Z{(Sklera)}: Die Lederhaut ist eine
  derbe Bindegewebsschicht. Sie ist zusammen mit dem
  Glaskörper verantwortlich für die Formgebung des Auges.
  \item \Z{Aderhaut (Choroidea): In der Aderhaut verlaufen
  alle Blutgefäße des Auges.}
  \item \T{Netzhaut} (Retina): In der Netzhaut befinden sich
  die Sinneszellen des Auges. Stäbchen für das Schwarz-Weiß-
  bzw. Dämmerungs-Sehen und Zapfen für das
  Farb-Sehen.
\end{itemize}
An der Vorderseite ist diese Schichten-Aufteilung anders:
\begin{itemize}
  \item \T{Hornhaut} \Z{(Cornea)}: Die Hornhaut ist eine
  durchsichtige Schicht von Bindegewebe. 
  \item \T{Regenbogenhaut} (Iris): Die Iris ist der
  verschiedenfärbige (blau, braun, grün) Bereich rund um die
  Pupille. Sie ist eine dünne Haut mit einer eigenen
  Muskulatur -- vegetativ gesteuert -- und sorgt dafür, dass
  bei Bedarf die Pupille vergrößert bzw. verkleinert wird.
\end{itemize}

Wenn man einem Lichtstrahl folgt, der in das Auge einfällt,
muss dieser zuerst die Hornhaut durchdringen, darauf folgt
die Vorderkammer (der Bereich zwischen Hornhaut und Iris).
Nun durchdringt er die \T{Linse}, sie bündelt das Licht und
sorgt dafür, dass wir scharf sehen können.
\Z{Die Linse hat einen eigenen Aufhängungsapparat, der sie
schmäler (für die Fernsicht) und breiter (für die Nahsicht)
machen kann. Nach der Linse folgt der Glaskörper. Er hat
eine gallertartige Konsistenz und verflüssigt sich mit dem
Alter. Zusammen mit der Lederhaut bestimmt er die pralle
Form des Auges.}

Nun landet der Lichtstrahl auf der \T{Netzhaut}. Die
Informationen der Sinneszellen (Stäbchen und Zapfen) werden
im Sehnerv gesammelt und an das Sehzentrum im Gehirn
weitergeleitet.

Wichtig für die Funktion des Auges sind weiters die
\I{Augenmuskeln}: Sie drehen das Auge in jede Richtung und
liegen ebenfalls in der Augenhöhle. Für die Befeuchtung des
Auges gibt es die Tränendrüse, die einen Flüssigkeitsfilm
auf dem Auge produziert und der über den Tränen-Nasengang
abfließen kann, und die \T{Bindehaut} (Konjunktiva), die bis
auf die Hornhaut die gesamte Vorderseite des Auges und die
Innenseite der Lider überzieht. Um das Auge beim Schlafen
vor Austrocknung zu schützen und den Tränenfilm gleichzeitig
verteilen zu können, verfügen wir über Augenlider, die
willkürlich und unwillkürlich geschlossen werden können.
}{%
\begin{itemize}
  \item Schichten des Auges:
  \begin{itemize}
    \item Lederhaut \Z{(Sclera)}
    \item \Z{Aderhaut (Choroidea)}
    \item Netzhaut (Retina)
  \end{itemize}
  \item Lichtstrahl von vorne nach hinten:
  \begin{itemize}
    \item Hornhaut \Z{(Cornea)}
    \item Vorderkammer
    \item Pupille, begrenzt von Iris
    \item Linse → Scharfstellen
    \item Glaskörper
    \item Netzhaut
    \item Sehnerv → Sehzentrum im Gehirn
  \end{itemize}
  \item Bewegung des Auges durch Augenmuskeln
  \item Schutzapparat des Auges:
  \begin{itemize}
    \item Tränendrüse und Bindehaut (Konjunktiva)
    \item Lider
  \end{itemize}
\end{itemize}
}{}{}{}



% \clearpage


\section{Das Ohr}
\index{Ohr}
\label{topic:ohr} 

\TextSlide{\TxBeschreibung}
{% 
Das Ohr ist das Organ für die akustische
Sinneswahrnehmung des Körpers. Zum Ohr gehört weiters das
Gleichgewichtsorgan.
Man kann das Ohr in drei verschiedene Untergruppen
unterteilen: In das \EE{Außenohr}, \EE{Mittelohr} und das 
\EE{Innenohr}.

Zum Außenohr gehören Ohrmuschel und äußerer Gehörgang. Die
Trennlinie ist hier das \T{Trommelfell}, das das Außenohr
hermetisch vom Mittelohr abschließt. Im Mittelohr befinden
sich die \EE{Gehörknöchelchen}. Es hat zum Druckausgleich
eine Verbindung zur Nase (Tube). So funktioniert es auch,
dass man beispielsweise bei einem Schnupfen durch Verstopfung des
Ohr-Nasen-Ganges und durch aufsteigende Keime eine
Mittelohrentzündung bekommen kann. Mit dem Übergang der
Gehörknöchelchen zur Schnecke passiert man die Grenze zum
Innenohr mit der Hörschnecke \Z{(Cochlea)} und dem
Gleichgewichtsorgan. \Z{Das Innenohr befindet sich im
Felsenbein.}
Wenn man einer Schallwelle folgt, die von der Ohrmuschel
eingefangen wurde, durchläuft man folgende Strukturen der
Reihe nach:
\begin{itemize}
  \item \E{Äußerer Gehörgang}: dieser geht nach vorne oben
  -- daher die Hörer des Stetoskopes immer in diese Richtung
  schauend einsetzen, sonst hört man
  nichts!
  \item \E{Trommelfell}: Es dient der Schallweiterleitung
  und zur Abgrenzung des Mittelohres und Innenohres gegen die Außenwelt.
  \item Nun folgen die Gehörknöchelchen:
  \begin{itemize}
    \item \I{Hammer} \Z{(Malleus)}: ist fest mit dem
    Trommelfell verwachsen, ist teilweise durch das
    Trommelfell hindurch sichtbar, verbunden mit dem
    \item \I{Amboss} \Z{(Incus)}: dieser gibt die erhaltenen
    Schwingungen weiter an den
    \item \I{Steigbügel} \Z{(Stapes)}: ist mit der Schnecke
    verbunden
  \end{itemize}
  \item \E{Schnecke} \Z{(Cochlea)}: ist mit einer eigenen
  Flüssigkeit gefüllt. In der Schnecke werden die
  empfangenen Schallwellen nach ihrer Frequenz aufgelöst und
  durch die Sinneszellen über den Gehörnerven an das
  Hörzentrum im Gehirn weitergeleitet.
\end{itemize}
Das \EE{Gleichgewichtsorgan}, als Teil des Innenohrs in
direkter Nähe der Schnecke im Felsenbein liegend, ist dafür
zuständig, dass man ein Gefühl der eigenen Lage im Raum hat
und Beschleunigungen feststellen kann. Hierfür hat es drei
Bogengänge, die nach den drei Achsen im Raum ausgerichtet
sind. Jeweils die komplementären Bogengänge links und rechts
sorgen für eine vollständige Informationserfassung. Wenn das
Gleichgewichtsorgan einer Seite eine Störung hat, resultiert
ein kaum aushaltbarer Schwindel.
}
{%
\begin{itemize}
  \item Gliederung des Ohres:
  \begin{itemize}
    \item Außenohr: Ohrmuschel und äußerer Gehörgang
    \item Mittelohr: Trommelfell, Hammer, Amboß, Steigbügel
    \item Innenohr: Schnecke, Gleichgewichtsorgan
  \end{itemize}
\end{itemize}
}{}{}{}

\begin{PictureSeries}{}{Das Ohr}
\PictureSeriesRow{2}
{./images/hirtler-aass/Ohr.pdf}
{Das Ohr im Querschnitt. Man sieht
den äußeren Gehörgang, Das
Trommelfell, die Gehörknöchelchen,
die Schnecke und die Bogengänge. \SourcePicture{Hirtler}}
{./images/hirtler-aass/Mittelohr.pdf} {Das Mittelohr: Im Detail sind hier
das Trommelfell, sowie die Gehörknöchelchen Hammer, Amboss
und Steigbügel zu sehen.
\SourcePicture{Hirtler}}
{}
{}
{}
{}
\end{PictureSeries}



% \clearpage % TEMP temp clearpage
\section{Die Haut}
\index{Haut}
\label{topic:haut}

\HeadingSub{\Lang{Term.} Derma}


\noindent Die Haut ist die Grenzfläche des Körpers zur
Außenwelt und gleichzeitig das größte Sinnesorgan des
menschlichen Körpers.
\Z{Sie macht ca. \nicefrac{1}{6} des Körpergewichts aus,
ihre Oberfläche beträgt beim Erwachsenen ca.
1,6\Meter*{\texttwosuperior}.}

\TextSlide{Aufgaben}
{%
Die Haut ist \E{die} Barriere gegen Einflüsse aus der
Umwelt. Sie verhindert Eindringen von Fremdkörpern und
Keimen. Sie sorgt für einen ausgeglichenen Wasser- und
Elektrolythaushalt, indem sie den Körper vor Austrocknung
schützt.

Durch Absonderung von Schweiß bzw. der Regelung der
Hautdurchblutung sowie durch ihre Isolationsfähigkeit sorgt
die Haut für eine gleichbleibende Körperkerntemperatur bei
wechselnden Witterungsverhältnissen.
Als größtes Sinnesorgan des Körpers dient die Haut der
Tastempfindung sowie der Temperaturempfindung.


\SlideCompensationNoParskip{1}

}
{
\begin{itemize}
  \item Schutz vor äußeren Einflüssen
  \item Schutz vor Wasserverlust/Austrocknung
  \item Temperaturregulation
  \begin{compactitem}
    \item Isolation
    \item Wärmeabgabe
  \end{compactitem}
  \item Sekretionsorgan (Schweiß, Talg)
  \item Sinnesorgan (Tast-, Temperatursinn)
\end{itemize}
}{}{}{}

\TextSlide{Die Haut ist in Schichten aufgebaut}
{\index{Oberhaut}\index{Lederhaut}\index{Unterhaut}%
Die oberflächlichste Schicht der Haut ist die \T{Oberhaut}
(Epidermis). Sie besteht aus einer mehr oder weniger dicken
\E{Hornhaut} (am dicksten an Handflächen und Fußsohlen, am
dünnsten an den Augenlidern) und besitzt keine eigene
Gefäßschicht.

Darauf folgt die \T{Lederhaut} (Dermis). In ihr sind alle
Hautanhangsgebilde (Schweißdrüsen, Talgdrüsen, Haare und
Nägel) befestigt. Weiters befinden sich in dieser
Hautschicht Tastkörperchen sowie freie Nervenendigungen
(\EE{Schmerzempfindungen}). Außerdem finden sich hier zwei
verschiedene Gefäßsysteme (eines oberflächlicher, eines
tiefer). Dieses dient wie bereits erwähnt der
Temperaturregulation.
Unter der Dermis befindet sich die \T{Unterhaut}
(\T{Subcutis}). Sie besteht hauptsächlich aus Fett. Dieses
dient der Wärmeisolierung.

\noindent\begin{minipage}{\linewidth}
\begin{center}
\includegraphics[width=0.48\linewidth]{./images/hirtler-aass/Haut-edited.png}
\end{center}

\CaptionofDiff{figure}{Die Haut im Querschnitt: \TabSub{A:} Oberhaut \Z{Epidermis},
\TabSub{B:} Lederhaut \Z{Dermis},
\TabSub{C:} Unterhaut (Subkutis),
\TabSub{d:} \Z{Meißner'sches Tastkörperchen},
\TabSub{e:} Talgdrüse,
\TabSub{f:} Haar mit Haarfollikel,
\TabSub{g:} Schweißdrüse,
\TabSub{h:} Venole,
\TabSub{i:} Freie Nervenenden (Schmerzempfindung),
\TabSub{j:} Arteriole,
\TabSub{k:} \Z{Pacini-Körperchen}.

\SourcePicture{Hirtler.}}{Die Haut im Querschnitt \SourcePicture{Hirtler.}}
\end{minipage}


}
{
\begin{itemize}
  \item \T{Oberhaut} \Z{(Epidermis)} \index{Oberhaut}
  \begin{itemize}
    \item Hornschicht, geschichtetes Plattenepithel
  \end{itemize}
  \item \T{Lederhaut} \Z{(Corium, Dermis)}
  \index{Lederhaut}
  \begin{itemize}
    \item Hautanhangsgebilde
    \begin{itemize}
      \item Schweißdrüsen
      \item Talgdrüsen
      \item Haarfollikel, Nagelmatrix
    \end{itemize}
    \item Tastkörperchen, freie Nervenendigungen (Schmerz)
    \item \Z{Kollagenfasern (Stütznetzwerk)}
    \item Arteriolen, Venolen (Versorgung der Haut, Temperaturregulation)
  \end{itemize}
  \item \T{Unterhaut} (\T{Subcutis}) \index{Unterhaut}
  \index{Subcutis}
  \begin{itemize}
    \item Depotfett, Wärmeisolierung
  \end{itemize}
\end{itemize}\index{Haut!Querschnitt}
\AuthNotes{Haut: Bild Querschnitt: muss vereinfacht werden,
lat. Ausdrücke weg bzw. grau. GABS, 2010-01-23: done. Bild von Hirtler
übernommen.}
}
{./images/hirtler-aass/Haut-edited.png}
{\AuthNotesPicLegalOk{}{Haut
Querschnitt}Die Haut im Querschnitt: \TabSub{A:} Oberhaut \Z{Epidermis},
\TabSub{B:} Lederhaut \Z{Dermis},
\TabSub{C:} Unterhaut (Subkutis),
\TabSub{d:} \Z{Meißner'sches Tastkörperchen},
\TabSub{e:} Talgdrüse,
\TabSub{f:} Haar mit Haarfollikel,
\TabSub{g:} Schweißdrüse,
\TabSub{h:} Venole,
\TabSub{i:} Freie Nervenenden (Schmerzempfindung),
\TabSub{j:} Arteriole,
\TabSub{k:} \Z{Pacini-Körperchen}.
}
{Hirtler.}




\TextSlide{Die Haut ist wichtig für die
\E{Thermoregulation}}
{%
Neben der Lunge dient vor allen Dingen  die Haut als
wichtigstes Organ zur \EE{Thermoregulation}: Die Blutgefäße 
der  Haut sind stark ineinander verwickelt und können sich
bei Bedarf stark erweitern. Die stark durchblutete  Haut
strahlt somit Wärme ab. Weiters wird durch die
Verdunstungskälte  des Schweißes  ebenfalls Wärme abgegeben.
}
{
\begin{itemize}
  \item Erweiterung - Kontraktion der Blutgefäße
  \item Verdunstungskälte des Schweißes
\end{itemize}
}{}{}{}

\TextSlide{Wundliegen}
{%
\index{Wundliegen}\index{Dekubitus}\label{topic:dekubitus}Wenn
der Mensch nicht fähig ist in ausreichendem Umfang
regelmäßig seine Lage zu ändern, kann es zum Wundliegen
kommen. Die Wunden nennt man dann \T{Dekubitus} oder
Druckgeschwüre.
\Z{Besonders gefährdet sind jene Stellen, an denen ein
Knochen knapp unter der Haut liegt (Druckstelle). Betroffen
sind vor allem alte oder pflegebedürftige bettlägrige
Menschen, bewusstlose Patienten bzw. sedierte
Intensivpatienten sowie Patienten mit Querschnittsläsionen.}
}
{
\begin{itemize}
  \item Unfähigkeit der Lageänderung
  \item besonders an prädestinierten Druckstellen
  \item Betroffene: Bettlägrige, Bewusstlose, Intensivpatienten,
  Querschnittsgelähmte
\end{itemize}
}{}{}{}

\clearpage
% \bigskip
\section{Die Atemwege}
 \index{Atemwege}
 \label{topic:atemwege-lunge}


\AuthNotes{KOCH: mag Nasenrachenraum und Mundrachenraum
streichen, evtl. wie Bild?}

\Layout{27}{}
{}
{}
{./images/anatomie/Respiratory_system_complete_de-edited.pdf}
{\AuthNotesPicLegalOk{}{PD}Die Atemwege, Übersicht.\label{fig:respiraktionstrakt-uebersicht}}
{Mariana Ruiz Villarreal, Public domain}


%\begin{multicols}{2}
Die Atemwege gliedern sich in die \EE{oberen} und die
\EE{unteren} Atemwege.

\TextSlide{Obere Atemwege}
{%
Der Weg der Einatemluft beginnt in der Nase bzw. im Mund. Die Nase ist
mit Schleimhaut ausgekleidet, knöcherne Vorsprünge sorgen für
eine Vorwärmung der Luft. In der Schleimhaut be\-finden sich die
Sinneszellen für den Geruchssinn. 

\Z{In den Schädelknochen befinden sich luft\-ge\-füllte
Hohlräume, die so\-genannten \T{Neben\-höhlen} (Stirn-,
Kiefer\-höhlen, Siebbeinzellen). Diese stehen in Verbindung
mit dem Nasenrachenraum.
Entzündet sich die Schleimhaut spricht man von der
\E{Neben\-höhlen\-ent\-zündung}, welche oft langwierig und
schmerzhaft sein kann.}


Die Luft gelangt nun über die Nasenhöhle in den \T{Rachen}
\Z{(\I{Pharynx})}. Den Rachen kann man in drei Etagen
einteilen:
\begin{itemize}
  \item \I{Nasenrachen} \Z{(\I{Nasopharynx},
  \I{Epipharynx})}
  \item \I{Mundrachen} \Z{(\I{Oropharynx}, \I{Mesopharynx})}
  \item \I{Schlundrachen} \Z{(\I{Laryngopharynx},
  \I{Hypopharynx})}
\end{itemize}
Diese drei Abschnitte gehen fließend ineinander über. Im
Nasenrachen mündet die \I{Ohrtrompete}\ZFootnote{(Tuba
auditiva, Eustachische Röhre)}, die für den Druckausgleich
des Ohres sorgt. Bei Verstopfung dieses Kanals kommt es zu
einer Abflussbehinderung und somit zu einem Rückstau des
Sekrets in das Mittelohr. Hierdurch entsteht eine
Entzündung.

Im Mundrachen \EE{kreuzt der Weg der
Nahrung mit dem Weg der Atemluft}: Die Nahrung gelangt vom
Mundraum in die hinten gelegene Speiseröhre, die Luft
vom Nasenrachenraum in die vorne gelegene Luftröhre.
Nasenrachen und Mundrachen können während des
Schluckvorgangs durch den weichen Gaumen voneinander
getrennt werden. Die Grenze zwischen Mundrachen und
Schlundrachen bildet der Kehldeckel, ein Teil des Kehlkopfs.

\subparagraph{Kehlkopf und Kehldeckel}
Der \T{Kehlkopf} besteht aus mehreren Teilen, die aus
Knorpel aufgebaut sind:
\begin{itemize}
  \item \T{Kehldeckel} \Z{(\I{Epiglottis})}: 
  verschließt die Luftröhre während dem Schlucken
  \Z{
  \item \IX{Schildknorpel} \Z{(\I{Cartilago thyroidea})}:
  bildet Vorderseite des Kehlkopfes, ist beim Mann größer und
  hervorstehend
  \item \IX{Ringknorpel} \Z{(\I{Cartilago cricoidea})}:
  bildet die untere Grenze des Rachenraumes. Zwischen
  Schild- und Ringknorpel befindet sich die Membran, die bei
  der Koniotomie (ugs. \enquote*{\I{Luftröhrenschnitt}}) zur
  notfallsmäßigen Sicherung der Atmung bei Zuschwellen der
  oberen Atemwege durchtrennt wird.
  \item \IX{Stellknorpel} \Z{(\I{Cartilago arytenoidea})}:
  befinden sich innerhalb des Schildknorpel-Bereichs. Über
  diese Knorpel (2) bewegen und spannen sich die
  Stimmbänder.
  }
\end{itemize}

\subparagraph{Schluckakt und Schluckrefelex}
Der Kehlkopf spielt nicht nur eine Rolle bei der
Stimmbildung über die \T{Stimmbänder}, sondern sorgt auch dafür,
dass während des Schluckens kein Speisebrei und keine
Flüssigkeiten in die vor der Speiseröhre liegende Luftröhre
gelangen können. Hierzu \EE{verschließt der \I{Kehldeckel}
am Anfang des \I[Schluckakt]{Schluckaktes} den Eingang zur
Luftröhre}.
Dies geschieht reflexartig und ist lebenswichtig.
Dieser Reflex ist bei Bewusstlosigkeit und während einer
Narkose außer Gefecht gesetzt.
Dadurch kann es zu einem Übertritt von
Erbrochenen, Blut o.\,ä. in die Luftröhre
(\I{Aspiration}) und in Folge zu einer lebensgefährlichen
Lungenentzündung kommen.
Um dieses Risiko zu minimieren muss man für eine
geplante Operation nüchtern sein.
Ein Ausfall dieses Reflex kann u.\,U. auch ein Zeichen des
Hirntods sein.

\Fazit{Im Rachenraum kreuzt der Speise- und der Luftweg!}

\Fazit{Der Kehldeckel verschließt die Luftröhre während
des Schluckaktes!} 

}{%
\begin{itemize}
    \item \EE{Nasen}- und \EE{Mundhöhle} \index{Nasenhöhle} 
    \index{Mundhöhle} 
    \item \EE{Rachen} (gemeinsamer Weg von Nahrung und
    Atem!) \index{Rachen} 
    \item \EE{Kehlkopf} (\T{Larynx}) mit
    \EE{Kehldeckel} \index{Kehlkopf} \index{Larynx} 
    \index{Kehldeckel}
\end{itemize}
}{}{}{}

\TextSlide{Untere Atemwege}
{%
\label{topic:alveolen}
In der \T{Luftröhre} (\T{Trachea}) wird die vom Rachen und
Kehlkopf kommende Luft zur Lunge weitergeleitet. Die
Luftröhre teilt sich in einen \E{rechten} und einen
\E{linken} \T{Hauptbronchus}, welche zu den jeweiligen
Lungenflügeln ziehen. Jeder Hauptbronchus verzweigt sich
immer weiter bis schlussendlich sehr dünne \E{Bronchioli} in
die \T{Lungenbläschen} (\T{Alveolen}) münden. Hier findet
der Gasaustausch statt.

Die Abschnitte der Atemwege, in denen kein Gasaustausch
stattfindet und die nur dem Transport dienen, bezeichnet man
als \T{Totraum}.
}{%
\begin{itemize}
    \item \index{Luftröhre}\EE{Luftröhre}
    (\index{Trachea}\T{Trachea})
    \item \T{Hauptbronchus} (li./re.)
    \item \IX{Bronchioli}
    \item \EE{Lungenbläschen}
    ({Alveolen})
\end{itemize}
}{}{}{}



%\end{multicols}


\subsection{Die Lunge ist für einen Teil des
Gasaustausches zuständig}
\index{Gasaustausch!Lunge}  
\index{Lunge}
\index{Lunge!-nflügel}

\TextSlide{\TxBeschreibung}
{
Die Lunge hat \EE{zwei Flügel}, welche sich jeweils in \EE{3
(rechts)} bzw.
\EE{2 (links) Lappen} gliedern:

\begin{center}\TableBaseModification
\begin{tabularx}{\linewidth}[H]{XX}
\toprule\addlinespace
\multicolumn{2}{c}{\TabHeading{Die Lunge}}
\\\addlinespace\midrule\addlinespace
\multicolumn{1}{c}{\TabSub{Rechter Flügel}}
 &
\multicolumn{1}{c}{\TabSub{Linker Flügel}}
\\\addlinespace\cmidrule(lr){1-1}\cmidrule(lr){2-2}\addlinespace
\begin{itemize}
    \item \EE{3} Lappen
    \begin{itemize}
        \item Oberlappen \index{Oberlappen} 
        \item Mittellappen \index{Oberlappen} 
        \item Unterlappen \index{Unterlappen} 
    \end{itemize}
\end{itemize}
&
\begin{itemize}
    \item \EE{2} Lappen
    \begin{itemize}
        \item Oberlappen
        \item Unterlappen
    \end{itemize}
\end{itemize}
\\\addlinespace\bottomrule
\end{tabularx}
\end{center}

Die Oberfläche ist mit dem \T{Lungenfell} \index{Lungenfell}
\index{Pleura!lungenseitige}, der \E{lungenseitigen}
\T{Pleura}, überzogen. In einem Lungenflügel verzweigt sich
der Hauptbronchus in immer weitere kleinere Bronchien bis
schließlich die Bronchioli in die \E{Lungenbläschen}
(\E{Alveolen}) enden.
}
{
\begin{itemize}
  \item 2 Flügel mit rechts 3 Lappen und links 2
  \item Lungenfell (Pleura)
  \item Lungenbläschen (Alveolen)
\end{itemize}
}{}{}{}

\subsubsection{Die Alveolen sind die Endstücke
des Atemwegs und sind für den Gasaustausch zuständig}

\TextSlide{\TxBeschreibung}
{
\emph{Siehe Abb.
\FullPrettyRef{fig:respiraktionstrakt-uebersicht}, rechts
oben.} In den \T{Lungenbläschen} (\T{Alveolen}) findet der
\index{Gasaustausch!Lunge}\E{Gasaustausch} zwischen der
Einatemluft und dem Blut statt. Sauerstoff gelangt dabei aus
den Alveolen in das Blut und Kohlendioxid wird im Gegenzug
vom Blut an die Alveolen abgegeben.

Dabei besteht zwischen Luft und Blut kein direkter Kontakt,
die Gase müssen Hindernisse wie die Alveolenwand und die
Blutgefäßwand überwinden. \E{Vergrößert sich dieser Weg
aufgrund einer Erkrankung (z.\,B. beim
Lungenödem\VARIANT{VAR-STANDARD}{,
\FullPrettyRef{topic:lungenoedem})}, ist die Sauerstoffaufnahme vermindert!} 

\Fazit{Lunge: Sauerstoff (\ce{O2}):   Alveolen ${\rightarrow}$ Blut}

\Fazit{Lunge: Kohlendioxid (\ce{CO2}):  Blut ${\rightarrow}$
Alveolen } }
{
\begin{itemize}
  \item Gasaustausch: Kohlendioxid gegen Sauerstoff
  \item Hindernisse: Alveolenwand, Blutgefäßwand
\end{itemize}
}{}{}{}

\clearpage
\subsubsection{Die Pleura}\label{topic:pleura}

\TODO{GABS}{2011-03-07}{Pleura: zuviel Text? Redundant? kürzen?}{}

\Layout{20}{\TxBeschreibung{} und Funktion}
{%
\index{Rippenfell}%
\index{Lungenfell}%
\index{Pleura!thoraxseitige}%
\index{Pleura!lungenseitige}%
Die Pleura hat zwei \protect\enquote{Blätter}. Eine Seite 
liegt an der Rippenseite (\T{Rippenfell}) an und ist fest
mit dem Zwerchfell und der Brustwand verwachsen. Die zweite
Seite  überzieht die Lungenoberfläche (\T{Lungenfell}).
Zwischen beiden Pleurablättern befindet sich ein
\E{Flüssigkeitsfilm}. Dieser sorgt dafür, dass die beiden
Blätter seitlich gleiten können. (Es gibt keine
Verwachsungen der beiden Seiten und die Pleuraflüssigkeit
fungiert sogar ein bisschen als
\protect\enquote{Schmiermittel}.)

Der Zwischenraum ist \EE{dicht abgeschlossen}, es kommt von
außen keine Luft dazu. Dadurch haften die beiden Blätter
aneinander, bei Zug (durch Unterdruck während der Einatmung)
an einem der beiden Blätter lassen sich diese nicht trennen,
folglich wird das Zweite passiv mitbewegt.\footnote{Wie bei
zwei Glasplatten mit Wasser dazwischen.} Der Vorteil dieser
Konstruktion ist, dass zwar die Lunge jeder Bewegung des
Brustkorbs folgt, jedoch nicht fixiert ist, sondern
verschieblich bleibt.

\Z{Beim Gesunden gibt es keine Verwachsungen zwischen diesen
beiden Pleurablättern. Bei Entzündungen kann es jedoch dazu
kommen. Die Pleura ist außerdem die einzige Struktur der
Lunge die schmerzempfindlich ist.} }
{
\begin{itemize}
    \item \E{Thoraxseitige} Pleura (Rippenfell)
    \item \E{Lungenseitige} Pleura (Lungenfell)
    \item Zwischenraum luftdicht, mit Schmierflüssigkeit
    \item Brustkorb mit Rippenfell \protect\enquote{zieht an} → Lungenfell mit Lunge folgt
    Bewegung     
\end{itemize}
}
% {./images/noauth/AASdev20090228-img41.png}
{./images/hirtler-aass/Pleura.pdf} {Pleura, Schema.} {Hirtler.}


\EE{Frage}: \Speech{Was würde passieren wenn der
Zwischenraum zwischen den Pleurablättern aufgrund einer
Verletzung (z.\,B. ein Loch aufgrund einer Stichwunde im
Brustbereich) nicht mehr luftdicht ist?} Antwort: s.
Fußzeile (\footnote{Antwort: Die Lunge würde in sich
zusammenfallen, da sie ja nicht mehr aufgespannt wird:
\I[Pneumothorax!anatomische
Erklärung]{Pneumothorax}\VARIANT{VAR-STANDARD}{
(\FullPrettyRef{topic:pneumothorax})}.})


% \TextSlide{Welche Funktion hat aber nun die Pleura?}
% {%
% Da die eine Seite mit dem Brustkorb, die andere Seite mit
% der Lunge verwachsen ist und der Raum zwischen den beiden
% Seiten luftdicht abgeschlossen ist, haften die beiden Seiten
% aneinander, aber sind gegeneinander verschieblich. (Es gibt
% keine Verwachsungen der beiden Seiten und die
% Pleuraflüssigkeit fungiert sogar ein bisschen als
% \protect\enquote{Schmiermittel}.)
% 
% Wenn sich nun der Brustkorb oder das Zwerchfell bei der
% Atmung bewegt, bewegt sich natürlich das Rippenfell mit (es
% ist ja mit dem Brustkorb und dem Zwerchfell verwachsen). 
% }
% {
% \begin{itemize}
%   \item Leine Luft zwischen beiden Pleuren
%   \item Pleuraflüssigkeit = Schmiermittel
%   \item Unterdruck bei Atmung
%   \item bei Verletzung bzw. Lufteindringen in den Pleuraspalt Pneumothorax
% \end{itemize}
% }{}{}{}

\Fazit{Wenn der Brustkorb mit dem daran verwachsenen
Rippenfell \protect\enquote{anzieht}, folgt das Lungenfell
mit der Lunge der Bewegung \Sign{→} die Lunge wird wie eine
Ziehharmonika aufgezogen und Luft angesaugt.}

\subsection{Atemmechanik}
\label{topic:atemmechanik} 
\index{Atemmechanik} 
\index{Zwischenrippenmuskulatur}

Wenn man sich den Thorax annähernd als Zylinder vorstellt,
so kann eine Volumenvergrößerung (Einatmung) durch eine
Vergrößerung des Querschnitts (Heben der
Zwischenrippenmuskulatur)
oder Vergrößerung der Höhe (Anspannung des Zwerchfells) erfolgen.

\begin{equation}
V = A \times h
\end{equation}
\begin{equation}
\mbox{Volumen} = \mbox{Fläche} \times \mbox{Höhe}
\end{equation}

\AuthNotes{KOCH: entweder wird das genauer erklärt wie die
Funktion der Pleura oder es wird ganz einfach mündlich
erklärt.

GABS: Sollte schon schriftlich erklärt werden \ldots wir
wollen ja auch immer dass sie die \protect\enquote{Atemhilfsmuskulatur}
etc. unterstützen, dann sollten sie wenigstens auch
nachlesen können \protect\enquote{wie man richtig atmet}

KOCH: soll erklärt werden.}

\begin{PictureSeriesWide}{}{Atemmechanik}
\PictureSeriesRowWide{3}
{./images/hirtler-aass/Thorax-Diaphragma.pdf}
{Das \IX{Zwerchfell} ist der wichtigste Atemmuskel und
trennt den Brust- und Bauchraum. \SourcePicture{Hirtler}}
{./images/hirtler-aass/Thoraxbewegung-frontal.pdf} {Bei der
Einatmung senkt sich das Zwerchfell und der Brustkorb hebt
sich. \SourcePicture{Hirtler}}
{./images/hirtler-aass/Thoraxbewegungen-Rippe.pdf}
{Die Rippen folgen der Atembewegung.
\SourcePicture{Hirtler}}
{}
{}
\end{PictureSeriesWide}





% \begin{figure}[H]
% \begin{center}
% % \includegraphics[width=\linewidth]{./images/noauth/AASdev20090228-img42.png}
% \includegraphics[width=\linewidth]{./images/hirtler-aass/Atmung.pdf}
% 
% \Captionof{figure}{Atemmechanik:  Das
% mittlere Bild zeigt die Zwischenrippenmuskulatur.
% \SourcePicture{Hirtler.}}
% \TODO{GABS}{0.14.0}{Atemmechanik: Bilder trennen und
% getrennt beschreiben.}{}
% \end{center}
% \end{figure}


\TextSlide{\TxBeschreibung}
{
Für die Atemmechanik ist ein gutes Zusammenspiel von
Skelett, Muskeln, Gewebe und Pleura(spalt) wichtig.

\Fazit{Ein bis zum Hals Verschütteter erstickt.}

\subparagraph{Atemmuskulatur} Der wichtigste Atemmuskel ist
das \EE{Zwerchfell}\index{Zwerchfell}. In Ruhe erledigt es
fast die gesamte Atemarbeit. Bei verstärkter Atmung kommen die
Zwischenrippenmuskeln zum Einsatz \cite{LippertAnatomie7}.

\Fazit{Der wichtigste Atemmuskel ist das \EE{Zwerchfell}!}

\subparagraph{Atemhilfsmuskulatur}% 
\label{topic:atemhilfsmuskulatur}%
\index{Atemhilfsmuskulatur}%
Neben den Hauptatemmuskeln (Zwerchfell,
Zwischenrippenmuskulatur) gibt es noch eine Reihe weiterer
Atemhilfsmuskeln die die Atmung unterstützen. Diese
verbinden den Thorax und die Oberarme. Fixiert man die Arme,
z.\,B. durch Aufstützen, können die Muskeln die Hebung des
Brustkorbes unterstützen.
\VARIANT{VAR-STANDARD,VAR-ERNAEHRUNGSPAED}{Vgl.
dazu
}\VARIANT{VAR-STANDARD,VAR-ERNAEHUNGSPAED}{\emph{Asthma},
\FullPrettyRef{topic:asthma-bronchiale};
}\VARIANT{VAR-STANDARD}{\emph{COPD},
\FullPrettyRef{topic:copd}}\VARIANT{VAR-STANDARD,VAR-ERNAEHRUNGSPAED}{.}

\subparagraph{Pleura} Die Funktion der Pleura ist unter \ref{topic:pleura}
erklärt. \Z{Der Unterdruck in der Pleura entsteht dadurch,
dass sich einerseits die (elastische) Lunge zusammenziehen
möchte, während andererseits der Brustkorb einen Gegenzug
ausübt. Diese beiden Kräfte stehen bei entspannter
Atemmuskulatur im Gleichgewicht.}


% \Z{Thoraxwand und Lunge besitzen elastische Eigenschaften: Sie werden gedehnt
% und entwickeln elastische Rückstellkräfte. An der Lunge können zwei Mechanismen ein
% Zusammenfallen des Lungengewebes bewirken. Erstens sind im Lungenzwischengewebe
% reichlich elastische Fasern eingelagert, die die Eigenelastizität der Lunge
% bewirken. Durch die Koppelung von Lunge und
% Thoraxwand am Pleuraspalt entsteht ein elastisches Gesamtsystem. 
% 
% Bei entspannter
% Atemmuskulatur und offenem Kontakt mit der Umgebungsluft wird sich ein
% Gleichgewicht der elastischen Kräfte einstellen; der Brustkorb befindet sich in
% der Atemruhelage. Durch den Zug der Lunge, die sich verkleinern möchte, und den
% Gegenzug der Thoraxwand entsteht ein ständiger Unterdruck im Pleuraspalt.}


\subparagraph{Atemzyklus}
Ein Atemzyklus besteht aus der Einatemphase und der
Ausatemphase. Dabei kommt es zu einem Ungleichgewicht
zwischen der nach außen gerichteten Zugkraft des Thorax und
der nach innen gerichteten Zugkraft der Lunge. Dieses
Ungleichgewicht wird beim Einatmen zunächst von einer
zusätzlichen Zugkraft des Zwerchfells und der äußeren
Zwischenrippenmuskulatur hervor gerufen und anschließend
beim Ausatmen durch Entspannung der Muskeln wieder abgebaut.

Während des Einatmens wird durch aktive Bewegung des
Zwerchfells und der äußeren Zwischenrippenmuskulatur das
Volumen des Brustkorbs und somit auch das der Lunge
vergrößert. Die in der Lunge befindliche Luft hat nun
deutlich mehr Platz, weswegen ihr Druck jetzt kleiner ist,
als der Luftdruck in der Umgebung (\EE{Unterdruck}). Dieser
Druckunterschied wird durch einströmende Luft ausgeglichen
-- man hat eingeatmet.
Da für das Einatmen Muskelarbeit verrichtet werden muss,
nennt man das Einatmen auch \E{aktive Phase} der Atmung.

Das Ausatmen erfolgt passiv (\E{passive Phase} der Atmung).
Das Zwerchfell und die äußere Zwischenrippenmuskulatur
entspannen wodurch sich das Thorax- bzw.
Lungenvolumen auf das ursprüngliche Volumen
verkleinert.\footnote{Man denke an eine Feder, welche
nachdem sie durch eine Zugkraft gedehnt wurde, nach dem
Loslassen durch ihre eigene (elastische) Rückstellkraft
wieder von selbst in die Ausgangslage zurück kehrt.} Durch
die Volumenverminderung steigt der Druck in der Lunge und
die Luft strömt nach außen bis die Druckunterschiede
zwischen Lunge und Umgebung wieder ausgeglichen sind. Die
Luft wird beim Ausatmen sozusagen wieder hinaus gedrückt.

Die Anzahl der Atemzyklen pro Minute nennt man \I{Atemfrequenz}
(\I{AF}). Sie liegt beim Erwachsenen, der sich in Ruhe befindet,
bei etwa 12 Atemzyklen pro Minute.

% Volumenschwankungen des Brustraums bewirken
% gleichzeitig auch Volumenschwankungen im Inneren der Lunge. Eine Vergrößerung
% des Thoraxinnenraums wird durch Kontraktion der Atemmuskulatur bewirkt. Dieser
% Vorgang wird als aktive Phase der Atmung bezeichnet. Dabei kommt es vor allem
% zu einer Abflachung und einem Tiefertreten des Zwerchfells, aber auch zu einer
% Anhebung der Rippen durch die äußeren Zwischenrippenmuskeln. Als Folge entsteht
% in der Lunge im Vergleich zur Umgebungsluft ein Unterdruck, der durch
% Einströmen von Luft ausgeglichen wird. Thoraxwand und Lunge werden dabei
% gedehnt.
% 
% Am Ende einer ruhigen Einatmung
% erschlafft die Einatmungsmuskulatur. Die elastischen Rückstellkräfte führen das
% Thoraxvolumen wieder in die Ausgangslage zurück. Dabei kommt es zu einer
% Druckerhöhung in der Lunge, die eine Luftbewegung hin zur Umgebungsluft bewirkt.
% Diese erfolgt in aller Regel passiv, d. h. ohne aktive
% muskuläre Anstrengung. Dieser Vorgang erfolgt in Ruhe bei einem Erwachsenen
% etwa zwölfmal in der Minute und wird als Atemfrequenz bezeichnet.
}
{
\begin{itemize}
  \item Atemmuskulatur: Zwerchfell, Zwischenrippenmuskeln
  \item Atemhilfsmuskulatur
  \item Elastische Eigenschaften von Thorax und Lunge
  \item Unterdruck für Atmung wichtig
\end{itemize}
}{}{}{}

\VARIANT{VAR-STANDARD}{\subsection{Atemfrequenz, -tiefe und -volumen}

Die Atemfrequenz, die Atemtiefe und das Atemvolumen
werden  unter
\FullPrettyRef{topic:atemvolumina} besprochen.
}


% \clearpage
\subsection{Atemgase}



\TextSlide{\TxBeschreibung}
{%
\index{Stickstoff}%
\index{Atemgas}%
\index{Sauerstoff!Raumluftgemisch}%
\index{Kohlendioxid!Raumluftgemisch}%
\DefinitionInPara{Als Atemgas bezeichnet man jenes Gasgemisch,
welches eingeatmet wird.} Normalerweise ist das ein Gasgemisch aus Stickstoff,
Sauerstoff, Kohlendioxid und Edelgasen
(\protect\enquote{Raumluft}).
Zwischen der normalen Einatem- und Ausatemluft besteht
besonders hinsichtlich des Sauerstoff- und
Kohlendioxid-Anteils ein Unterschied, siehe
\prettyref{tab:luft-gas}.

\begin{table}[H]
\TableBaseModification
\Captionof{table}{Gasgehalt in der Raum- und
Ausatemluft.\label{tab:luft-gas}}

\begin{tabular}[H]{ccc}
\addlinespace\toprule\addlinespace
\centering \TabHeading{Raumluft} & & \TabHeading{Ausatemluft}
\\\addlinespace\midrule\addlinespace
\EE{21\,\PC*}  & Sauerstoff (\ce{O2}) & \EE{16\,\PC*}
\\\addlinespace
78\,\PC* & Stickstoff (\ce{N}) & 78\,\PC*
\\\addlinespace
\EE{0,03\,\PC*} & Kohlendioxid (\ce{CO2}) & \EE{4\,\PC*} 
\\\addlinespace
1\,\PC* & Edelgase  & 1\,\PC*
\\\addlinespace\bottomrule
\end{tabular}
\end{table}
}
{
\begin{itemize}
  \item Raumluft
  \item Atemgase: Sauerstoff und Kohlendioxid
\end{itemize}
}{}{}{}

\clearpage
\section{Der Verdaungstrakt}
\index{Magen-Darm-Trakt}
\index{Gastrointestinaltrakt}
\index{GIT}
\label{topic:gastrointestinaltrakt}

\HeadingSub{\Lang{Term.}  \T{Gastrointestinaltrakt} (\T{GIT}), 
\Lang{Syn.} {Magen-Darm-Trakt}}


\TextSlide{\TxBeschreibung}
{%
Der Verdauungstrakt dient der Verstoffwechselung der mit dem
Essen aufgenommenen Nährstoffe.
Er beginnt im Mund mit der Nahrungsaufnahme und deren
Zerkleinerung. Durch diverse Enzyme erfolgt die
Weiterverarbeitung im Magen-Darm-Trakt. Wenn sie klein genug
zerlegt sind, werden sie durch die Darmzellen aufgenommen
und in das Blut überstellt. Über das Blut werden die
Nährstoffe an ihre Zielzellen transportiert.
Abfallstoffe landen schlußendlich wieder im Darm (bzw. den
Nieren) um ausgeschieden zu werden.
}
{
\begin{itemize}
  \item Nahrungsaufnahme
  \item Weiterverarbeitung im Magen-Darm-Trakt
  \item Transport der Nährstoffe und Energieträger ins Blut
  \item Zellaufbau und Abbau von Abfallstoffen
\end{itemize}
}{}{}{}

\bigskip

\bigskip

\Layout{27}
{Übersicht}
{}
{}
{./images/anatomie/Digestive_system_diagram_de-edited.pdf}
{Schema des Verdauungstraktes.}
{Mariana Ruiz Villarreal, Public domain}

\Layout{57}
{}
{}
{}
{\begin{tabulary}{\linewidth}[H]{LLL}
\toprule\addlinespace
\noindent\multirow{3}{8em}{\TabHeading{Mundhöhle}}
 &
 Zunge
 &
\noindent\multirow{2}{15em}{Zerkleinerung der Nahrung und Vermischung
mit Speichel}
\\\addlinespace\hhline{~-~}\addlinespace
 &
Zähne 
&
\\\addlinespace\hhline{~-~}\addlinespace
 &
Speicheldrüsen
&
\\\addlinespace\hhline{---}\addlinespace
Rachen
 &
 &
Kreuzung des Speisebreis mit dem Luftweg

Kehldeckel verschließt beim Schlucken die Luftröhre
\\\addlinespace\hhline{---}\addlinespace
\noindent\EE{Speiseröhre} (\T[Oesophagus@Ösophagus]{Ösophagus})
 &
 &
Speisebrei wird mittels Muskelbewegung zum Magen befördert
\\\addlinespace\hhline{---}\addlinespace
\noindent\TabHeading{Magen}
 &
 &
Vermischung mit Magensäure

Magensäure zerlegt Nahrung

Pförtner gibt portionsweise Speisebrei an den Zwölffingerdarm ab

\\\addlinespace\hhline{---}\addlinespace
\multirow{4}{8em}{\TabHeading{Dünndarm}}
&
\noindent\TabSub{Zwölffingerdarm} \Z{Duodenum}
 &
Gallenflüssigkeit (${\leftarrow}$Leber) und Bauchspeichel
(${\leftarrow}$Pankreas) wird zugesetzt
\\\addlinespace\hhline{~--}\addlinespace
 &
\noindent\TabSub{Leerdarm} \Z{Jejunum}
 &
Spaltung von Nährstoffen (Eiweiß, Kohlehydrate, Fett) und
Nährstoffaufnahme
\\\addlinespace\hhline{~-~}\addlinespace
 &
\noindent\TabSub{Krummdarm} \Z{Ileum}
 &
\\\addlinespace\hhline{---}\addlinespace
\multirow{6}{8em}{\TabHeading{Dickdarm}\\(\TabHeading{Colon})}
&
\noindent\TabSub{Blinddarm}
&
Sackgasse
\\\addlinespace
&
\noindent\TabSub{Wurmfortsatz} (\TabSub{Appendix}) 
 &
Anhängsel der Sackgasse, kann sich entzünden.
\\\addlinespace\hhline{~--}\addlinespace
 &
\noindent\TabSub{Aufsteigender} Teil
 &
\multirow{4}{\columnwidth}{Wasserentzug, \\Eindickung, \\
Bakterienbesiedlung, \\Produktion und Speicherung des
Stuhls} 
\\\addlinespace\hhline{~-~}\addlinespace
 &
\noindent\TabSub{Querender} Teil
 &
\\\addlinespace\hhline{~-~}\addlinespace
 &
\noindent\TabSub{Absteigender} Teil
 &
\\\addlinespace\hhline{~-~}\addlinespace
 &
\noindent\TabSub{S-förmig} gebogener Teil (\Z{Sigma})
 &
\\\addlinespace\midrule\addlinespace
\noindent\TabHeading{Mastdarm} (\TabHeading{Rektum})
 &
 &
\protect\enquote{Speicherung} des ausscheidungsbereiten
Stuhls
\\\addlinespace\hhline{---}\addlinespace
\noindent\TabHeading{Anus}
 &
 &
Absetzen des Stuhls

\\\addlinespace\bottomrule
\end{tabulary}}
{Übersicht über den Verdauungstrakt --
\emph{\protect\enquote{Von der Schüssel in die Schüssel.}}\label{tab:verdauungstrakt-uebersicht}}
{}


\clearpage % TEMP temp clearpage

\Layout{20}{Mundhöhle}
{\index{Mundhöhle}%
In der Mundhöhle kommt es durch die Zähne zur
\EE{Zerkleinerung} der Nahrung. Durch die Kaubewegungen der
Zähne und durch die Bewegungen der Zunge wird der Brei mit
Speichel aus den Speicheldrüsen vermischt.
\VARIANT{VAR-ERNAEHRUNGSPAED}{Insgesamt 3 verschiedene
Speicheldrüsen fügen ihr Sekret dem Speisebrei zu
(Ohrspeicheldrüse, Unterzungenspeicheldrüse,
Unterkierferspeicheldrüse).} Der Speichel beginnt schon im
Mund mit der Aufspaltung von Stärke und Zucker der Speise.

Der Erwachsene besitzt bis zu 32 Zähne.
\VARIANT{VAR-ERNAEHRUNGSPAED}{Davon befinden sich jeweils 16
in Ober- bzw. Unterkiefer (4x 2 Schneidezähne, 1 Eckzahn, 2
Vormahlzähne, 3 Mahlzähne).} Diese Anzahl kann insbesondere
durch die unterschiedliche Ausprägung der hintersten
Mahlzähne (\protect\enquote{Weisheitszähne}) variieren. Das
Milchgebiss unterscheidet sich durch die Anzahl
\VARIANT{VAR-ERNAEHRUNGSPAED}{(insg. 20, 4x 2 Schneidezähne,
1 Eckzahn, 2 Mahlzähne)} von den bleibenden Zähnen.

Die Zunge schiebt den Bissen Schritt für Schritt in den
Rachen. Beim Schlucken sorgt sie dafür, dass der Kehldeckel
verschloßen wird, da es sonst zu einem Übertritt der Nahrung
in die Luftröhre und damit zu einer \I{Aspiration} kommen
kann. }
{
\begin{itemize}
    \item Zähne: Nahrungszerkleinerung
    \item Speicheldrüsen
    \item Zunge: Nahrungsfortbewegung, Kehldeckelverschluss
\end{itemize}
}
{./images/hirtler-aass/needs-work/Mundhoehle}
{Die Mundhöhle}
{Hirtler}

\VARIANT{VAR-ERNAEHRUNGSPAED}{
\subsection{Allgemeiner Wandaufbau des Verdauungstrakts}
\TextSlide{\TxBeschreibung}{
Der gesamte Verdauungstrakt ist durch einen grundsätzlich
einheitlichen Wandaufbau gekennzeichnet. Dieser besteht aus
4 Schichten (von innen nach außen):

\begin{itemize}
  \item Schleimhaut (Mucosa)
  \item Verschiebeschicht (Submucosa)
  \item Muskelschicht (Muscularis): hier kann man zwei Faserrichtungen
  unterscheiden - zirkulär- und längsgerichtet
  \item Bauchfell (Peritonaeum)
\end{itemize}

Oberhalb des Zwerchfells findet man kein Bauchfell, da dies nur den Bauchraum
auskleidet.
}{
\begin{itemize}
  \item Schleimhaut
  \item Verschiebeschicht
  \item Muskelschicht
  \item Bauchfell
\end{itemize}
}
}



\Layout{20}{Speiseröhre}
{%
\index{Speiseröhre}%
\label{topic:speiseroehre}%
\Lang{Term.} \T[Oesophagus@Ösophagus]{Ösophagus}.
Die Speiseröhre (\T[Ösophagus@Oesophagus]{Ösophagus}) ist
ein ungefähr 25\CM* langer Muskelschlauch. Er dient dem
\EE{Transport} über den Mund aufgenommener Nahrung vom
Rachen in den Magen. Diese Fortbewegung funktioniert über
eine wellenförmige Kontraktion der Wandmuskeln der
Speiseröhre. Im Ösophagus werden keine verdauenden Enzyme
ausgeschieden -- er dient rein zum Weitertransport des
Speisebreies.

\Z{Die Speiseröhre weist \E{drei Engstellen} auf, welche die
Nahrung überwinden muss. Die erste ist hinter dem Kehlkopf
(wenn hier ein Speisebrocken stecken bleibt, kann er
außerdem die Luftröhre verschließen. Landläufig nennt man
das den \Speech{\protect\enquote{Wiener-Würschtel-Tod}}.).
Die zweite Enge liegt dort, wo die Aorta die Speiseröhre
kreuzt. Die dritte Enge entspricht dem Durchtritt des
Ösophagus durch das Zwerchfell.}

In der Wand des Ösophagus finden sich zahlreiche Venen. Bei
manchen Lebererkrankungen kann es zu einer Stauung kommen;
es entwickeln sich sog.
\I[Oesophagusvarizen@Ösophagusvarizen]{Ösophagusvarizen}.\ZFootnote{Z.\,B. fortgeschrittene
Leberzirrhose: Umbau der Leber in Bindegewebe durch z.\,B.
chronische Entzündung und damit resultierendem Verschluss
der Pfortader. Dabei fungieren diese Venen als
Umgehungskreislauf in die Vena cava.} Wenn eine solche
Ösophagusvarize platzt, kann es zu einem großen Blutverlust
kommen (\prettyref{topic:oesophagusvarizenblutung}).
}{
\begin{itemize}
    \item $\sim$ 25\CM* langer Muskelschlauch
    \item Nahrung: Rachen $\rightarrow$ Magen
    \item Keine Verdauung
    \item Erkrankung: Ösophagusvarizen
\end{itemize}
}
{./images/wikipedia-pd/stomach.pdf}
{Ösophagus, Schema.}
{PD}






\TextSlide{Magen}
{%
\index{Salzsäure}%
\index{Magen!-pförtner}%
\index{Magen}%
\Z{\Lang{Term.} Gaster, Ventriculus.}.
Der Magen ist ein muskuläres Hohlorgan und dient der
Nahrungsspeicherung, Nahrungsdurchmischung sowie der
kontrollierten Weiterleitung des durchmischten Speisebreis.

Verantwortlich für den Beginn der Verdauung der Nahrung ist
der durch die Magenschleimhaut gebildete saure, aus
\EE{Salzsäure} bestehende, \EE{Magensaft}. Die
Magenschleimhaut produziert ausserdem eine \EE{schützende
Schleimschicht}, die die Oberfläche vor der starken Säure
schützt.

Nachdem die Wandmuskulatur den Speisebrei mit dem Magensaft
\EE{durchmischt} hat, gibt der \EE{Magenpförtner}
\Z{(Pylorus)} den Speisebrei portionsweise an den Dünndarm
-- genauer: an den Zwölffingerdarm -- weiter.
}{
\begin{itemize}
  \item Magensaft mit Salzsäure
  \item Schleim schützt Magenschleimhaut
  \item Wandmuskulatur durchmischt Speisebrei
  \item Portionsweise Abgabe von Speisebrei durch den
  Magenpförtner in
  den Zwölffingerdarm
\end{itemize}
}{}{}{}






\TextSlide{Zwölffingerdarm}{%
\index{Zwölffingerdarm}%
\Lang{Term.} \T{Duodenum}.%
Der Zwölffingerdarm ist der Beginn des Dünndarmes ab dem
Pförtner des Magens und misst ca. 25\CM*. Sein Verlauf
ähnelt einem C.

Hier \EE{münden 2 Ausführungsgänge} ein, der \EE{Gallengang}
und der \EE{Pankreasgang}. Aus dem Gallengang wird die
\EE{Gallenflüssigkeit} aus der \EE{Leber} abgegeben, welche
für die \EE{Fettverdauung} wichtig ist (Emulsion der Fette);
aus dem \EE{Pankreas} wird der \EE{Bauchspeichel}, der
Verdauungsenzyme beinhaltet, zugeführt.

Die Kenntnis des gemeinsamen Ausführungspunktes von Galle
und Pankreas ist wichtig, um pathologische Vorgänge dieser
beiden Organe zu verstehen. Ein Gallenstein beispielsweise,
welcher im Ausführungsgang nahe der Mündung in den
Zwölffingergang stecken bleibt, führt durch Rückstau der
Pankreasflüssigkeit zu einer akuten Entzündung der
Bauchspeicheldrüse.
}
{%
\begin{itemize}
  \item Beginn des Dünndarmes, ca. 25\CM* lang
  \item 2 Gänge münden:
  \begin{itemize}
    \item Gallengang (Gallenflüssigkeit aus Leber, Fettverdauung)
    \item Pankreasgang (Bauchspeichel, Enzyme)
  \end{itemize}
\end{itemize}
}{}{}{}



\opt{STATUS-EXPERIMENTAL}{
\Parallel{

\TODO{GABS}{2010-04-11}{Lage des Duodenums: Bildfehlt.}{}

% \includegraphics[width=\linewidth]{./images/noauth/AASdev20090228-img48.png}
\includegraphics[width=\linewidth]{./icons/under-construction.png}

\Captionof{figure}{\AuthNotesPicLegalNotOk{}{}Die Lage des
Zwölffingerdarms im Körper.} 

}{

\TODO{GABS}{2010-04-11}{Duodenum und Pankreas: fehlt.}{}

\includegraphics[width=\linewidth]{./icons/under-construction.png}
% \includegraphics[width=\linewidth]{./images/noauth/AASdev20090228-img49.png}
\Captionof{figure}{\AuthNotesPicLegalNotOk{}{}Der
Zwölffingerdarm und die Bauchspeicheldrüse.} 

}
}


\TextSlide{Leber und Gallenblase}{%
\index{Leber}%
\index{Gallenblase}%
\label{topic:gallenblase}%
\Lang{Term.} Hepar. Die Leber ist ein solides Organ und befindet sich im
\EE{rechten Oberbauch}. Sie hat einen linken und rechten
Lappen. Sie ist die \protect\enquote{Chemiefabrik} des Körpers: Sie
produziert wichtige Stoffe, entgiftet und speichert Fette
und Zucker.

Die Leber produziert u.\,a. wichtige, vom Körper benötigte,
\EE{Eiweißstoffe} (Proteine).
Weiters sorgt sie für den \EE{Abbau von Giftstoffen}
(\EE{Entgiftung}). Die Leber nimmt dabei durch die
Giftstoffe mit der Zeit selbst auch Schaden, welches zu
chronischen Leberschädigungen, wie z.\,B. einer
Leberentzündung (Hepatitis) oder Leberzirrhose (könnte man
als fortschreitende \protect\enquote{Vernarbung} der Leber
bezeichnen), führen kann. Über die \T{Pfortader} (Vena
portae, \ref{topic:pfortader}) erhält die Leber
nährstoffreiches Blut von Teilen des Darms um es zu
entgiften, bevor es in den normalen Blutkreislauf
zurückkehrt.
Außerdem dient die Leber der \EE{Speicherung} von Fetten und Zucker.

Die Leber ist auch eine \EE{Verdauungsdrüse}: Sie produziert
die \T{Gallenflüssigkeit}
(\Speech{\protect\enquote{Galle}}), welche der 
\EE{Fettverdauung} dient.
Von der Leber leiten die Gallengänge\index{Gallengang} die Gallenflüssigkeit
in den Zwölffingerdarm. An den Gallengängen hängt die
\T{Gallenblase}, in welcher die Gallenflüssigkeit
\E{gespeichert (!)} -- nicht produziert -- wird. Dadurch
entsteht dem Körper kein Schaden, wenn man diese
beispielsweise wegen Gallensteinen entfernen muss.
}{
\protect\enquote{Chemiefabrik} des Körpers
\begin{itemize}
  \item \protect\enquote{Chemiefabrik}:
  \begin{itemize}
    \item Produktion von  Eiweißen
    \item Entgiftung
    \item Nährstoffreiches Blut vom Darm zwecks Entgiftung
  \end{itemize}
  \item Speicherung von Fett und Zucker
  \item \protect\enquote{Verdauungsdrüse}
  \begin{itemize}
    \item Produziert Gallenflüssigkeit  →
    Fettverdauung
    \item Gallengänge Richtung Zwölffigedarm
    \item An Gallengängen: Gallenblase (Speicherung)
  \end{itemize}
\end{itemize}
}{}{}{}

 


\begin{PictureSeriesWide}{}{Bilderserie: Leber und
Gallenblase}
\PictureSeriesRowWide{3}
{./images/wikipedia-pd/Liver_and_pangreas_-_transparent.png}
{Leber, Gallenblase und
Pankreas.
\SourcePicture{WMC}}
{./images/hirtler-aass/Gallenblase.pdf}
{Gallengang, Gallenblase und Einmündung in den
Zwölffingerdarm.
\SourcePicture{Hirtler}}
{./images/hirtler-aass/Pankreas-1.pdf}
{Lage des Pankreas zum Zwöffingerdarm. Deutlich zu erkennen
ist die Einmündung des Pankreas- und des Gallenganges in
den Darm. \SourcePicture{Hirtler}}
{empty}
{}
\end{PictureSeriesWide}




\Layout{60}{Bauchspeicheldrüse}
{%
\index{Bauchspeicheldrüse}%
\index{Pankreas}%
\Lang{Term.} \T{Pankreas}. Das Pankreas ist ca. \VonBis{15}{22}\CM* lang und liegt
hinter dem Magen. Es hat einen \EE{Ausführungsgang}, welcher
zusammen mit dem Gallengang \EE{in den Zwölffingerdarm}
(Duodenum) mündet. Man kann es grob unterteilen in Kopf und
Schwanz. Der Pankreaskopf liegt in der C-Schlaufe des
Zwölffingerdarmes, die Spitze des Pankeas reicht bis zur
Milz.

\subparagraph{Aufgaben} 
Das Pankreas hat \EE{zwei wesentliche Funktionen}.
Einerseits ist es verantwortlich für die \EE{Enzymproduktion für
die Verdauung}. Diese werden über den
Pankreas-Ausführungsgang in den Zwölffingerdarm abgegeben.
Andererseits produziert das Pankreas zwei wichtige \EE{Hormone},
\I{Insulin} \Z{und Glukagon}, die  den
\EE{Zuckerstoffwechsel} des Körpers regulieren. Insulin sorgt eine
Verringerung des Blutzuckers durch eine Förderung der
Zuckerentfernung aus dem Blut. \Z{Glukagon hingegen ist
verantwortlich für die Steigerung der Menge an Zucker im
Blut.}
}
{
\begin{itemize}
  \item Produktion von
  \begin{itemize}
    \item Verdauungsenzymen
    \item Hormonen für den Zuckerstoffwechsel:
    
    Insulin \Z{und Glukagon}
  \end{itemize}
\end{itemize}
}
{./images/hirtler-aass/Pankreas-1.pdf}
{Pankreas}
{Hirtler}




\Layout{60}{Leer- und Krummdarm}
{%
\index{Krummdarm}%
\index{Leerdarm}%
\Lang{Term.} \T{Jejunum}, \T{Ileum}. Jejunum (Leerdarm) und Ileum (Krummdarm) sind die zwei
weiteren Teile des Dünndarmes. Ihre Gesamtlänge beträgt ca.
3\Meter* . Ihre Hauptaufgabe besteht in der weiteren
Verdauung der Nahrung und der Nährstoffaufnahme. Hierfür
dienen Darmzotten -- kleine \protect\enquote{Ausstülpungen
der Oberfläche} -- zur Oberflächenvergrößerung.

\label{topic:mesenterium}
Jejunum und Ileum sind ab ihrem Abgang aus dem Duodenum bis
zum Übergang in den Dickdarm relativ frei beweglich. Die
einzige Befestigung dieses Darmabschnittes ist das
\T{Mesenterium} \Z{(Darmgekröse)}. Es befindet sich ungefähr
auf Höhe des Nabels und enthält die \E{Blutgefäße}, die zum
Darm hin und vom Darm weg führen.

\Z{Das besondere am Ileum ist, dass es einzelne Zellbezirke
des Abwehrsystems (Peyer'sche Plaques) enthält. Sie sorgen
dafür, dass Bakterien aus dem Dickdarm nicht den weiteren
Dünndarm bevölkern können. Wenn dies nicht gelingt,
resultiert eine Entzündung des Darmes (Enteritis bis
Gastroenteritis) mit den Symptomen von Durchfall, Übelkeit
und Erbrechen.}
}
{
\begin{itemize}
    \item Ca. 3\Meter* lang, Oberflächenvergrößerung durch Darm\-zotten
    \item \EE{Nährstoffaufnahme}\index{Nährstoffaufnahme}
    \item \Z{beweglich, am Darmgekröse (Mesenterium) aufgehängt}
    \item \Z{Abwehr von Darmbakterien}
\end{itemize}
}
{./images/wikipedia-pd/Small_intestine.pdf}
% {./images/noauth/AASdev20090228-img56.png}
{Dünndarm.}
{WmCo/Mikael Häggström, PD.}

\begin{PictureSeriesWide}{}{Bilderserie: Bauchsitus und
Mesenterium}
\PictureSeriesRowWide{3}
{./images/hirtler-aass/Bauchfilm_4032-AASS-0112mm}
{Der Bauchsitus mit aufgeklappter
Bauchdecke (links = kopfwärts). Am Rand sieht man deutlich
den Dickdarm, unter der Fettschürze zeichnen sich die
Schlingen des Dünndarms ab. \SourcePicture{Hirtler}}
{./images/hirtler-aass/Duenndarm-Mesenterium}
{Schema: Dünndarm mit Aufhängung
(Mesenterium) und Blutgefäßen. \SourcePicture{Hirtler}}
{./images/hirtler-aass/Bauchfilm_Image_022}
{Foto: Dünndarm mit Aufhängung
(Mesenterium) und Blutgefäßen. \SourcePicture{Hirtler}}
{}
{}
\end{PictureSeriesWide}





\Layout{60}{Dickdarm}{%
\index{Dickdarm}%
\Lang{Term.} \T{Colon}. 
Der Dickdarm (Colon) ist ca. 1,5\Meter* lang und
wird in verschiedene Abschnitte unterteilt: 

\begin{center}
\TableBaseModification
\begin{tabulary}{\linewidth}[H]{LLL}
\toprule\addlinespace
\multirow{6}{8em}{\EE{Dickdarm}\\(\T{Colon})}
&
\EE{Blinddarm} (\T{Coecum})\index{Blinddarm}
&
Sackgasse
\\\addlinespace
&
\EE{Wurmfortsatz} (\T{Appendix})
\index{Wurmfortsatz}\index{Appendix} &
Anhängsel der Sackgasse, kann sich entzünden.
\\\addlinespace\hhline{~--}\addlinespace
 &
\EE{Aufsteigender} Teil (\T{Colon ascendens})
&
\multirow{4}{\columnwidth}{Wasserentzug, Eindickung, \\
Bakterienbesiedlung, \\Speicherung} 
\\\addlinespace\hhline{~-~}\addlinespace
 &
\EE{Querender} Teil (\T{Colon transversum})
 &
\\\addlinespace\hhline{~-~}\addlinespace
 &
\EE{Absteigender} Teil (\T{Colon descendens})
 &
\\\addlinespace\hhline{~-~}\addlinespace
 &
\EE{S-förmig} gebogener Teil (\T{Sigma})
\index{Dickdarm!S-förmig gebogener Teil} &
\\\addlinespace\hhline{~-~}\addlinespace
\EE{Mastdarm} (\T{Rektum})\index{Mastdarm}&\protect\enquote{Speicherung} des
ausscheidungsbereiten Stuhls
\\\addlinespace\bottomrule
\end{tabulary}
\end{center}

Der Leerdarm (Ileum) mündet in den Dickdarm. Diese Mündung
ist wie eine \E{T-Kreuzung} gebaut: \E{Kopfwärts} geht es in
den \EE{aufsteigenden Teil} des Dickdarmes, \E{zehenwärts}
in den \EE{Blinddarm}. Wie der Name des Blinddarmes schon
sagt endet er blind, er ist daher eine Sackgasse. Am
Blinddarm angehängt findet man den \EE{Wurmfortsatz}
(\T{Appendix}) welcher sich häufig entzündet
(\T{Appendiz\underline{itis}}\footnote{Das Suffix
\emph{\protect\enquote{-itis} zeigt immer eine Entzündung
an, vgl. Kapitel Hygiene}}). Der Wurmfortsatz erfüllt beim
Menschen keine nennenswerte Funktion, außer der
Arbeitsplatzsicherung von Anästhesisten und Chirurgen (wegen
der vielen Blinddarm-Operationen). Es wird vermutet, dass er
eine immunologische Funktion hat.

Vom Blinddarm geht es weiter in den \EE{aufsteigenden}, den
\EE{querenden} und den \EE{absteigenden Teil des Dickdarms}
\Z{(Colon ascendens, transversum und descendens)}. Hierbei
durchquert der Dickdarm alle Quadranten des Bauches, von
rechts unten über rechts oben nach links oben bis links
unten. Es folgt der \EE{S-förmig gebogene Teil}
\Z{(Sigmoideum)} und danach das \EE{Rektum}, welches den
Stuhl über den \EE{Anus} nach außen entlässt.

Die Hauptfunktion des Dickdarms ist das Eindicken des
Stuhls. Der Großteil der Wasseraufnahme aus der verdauten
Nahrung findet hier statt. \Z{Im Dickdarm befinden sich
Darmbakterien (Escherichia coli).}
}
{
\begin{itemize}
  \item T-Kreuzung
  \item Blinddarm = Sackgasse
  \item Wurmfortsatz
  \item Funktion des Dickdarms: Stuhleindickung
\end{itemize}
}
{./images/wikipedia-pd/Colon_anatomy.pdf}
{Der Dickdarm}
{Mariana Ruiz Villarreal,
    Public domain}


% \clearpage % TEMP temp clearpage
\section[Sonstige Strukturen im Bauchraum]{Sonstige Organe und Strukturen im Bauchraum}
\label{topic:sonstige-strukturen}
\label{topic:sonstige-bauchorgane}

\subsection{Bauchfell}
\index{Bauchfell}

\HeadingSub{\Lang{Term.} \T{Peritoneum}}
% REVIEW Bauchfell: Überarbeitung
\A{Review: Bauchfell: Überarbeitung}
\A{Review: Kommentare eingefügt.}


\Layout{20}
{\TxBeschreibung}
{%
Das Bauchfell (\T{Peritoneum}) ist eine dünne Haut, welche
\EE{die gesamte Bauchhöhle auskleidet}. Dieser Überzug
bedeckt die meisten Bauchorgane\ZFootnote{außer z.\,B. die
Nieren} \Z{und das große und kleine Netz (Omentum majus und
minus). Das große Netz bedeckt die Bauchorgane und ist die
erste Struktur, die man bei einer Baucheröffnung sieht. Es
ermöglicht die Abkapselung von Entzündungen, sodass sich
diese nicht sofort auf den ganzen Bauchraum ausbreiten
können.
}
\ACh{GABS}{2010-03-09}{Unter großem und kleinen Netz wird sich kaum wer was vorstellen können \ldots} 

Es sorgt durch seine feuchte und glatte Oberfläche dafür,
dass die Organe, wie z.\,B. Darm, ihre Bewegungen zur
Verdauungen besser durchführen können. Das Bauchfell ist
auch verantwortlich für die \EE{Schmerzen}, die z.\,B. bei
einer Entzündung der Verdauungsorgane (z.\,B.
Blinddarmentzündung (Appendizitis)) entstehen. Es kann sich
auch selbst entzünden (\E{Bauchfellentzündung} (Peritonitis,
\FullPrettyRef{topic:bauchfellentzuendung})), z.\,B. durch eine
Infektion oder wenn es bei einem Durchbruch
(\E{Perforation}) des Darmes zu einem Eintritt von
Darminhalt in die Bauchhöhle kommt.
}{
\begin{itemize}
    \item Auskleidung der gesamten Bauchhöhle
    \item dünner Überzug
    \item Bewegung der Organe
    \item schmerzempfindlich
    \item Bauchfellentzündung (Peritonitis)
\end{itemize}
}
{./images/gabriel-sebastian-aass/Torso_GABS_IMG_1251_00800_frei.jpg}
{Der Brust- und Bauchraum, geöffnet am Torso in Aufsicht. Zu
sehen sind die Lunge, dahinter das Herz, darunter die Leber
(braun) mit einem kleinen Stück der Gallenblase (grün), der
Magen mit einem Teil der gelben Fettschürze
\Z{\protect\enquote{großes Netz}}, sowie der Dickdarm und
der Dünndarm. } 
{GABS.}


\subsection{Die Milz} 
\index{Milz}

\HeadingSub{\Z{\Lang{Term.} Lien, Splen}}

\TextSlide{\TxBeschreibung}
{\index{Kapsel!Milz}
Die Milz ist ca. 12\CM* groß, oval geformt und und liegt im
linken Oberbauch, verdeckt von den Rippen. Sie hat nichts 
mit der Verdauung  zu tun.

\Z{Die Milz dient der immunologischen Überwachung des
Blutes und ist daher Teil des Immunsystems. Sie filtert das
Blut, zu alte Erythrozyten werden durch weiße Blutkörperchen in der
Milz erkannt und entfernt. Beim Embryo ist sie auch an
der Blutbildung beteiligt. Ihre Vene mündet in
die Pfortader.} Damit sie ihre Aufgaben erfüllen
kann, ist sie \EE{sehr gut durchblutet}! 
Sie ist weiters von einer festen \EE{Bindegewebskapsel}
umgeben und in das Bauchfell gehüllt.
}
{
\begin{itemize}
  \item 12\CM*, oval, im linken Oberbauch unter den Rippen gelegen
  \item Von \E{Kapsel} umgeben
  \item Teil des Immunsystems, entfernt alte Erythrozyten
  \item Gut durchblutet!
\end{itemize}
}
{./icons/under-construction.png}
% {./images/noauth/AASdev20090228-img60.png}
{\AuthNotesPicLegalNotOk{img60}{}Milz:
Lage.}
{Quelle nicht eruierbar.}
 
\TextSlide{\TxGefahren{}: \protect\T{Zweizeitiger Milzriss}}
{%
\index{Zweizeitiger Milzriss}%
\index{Milz!-riss, zweizeitiger}%
Aus den Tatsachen, dass die Milz gut durchblutet ist und
über eine feste Kapsel verfügt, ergibt sich bei Verletzungen
eine besondere Gefahr:
Es kann vorkommen, dass bei einem Unfall die Milz einreisst,
die Kapsel aber erhalten bleibt: Es blutet zuerst in die
Kapsel ein und der Druck steigt, die Kapsel wirkt als eine
Art Druckverband. Nach einer gewissen Zeit reißt jedoch auch
die Kapsel, der Patient blutet nun \E{plötzlich} ungehindert
in den Bauchraum. \Sign{→} \Ca{Der Patient wirkt zuerst
\protect\enquote{ganz normal}, unerwartet und plötzlich
\protect\enquote{verfällt} er.} }{%
\begin{itemize}
  \item Milz reisst, Kapsel intakt: langsame Blutung
  \item Druck in Kapsel steigt: Kapsel reisst auch
  \item Plötzlich massive, ungehinderte Blutung im Bauchraum
  \item \protect\enquote{Aus heiterem Himmel}
\end{itemize}
}{}{}{}


\subsection{Die Nieren}

\HeadingSub{\Lang{Term.} Ren, \Z{Nephros}}
\index{Niere}

\Layout{20}{\TxBeschreibung}
{%
Die Niere ist ein paarig angelegtes Organ und befindet sich
links und rechts neben der Wirbelsäule auf Höhe der unteren
Rippenpaare. Durch die Lage der Leber liegt die rechte Niere
etwas weiter fußwärts als die linke Niere. Die Nieren
besitzen wie die Milz eine Bindegewebskapsel.

Die Niere \EE{filtert ausscheidungspflichtige Stoffe}
(z.\,B. Stickstoff und Abfallstoffe von Medikamenten) aus
dem Blut. Daher ist auch sie äußerst gut durchblutet und
reagiert sehr empfindlich auf einen Blutdruckabfall. Aus
diesem Grund hat sie auch einen Einfluss auf die
Blutdruckregulation.
Weiters ist sie verantwortlich für die Verhältnisse der
\EE{Elektrolytkonzentrationen} im Blut und in den Geweben
und damit für den \EE{Säure-Basen-Haushalt} im Körper (pH)
sowie für den \EE{Wasserhaushalt}.
Außerdem bildet die Niere Hormone für die
Blutdruckregulierung \Z{(Renin)} und für die Blutbildung
\Z{(Erythropoetin)}.

Mit zunehemendem Alter, insbesondere bei Diabetikern, ist
die Niere eines der Hauptorgane, die \protect\enquote{den
Geist aufgeben}. Daraus resultiert eine \E{chronische
Niereninsuffizienz}, die Patienten müssen regelmäßig zur
Dialyse (\protect\enquote{Blutwäsche}, vgl.
\FullPrettyRef{topic:niereninsuffizienz}).
Außerdem zeigen diese Personen häufig eine Anämie
(Blutarmut) durch den Mangel an dem Blutbildungshormon
\Z{Erythropoetin}.

}{
\begin{itemize}
  \item paarig angelegtes Organ
  \item stark durchblutet, empfindlich auf RR- Abfall ${\rightarrow}$
  Funktion ${\downarrow}$
  \item reguliert
  \begin{itemize}
    \item \EE{Wasser-} und
    \EE{Elektrolythaushalt}
    \index{Wasserhaushalt!Niere}
    \index{Elektrolythaushalt!Niere}
    \item \EE{Säure/Basen-Haushalt}
    \index{Säure-Basen-Haushalt!Niere}
    \item Blutdruck
  \end{itemize}
  \item \Z{Hormonproduktion} (Renin, Erythropetin)
\end{itemize}
}
{./images/wikipedia-pd/Kidneys.png}
% {./images/noauth/AASdev20090228-img61.png}
{Die Nieren.}
{WmCo/Mikael Häggström and Madhero88, PD.}




\TextSlide{Aufbau der Nieren}
{%
Die Niere hat eine bohnenförmige Form (ca. $11 \times 5,5
\times 4$\CM*). Wenn man sie längs durchschneidet kann man
ihren Aufbau erkennen. Sie ist unterteilbar in ein
Nierenmark und eine Nierenrinde. Das Nierenmark ist
pyramidenförmig, mit der Spitze nach innen. In ihm sammeln
sich die Nierenkanäle. An der Spitze des Marks folgt der
Übergang in das Nierenbecken und somit der Anschluss an den
Harntrakt. Als Stiel \Z{(Hilus)} wird der Bereich der Niere
bezeichnet, wo die Nierenarterie, die Nierenvene und der Harnleiter
befestigt sind.

\subparagraph{\Z{Mikroskopischer Aufbau}}
\Z{Die funktionelle Einheit der Niere ist das Nephron. Es
gibt davon ca. 1 Mio in jeder Niere. Im Nephron kann man unterscheiden:
\begin{itemize}
  \item Glomerulus: Hier wird das Blut gefiltert und Plasmawasser
  abgepresst (ca.
  150\,\nicefrac{\Liter}{\Day})
  \item Schleifenapparat (Tubulus): Das abgepresste
  Plasma\-wasser wandert den Schleifen\-apparat entlang.
  Hier werden die meisten wichtigen Stoffe und der
  Groß\-teil des Wassers wiederum zurückgeholt, sodass am
  Ende nur mehr \VonBis{1,5}{2}\Liter* Harn pro Tag in der
  Harnblasen landen.
\end{itemize}
}
}
{%
\begin{itemize}
  \item Rinde, Mark
  \item Hilus (Stiel)
  \item Gefäße
  \item Nierenbecken mit Anschluss an den
  Harntrakt\index{Nierenbecken}
\end{itemize}
}
{}{}{}



\Text
{Funktionsweise}
{\Z{%
In der Niere werden pro Tag 1\,500\Liter* Blut
\protect\enquote{bearbeitet}. Das Blut wird in den Glomerula
zunächst gefiltert. Dadurch entstehen 150\Liter* Primärharn
(inkl.
Elektrolyte usw.). Dieser fließt dann weiter durch den
Schleifenapparat. Dort werden für den Körper wichtige Stoffe
wieder aufgenommen und Wasser rückresorbiert, sodass aus
150\Liter* Primärharn schlußendlich \VonBis{1,5}{2}\Liter*
Endharn pro Tag den
Körper verlassen. 
\cite{LoefflerBiochemie7}}}
{}
{}
{}
{}


\subsection{Harnableitendes System -- Harntrakt}
\index{Harn!-ableitendes System}\index{Harn!-trakt}

\Layout{20}{\TxBeschreibung}
{%
Der Harntrakt besteht aus den paarig angelegten
Nieren mit den Nierenkelchen und den Nierenbecken, auf jeder
Seite je ein  \II{Harnleiter} \Z{(Ureter)}, welcher ca.
30\CM* lang ist und in die
\II{Harnblase} mündet. Von der Blase leitet die
\II{Harnröhre} (\Z{Urethra}) den Harn in die Aussenwelt.

Die Harnröhre ist beim \EE{Mann} ca. 25\CM* lang und
verläuft von der Harnblase durch die \EE{Prostata}
\index{Harn!-röhre!Verlauf} und den Penis und mündet an der
Eichel in die Außenwelt. Bei einer Vergrößerung der Prostata
kann der Harnabfluss bis zur Unpassierbarkeit erschwert sein
(\E{Harnverhalt}\VARIANT{VAR-STANDARD}{,
\prettyref{topic:harnverhalten}}).

Bei der \EE{Frau} ist die Harnröhre mit ungefähr 4\CM*
\E{deutlich kürzer} als beim Mann und mündet oberhalb des
Scheideneinganges. Durch die kürzeren Harnwege bei der Frau kann es 
leichter zu \E{Harnwegsinfektionen} kommen.
}
{
\begin{itemize}
    \item Nieren inkl. Nierenbecken
    \item \index{Harn!-leiter}je 1 Harnleiter 
    \Z{(Ureter)}, 30\CM* lang, mündet  in die 
    \item \index{Blase}Blase \index{Harn!-blase}
    \item Harnröhre  (\Z{Urethra,} Mann: 25\CM*, Frau
    4\CM*)\index{Harn!-röhre}
\end{itemize}
}
{./images/hirtler-aass/Niere.pdf}
{\AuthNotesPicLegalOk{}{Lena}Harntrakt
mit Nieren (rechte Niere im Querschnitt mit Nierenbecken),
Harnleiter, Blase und Harnröhre. \label{fig:harntrakt}}
{Lena Hirtler}




\subsection{Die Nebenniere(n) (Adren)}


\TextSlide{\TxBeschreibung}
{%
Sie sind ein paariges, jeweils am oberen Nierenpol
gelegenes Organ und überdecken die Niere \enquote{mützenformig}.
Mit der Harnproduktion haben sie \E{nichts} zu tun. 

Die Hauptaufgabe der Nebennieren ist die
\EE{Hormonproduktion}. \Z{Man kann die Nebenniere im
Querschnitt in Nebennierenrinde und -mark unterteilen. 
In den jeweiligen Abschnitten werden jeweils
verschiedene Hormone produziert.} Die wichtigsten
produzierten Hormone sind:

\begin{itemize}
    \item Adrenalin, Noradrenalin
    \item Sexualhormone \Z{(Androgene)}
    \item \Z{Kortisol}
    \item \Z{Aldosteron}
\end{itemize}

}{%
\begin{itemize}
  \item Hormonproduktion (Adrenalin, Sexualhormone, \ldots)
\end{itemize}
}{}{}{}




\subsection{Was gibt es noch?}

Außer den \protect\enquote{klassischen Bauchorganen} liegen auch andere
Organsysteme im Bauchraum auf die nicht vergessen werden
sollte. Besonders betrifft das die weiblichen und männlichen
Fortpflanzungsorgane wie Eierstöcke, Eileiter und Uterus
bzw. Prostata und Samenbläschen. Diese werden allerdings im
Kapitel \enquote{Geschlechtsorgane} gesondert besprochen.

\clearpage
\section{Geschlechtsorgane}

\VARIANT{VAR-STANDARD}{
Bei Mann und Frau kann man zwischen äußeren und inneren Geschlechtsorganen
unterscheiden.
}


\VARIANT{VAR-ERNAEHRUNGSPAED}{
Global gibt es zwei Möglichkeiten, die Geschlechtsorgane von Mann und Frau
einzuteilen:

\begin{itemize}
  \item innere und äußere Geschlechtsorgane: Hier wird
  einzig auf ihre Lage beim Erwachsenen Bezug genommen.
  (außen = sichtbar, innen = nicht sichtbar)
  \item primäre und sekundäre Geschlechtsorgane: Bei dieser Einteilung
  unterscheidet man zwischen Merkmalen, die bereits bei
  Geburt vorhanden sind (=primär), und Merkmalen, die sich
  erst in der Pubertät entwickeln
  (=sekundär).
  \begin {itemize}
    \item primäre Geschlechtsmerkmale: Gonaden (Keimdrüsen)
    \item sekundäre Geschlechtsmerkmale: ableitende Geschlechtsorgane, äußeres
    Genitale, außerhalb des Genitalbereichs auftretende Geschlechtsmerkmale
  \end{itemize}
\end{itemize}
}


\subsection{Männliche Geschlechtsorgane}

Beim Mann ist die Unterteilung folgendermaßen:
\begin{itemize}
  \item innen: Hoden, Nebenhoden, Samenstrang, Prostata und Samenbläschen
  \item außen: Glied, Hodensack
\end{itemize}


\Layout{27}
{}
{}
{}
{./images/hirtler-aass/Genitale-m.pdf}
{Die männlichen Geschlechtsorgane im
Querschnitt}
{Lena Hirtler}


\TextSlide{Hoden}
{%
Im Hoden findet die Spermiogenese, also die Bildung der
Spermien, statt. Der Hoden hat eine Pflaumenform mit einem
ungefähren Längsdurchmesser von 5\CM* und ist paarig 
angelegt. Er steht über die Nebenhoden in Verbindung mit den
ableitenden Samenwegen.

\Z{Da der Hoden während der Schwangerschaftswochen im
Ungeborenen von der Bauchhöhle seinen Weg in den Hodensack
finden soll, muss dies nach der Geburt kontrolliert werden.}
}
{
\begin{itemize}
  \item Bildung von Spermien
  \item Paarig
  \item Sollte nach Geburt im Hodensack angekommen sein
\end{itemize}
}
{}{}{}


\TextSlide{Nebenhoden}
{%
Der Nebenhoden liegt dem Hoden kappenförmig an. Er befindet
sich daher auch im Hodensack. Er ist verantwortlich für
einige Enzyme die z.\,B. den Spermien helfen sollen, die
Eizelle zu finden. Weiters verfügt die Wand des Nebenhodens
über eine starke Muskulatur, die mit einer eigenen
Peristaltik den Transport der Spermien bewirkt. Die Spermien
bleiben im Nebenhoden ca. \VonBis{7}{14} Tage und erlangen
im Nebenhoden ihre volle Befruchtungsfähigkeit.
}{
\begin{itemize}
  \item Im Hodensack
  \item Enzymbildung
  \item Transport der Spermien
\end{itemize}

}{}{}{}


\TextSlide{Samenstrang}
{
Unter Samenstrang versteht man die Struktur, die vom Hoden durch den
Leistenkanal in Richtung Prostata und Penis führt.
In ihm verlaufen folgende Strukturen:
\begin{itemize}
  \item Hodengefäße
  \item Samenleiter
  \item Nervenfasern des vegetativen Nervensystems
\end{itemize}

In der Wand des Samenstrangs verläuft der Hodenheber-Muskel,
ein Ausläufer der schrägen Bauchmuskeln. Er ist dafür
verantwortlich, den Hoden vor zu großen
Temperaturunterschieden zu schützen (z.\,B. zieht er den
Hoden bei Kälte Richtung Bauchraum).
}
{
\begin{itemize}
  \item Enthält Samenleiter
  \item Läuft im Leistenkanal
  \item Hodenheber-Muskel
\end{itemize}
}{}{}{}


\TextSlide{Prostata und Bläschendrüse}
{%
Die Prostata ist eine etwa kastaninengroße Drüse. Sie liegt
direkt der Harnblase an und die Harnröhre läuft durch sie
hindurch. Die Prostata liegt außerdem dem Enddarm an -- sie
ist so tastbar und jeder Mann über 50 sollte sie regelmäßig
untersuchen lassen. Mit zunehmendem Alter kann sie
nämlich entarten. Postatakrebs ist eine der häufigsten
Tumorarten beim Mann.
Die Prostata ist zuständig für ca. \Z{30\PC*} der
Ejakulatflüssikeit und erzeugt Enzyme, die die Beweglichkeit
der Spermien fördern.
Die Bläschendrüse liegt der Prostata an beiden Seiten an.
Sie sorgt für ca. \Z{60\PC*} der Ejakulatflüssigkeit.
}
{
\begin{itemize}
  \item Prostata → \Z{30\PC*} des Ejakulates, Beweglichkeit der Spermien
  \item Bläschendrüse → \Z{60\PC*} des Ejakulates
\end{itemize}
}{}{}{}


\TextSlide{Penis}
{%
\Lang{Syn.} Glied.
Der wichtigste Bestandteil des Penis ist der Schwellkörper.
\Z{Er besteht aus zwei Untereinheiten:
\begin{itemize}
  \item Corpus cavernosum: Ist in zwei Teile geteilt und
  befindet sich an der Oberseite. Es sorgt durch
  Blutaufnahme für das Steif-werden des Penis
  \item Corpus spongiosum: liegt an der Unterseite. In ihm
  liegt die Harnröhre, die in der Eichel endet. Dadurch dass
  sich dieser Körper nicht so sehr mit Blut füllt, wird die
  Harnröhre nicht abgedrückt und die Ejakulation ermöglicht.
\end{itemize}
}

Der berührungsempfindlichste Teil des Penis ist die Eichel.
Sie ist im erektionslosen Zustand unter der Vorhaut
verborgen, die sich bei Erektion zurückzieht.
\Z{Erektion und Ejakulation erfolgen reflektorisch. Zur
Erektion kommt es durch eine Kombination aus physischen und
psychischen Reizen. Die Ejakulation erfolgt durch
Kontraktionen der Beckenbodenmuskulatur.
Unter einer \I{Potenzstörung} versteht man die Unmöglichkeit zur
Durchführung eines Beischlafes. Bei einer Potenzstörung kann
entweder die Erektion gestört sein oder die Ejakulation
verhindert sein.}
}
{
\begin{itemize}
  \item Schwellkörper - bestehend aus zwei Teilen
  \item Vorhaut - Eichel
  \item Erektion vs. Ejakulation
  \item Potenzstörung
\end{itemize}
}{}{}{}


\TextSlide{Hodensack}
{%
Der Hodensack umschließt \EE{Hoden}, \EE{Nebenhoden} und
\EE{Samenstränge}. Er ist aus einzelnen Schichten aufgebaut.
Am wichtigsten ist, dass auch er Fasern des
Hodenheber-Muskels enthält.

Der Hoden kann sich im Hodensack um seine eigene Achse
drehen. Wenn dies passiert und er sich nicht zurückdreht,
spricht man von einer \T{Hodentorsion}. Es werden die
wichtigen \E{Blutgefäße abgeschnürt}, es kommt zum Absterben
des Hodens begleitet mit unvorstellbaren Schmerzen und
vegetativen Symptomen wie Übelkeit und Erbrechen. Wenn diese
Torsion nicht schnellstmöglich operiert wird, stirbt dieser
Hoden vollständig ab. \Z{Bei der Operation wird der Hoden an
den Hodensack festgenäht, damit dies nicht nocheinmal
vorkommen kann.}
}{%
\begin{itemize}
  \item Hodenheber-Muskel
  \item CAVE: Hodentorsion
\end{itemize}
}{}{}{}

\subsection{Weibliche Geschlechtsorgane}

Die Unterteilung der Geschlechtsorgane der Frau ist:
\begin{itemize}
  \item innen: Eierstöcke, Eileiter, Gebärmutter, Scheide
  \item außen: Schamlippen, Klitoris
\end{itemize}

\Layout{27}
{}
{}
{}
{./images/hirtler-aass/Genitale-w.pdf}
{Weibliche Geschlechtsorgane,
Querschnitt}
{Lena Hirtler}


\TextSlide{Eierstöcke}
{%
In den Eierstöcken (Ovarien) reifen die Eizellen heran. Sie
sind paarig und liegen im Beckenbereich und sind ca. $4
\times 2 \times 1$\CM* groß. Dadurch dass in den Ovarien
regelmäßig Eizellen heranreifen, kann man in ihnen alle
Entwicklungsstadien finden. Nach dem Eisprung ist das Ovar
auch zuständig für die Bildung von Hormonen, die im Falle
einer Befruchtung der Eizelle schwangerschaftserhaltend
wirken.
}
{
\begin{itemize}
  \item 2 Stück, im Beckenbereich
  \item Heranreifen von Eizellen
  \item Hormonbildung
\end{itemize}
}{}{}{}


\TextSlide{Eileiter}
{%
Der Eileiter \Z{(Tuba uterina / Salpinx)} ist ein ca.
\VonBis{10}{18}\CM* langer Schlauch mit einer freien Öffnung
in die Bauchhöhle. Nahe dem Ovar zeigt der Eileiter einen
Trichter mit fingerförmigen Ausläufern, worin die Eizelle
normalerweise nach dem Eisprung
\protect\enquote{aufgefangen} wird.
Die \E{Befruchtung} der Eizelle findet im Eileiter statt,
die befruchtete Eizelle gelangt dann normalerweise über den
Eileiter in die Gebärmutter, wo sie sich einnistet.

Wenn die Eizelle den \protect\enquote{falschen} Weg nimmt,
oder sich nach der Befruchtung nicht in die Gebärmutter
weiterbewegt, kann es zu einer \EE{Eileiterschwangerschaft}
kommen. Dies kann zu einer Lebensgefährdung der Mutter
führen, wenn durch das Wachstum des Embryos der Eileiter
\protect\enquote{gesprengt} wird und es zu massiven
Blutungen kommt.
}
{
\begin{itemize}
  \item \VonBis{10}{18}\CM*, freie Öffnung in die Bauchhöhle
  \item Trichterförmig
  \item Befruchtung
  \item Gefahr: Eileiterschwangerschaft
\end{itemize}
}{}{}{}


\TextSlide{Gebärmutter}
{%
Die \T{Gebärmutter} (\T{Uterus}) kann grob in zwei Teile,
dem \E{Körper} und dem \E{Hals} (Zervix), eingeteilt werden.
In den Körper münden die beiden Eileiter, die die
befruchtete Eizelle in den Uterus führen sollen.
Der Körper der Gebärmutter ist äußerst muskulös um bei der
Geburt die Austreibung des Fetus zu ermöglichen.
Die Schleimhaut unterliegt bei der geschlechtsreifen Frau
zyklischen Veränderungen (siehe weiblicher Zyklus,
\FullPrettyRef{topic:weiblicher-zyklus}).

Über den Hals mündet die Gebärmutter in die Scheide, die
Mündung wird als \T{Muttermund} bezeichnet.
Der Gebärmutterhalskrebs \Z{(Zervixkarzinom)} ist eine
häufige Krebsart bei der Frau. 
Daher sollte sie regelmäßig beim Gynäkologen einen Abstrich
machen lassen um eine Früherkennung zu ermöglichen.
}
{
\begin{itemize}
  \item 2 Teile: Körper, Hals
  \item Körper → muskulös, Schleimhaut unterliegt zyklischen Veränderungen
  \item CAVE: Gebärmutterhalskrebs
\end{itemize}
}{}{}{}


\TextSlide{Scheide}
{%
Auf den Muttermund der Gebärmutter folgt die Scheide. Sie
ist erheblich dehnbar und ihr Epithel unterliegt ebenfalls
hormonellen Zyklusschwankungen. Ihre Dehnbarkeit ermöglicht
die Anpassung an die Passage des kindlichen Kopfes während
des Geburtsvorganges.
}
{
\begin{itemize}
  \item Dehnbarkeit → Geburtsvorgang
  \item hormonelle Zyklusschwankungen
\end{itemize}
}{}{}{}


\TextSlide{Schamlippen und Klitoris}
{%
Bei der Frau gibt es insgesamt vier Schamlippen, jeweils
zwei große und zwei kleine. \Z{Die großen Schamlippen liegen
außen, die kleinen innen und umrahmen den Scheideneingang
mit den Öffnungen der Scheide und der Harnröhre.
An den kleinen Schamlippen münden Drüsen, die für den
Feuchtigkeitserhalt dieser Region zuständig sind.
Die Frau besitzt, wie der Mann auch, Schwellkörper. Diese
liegen in der Klitoris und im Bereich der großen
Schamlippen. Vergleichsweise mit der Eichel ist die Klitoris
die berührungsempfindlichste Stelle der Frau.}
}
{
\begin{itemize}
  \item 4 Schamlippen
  \item Drüsen
  \item Klitoris + Schwellkörpers
\end{itemize}
}{}{}{}

\subsubsection{Der weibliche Zyklus}
\label{topic:weiblicher-zyklus}\index{Zyklus, weiblicher}

\TextSlide{\TxBeschreibung}
{%
Der weibliche Zyklus dauert \EE{28 Tage}, dies entspricht
einem Lunarmonat (Zeit von Vollmond bis Vollmond). Er wird
von Menstruation (Monatsblutung) zu Menstruation gerechnet.
Während eines Zyklus wird die Uterusschleimhaut aufgebaut
und bei der Menstruation wieder abgestoßen.
Der Zyklus ist durch Östrogen und Progesteron
hormongesteuert. Der \EE{Eisprung} findet \E{in der Mitte}
des Zyklus statt, ca. am 14. Tag. Wenn es zu keiner
Befruchtung der Eizelle kommt, läuft der Zyklus normal
weiter, es kommt zur Menstruation.
Wenn es zu einer Befruchtung der Eizelle kommt, wird der
Zyklus hormonell angehalten und die Eizelle kann sich in die
Uterusschleimhaut einnisten.
}{
\begin{itemize}
  \item 28 Tage
  \item Eisprung am 14. Tag 
  \item keiner Befruchtung → Menstruation
  \item Befruchtung → Schwangerschaft
\end{itemize}
}
{}{}{}



% \begin{WidePar}
% \noindent\begin{minipage}{\linewidth}
% \noindent\includegraphics[width=\linewidth]{}
% 
% \Captionof{figure}{
% \SourcePicture{Hirtler}}
% \end{minipage}
% \end{WidePar}

\Layout{37}
{}
{}
{}
{./images/hirtler-aass/Zyklus.pdf}
{Der weibliche Zyklus. Oben:
Follikelreifung. Unten: Auf- und Abbau der
Uterusschleimhaut. Am Tag 0 des Zyklus wird die
Uterusschleimhaut abgestoßen, es kommt zur Regelblutung,
danach wird die Schleimhaut wieder aufgebaut. Am Tag 14
erfolgt der Eisprung.
Kommt es zu keiner Befruchtung, wird die Schleimhaut erneut
am Tag 28 abgestoßen. Nach einer Befruchtung wandert die
Eizelle weiter in den Uterus und nistet sich dort ein.}
{Lena Hirtler}


\ChapterEnd